% Shared TikZ macros for the architecture figures (Medverse vs. Ours).
% Load with % Shared TikZ macros for the architecture figures (Medverse vs. Ours).
% Load with % Shared TikZ macros for the architecture figures (Medverse vs. Ours).
% Load with % Shared TikZ macros for the architecture figures (Medverse vs. Ours).
% Load with \input{arch-macros.tex} after \usepackage{tikz,ifthen} and
% \usetikzlibrary{arrows.meta,calc,positioning}. All macro names are
% prefixed `arch` to avoid clashing with pgf/tikz internals.

% ---- palette (approximating the PowerPoint drafts) ----
\definecolor{archTargetGreen}{HTML}{9BC7AC}
\definecolor{archContextPink}{HTML}{E8B8B0}
\definecolor{archAttnBlue}{HTML}{BBD6EC}
\definecolor{archEncGray}{HTML}{E4E4E4}
\definecolor{archArrowPale}{HTML}{CFE6D6}
\definecolor{archCtBg}{HTML}{2B2B2B}
\definecolor{archCtBlob}{HTML}{8C8C8C}
\definecolor{archCtxYellow}{HTML}{E8D24C}
% darker, more saturated siblings of archTargetGreen/archContextPink --
% pastel fills read poorly as 1pt strokes, and mask/label content needs to
% read as "which volume owns this" at a glance, so masks/labels and any
% line (arrow, brace) touching that volume's row use these instead of the
% pastel role color or the old flat archMaskGold.
\definecolor{archTargetAccent}{HTML}{2F6B4A}  % target: darker green
\definecolor{archContextAccent}{HTML}{A8453E} % context: darker red

% fake-isometric depth skew shared by every cube
\newcommand{\archdx}{0.30}
\newcommand{\archdy}{0.20}

% ---- shared text-role macros ----
% Every arch_*.tex figure is compiled standalone at the SAME base size
% (documentclass[tikz,border=4pt,12pt]{standalone}, matching the thesis
% body's 12pt), so \normalsize etc. mean the same absolute point size in
% every figure. Using these three named roles instead of ad hoc \Large /
% \bfseries\normalsize / \scriptsize per figure keeps "a box title" (etc.)
% the same size everywhere it appears, regardless of which figure it's in
% or how that figure's canvas happens to be scaled by \includegraphics.
%   \archtitle    -- figure/method/panel titles, the "xL" repeat marker
%   \archboxlabel -- sub-block header labels (Fuse, Encoder, attention, ...)
%   \archnote     -- small secondary annotations (placeholder captions)
% Plain, unlabeled node text (cube/token labels, row/lane labels) stays at
% the bare \normalsize default -- no macro needed for that role.
%
% DECLARATIONS, not argument-taking macros -- deliberately, so they can
% precede a raw `\\` inside a tikz node's text (as `{\archboxlabel line
% one\\line two}`) without hiding that `\\` one macro-expansion level deep.
% A `\\` buried inside a wrapped {\archboxlabel{...}} argument breaks
% TikZ's own align=center line-splitting (it scans for `\\` before the
% argument expands) -- use these the same way `\bfseries`/`\Large` are
% already used directly inline elsewhere in these figures.
\newcommand{\archtitle}{\Large\bfseries}
\newcommand{\archboxlabel}{\bfseries\normalsize}
\newcommand{\archnote}{\scriptsize}

% ---------------------------------------------------------------
% \archvolcube{x}{y}{w}{h}{topcolor}{content}{label}
%   content: ct | ctmask | predmask | <flat color name>
% ---------------------------------------------------------------
\newcommand{\archvolcube}[7]{%
  \begin{scope}
    \pgfmathsetmacro{\cx}{#1}\pgfmathsetmacro{\cy}{#2}%
    \pgfmathsetmacro{\w}{#3}\pgfmathsetmacro{\h}{#4}%
    \fill[#5] (\cx,\cy+\h) -- ++(\archdx,\archdy) -- ++(\w,0) -- ++(-\archdx,-\archdy) -- cycle;
    \fill[#5!60!black] (\cx+\w,\cy) -- ++(\archdx,\archdy) -- ++(0,\h) -- ++(-\archdx,-\archdy) -- cycle;
    \begin{scope}
      \clip (\cx,\cy) rectangle ++(\w,\h);
      \ifthenelse{\equal{#6}{ct}}{%
        \fill[archCtBg] (\cx,\cy) rectangle ++(\w,\h);
        \fill[archCtBlob] plot[smooth cycle] coordinates
          {(\cx+.25*\w,\cy+.55*\h) (\cx+.45*\w,\cy+.75*\h) (\cx+.7*\w,\cy+.6*\h)
           (\cx+.6*\w,\cy+.3*\h) (\cx+.35*\w,\cy+.25*\h)};
      }{\ifthenelse{\equal{#6}{ctmask}}{%
        \fill[archCtBg] (\cx,\cy) rectangle ++(\w,\h);
        \fill[archCtBlob] plot[smooth cycle] coordinates
          {(\cx+.25*\w,\cy+.55*\h) (\cx+.45*\w,\cy+.75*\h) (\cx+.7*\w,\cy+.6*\h)
           (\cx+.6*\w,\cy+.3*\h) (\cx+.35*\w,\cy+.25*\h)};
        \fill[archCtxYellow] plot[smooth cycle] coordinates
          {(\cx+.42*\w,\cy+.42*\h) (\cx+.52*\w,\cy+.5*\h) (\cx+.5*\w,\cy+.35*\h) (\cx+.44*\w,\cy+.3*\h)};
      }{\ifthenelse{\equal{#6}{predmask}}{%
        \fill[black] (\cx,\cy) rectangle ++(\w,\h);
        \fill[white] plot[smooth cycle] coordinates
          {(\cx+.4*\w,\cy+.65*\h) (\cx+.6*\w,\cy+.7*\h) (\cx+.62*\w,\cy+.45*\h)
           (\cx+.48*\w,\cy+.3*\h) (\cx+.38*\w,\cy+.4*\h)};
      }{%
        \fill[#6] (\cx,\cy) rectangle ++(\w,\h);
      }}}%
    \end{scope}
    \draw[black,thin] (\cx,\cy) rectangle ++(\w,\h);
    \draw[black,thin] (\cx,\cy+\h) -- ++(\archdx,\archdy) -- ++(\w,0) -- ++(-\archdx,-\archdy) -- cycle;
    \draw[black,thin] (\cx+\w,\cy) -- ++(\archdx,\archdy) -- ++(0,\h) -- ++(-\archdx,-\archdy) -- cycle;
    \node[anchor=north,align=center] at (\cx+\w/2,\cy-0.1) {#7};
  \end{scope}
}

% ---------------------------------------------------------------
% \archvolcubeimg{x}{y}{w}{h}{topcolor}{imgpath}{label} -- same cube as
% \archvolcube, but the front face shows a real raster image (clipped to
% exactly fill the face, aspect ratio not preserved) instead of one of
% \archvolcube's procedural ct/ctmask/predmask blobs. Needs graphicx
% loaded by the including figure. A separate macro rather than a new
% \archvolcube content mode, so existing ct/ctmask/predmask figures are
% untouched.
% ---------------------------------------------------------------
\newcommand{\archvolcubeimg}[7]{%
  \begin{scope}
    \pgfmathsetmacro{\cx}{#1}\pgfmathsetmacro{\cy}{#2}%
    \pgfmathsetmacro{\w}{#3}\pgfmathsetmacro{\h}{#4}%
    \fill[#5] (\cx,\cy+\h) -- ++(\archdx,\archdy) -- ++(\w,0) -- ++(-\archdx,-\archdy) -- cycle;
    \fill[#5!60!black] (\cx+\w,\cy) -- ++(\archdx,\archdy) -- ++(0,\h) -- ++(-\archdx,-\archdy) -- cycle;
    \begin{scope}
      \clip (\cx,\cy) rectangle ++(\w,\h);
      \node[anchor=south west, inner sep=0pt] at (\cx,\cy)
        {\includegraphics[width=\w cm,height=\h cm]{#6}};
    \end{scope}
    \draw[black,thin] (\cx,\cy) rectangle ++(\w,\h);
    \draw[black,thin] (\cx,\cy+\h) -- ++(\archdx,\archdy) -- ++(\w,0) -- ++(-\archdx,-\archdy) -- cycle;
    \draw[black,thin] (\cx+\w,\cy) -- ++(\archdx,\archdy) -- ++(0,\h) -- ++(-\archdx,-\archdy) -- cycle;
    \node[anchor=north,align=center] at (\cx+\w/2,\cy-0.1) {#7};
  \end{scope}
}

% ---------------------------------------------------------------
% \archfeatcube{x}{y}{w}{h}{color}{label} -- flat-colored feature volume
% ---------------------------------------------------------------
\newcommand{\archfeatcube}[6]{%
  \begin{scope}
    \pgfmathsetmacro{\cx}{#1}\pgfmathsetmacro{\cy}{#2}%
    \pgfmathsetmacro{\w}{#3}\pgfmathsetmacro{\h}{#4}%
    \fill[#5] (\cx,\cy) rectangle ++(\w,\h);
    \fill[#5!80!white] (\cx,\cy+\h) -- ++(\archdx,\archdy) -- ++(\w,0) -- ++(-\archdx,-\archdy) -- cycle;
    \fill[#5!60!black] (\cx+\w,\cy) -- ++(\archdx,\archdy) -- ++(0,\h) -- ++(-\archdx,-\archdy) -- cycle;
    \draw[black,thin] (\cx,\cy) rectangle ++(\w,\h);
    \draw[black,thin] (\cx,\cy+\h) -- ++(\archdx,\archdy) -- ++(\w,0) -- ++(-\archdx,-\archdy) -- cycle;
    \draw[black,thin] (\cx+\w,\cy) -- ++(\archdx,\archdy) -- ++(0,\h) -- ++(-\archdx,-\archdy) -- cycle;
    \node[align=center] at (\cx+\w/2,\cy+\h/2) {#6};
  \end{scope}
}

% ---------------------------------------------------------------
% \archfunnel{x}{y}{h}{wwide}{wnarrow}{dir}{color}{label}
%   dir=1  : wide on left, narrow on right  (encoder)
%   dir=-1 : narrow on left, wide on right  (decoder)
% ---------------------------------------------------------------
\newcommand{\archfunnel}[8]{%
  \begin{scope}
    \pgfmathsetmacro{\cx}{#1}\pgfmathsetmacro{\cy}{#2}\pgfmathsetmacro{\h}{#3}%
    \pgfmathsetmacro{\ww}{#4}\pgfmathsetmacro{\wn}{#5}%
    \pgfmathsetmacro{\inset}{(\ww-\wn)/2}%
    \ifnum#6=1
      \fill[#7] (\cx,\cy) -- (\cx+\ww,\cy+\inset) -- (\cx+\ww,\cy+\h-\inset) -- (\cx,\cy+\h) -- cycle;
      \draw[black,thin] (\cx,\cy) -- (\cx+\ww,\cy+\inset) -- (\cx+\ww,\cy+\h-\inset) -- (\cx,\cy+\h) -- cycle;
    \else
      \fill[#7] (\cx,\cy+\inset) -- (\cx+\ww,\cy) -- (\cx+\ww,\cy+\h) -- (\cx,\cy+\h-\inset) -- cycle;
      \draw[black,thin] (\cx,\cy+\inset) -- (\cx+\ww,\cy) -- (\cx+\ww,\cy+\h) -- (\cx,\cy+\h-\inset) -- cycle;
    \fi
    \node[align=center] at (\cx+\ww/2,\cy+\h/2) {#8};
  \end{scope}
}

% ---------------------------------------------------------------
% \archbowtie{x}{y}{w}{h}{pinch}{color}{label} -- Medverse-style
% single-row "dual U-Net" hourglass (funnel-in then funnel-out).
% ---------------------------------------------------------------
\newcommand{\archbowtie}[7]{%
  \begin{scope}
    \pgfmathsetmacro{\cx}{#1}\pgfmathsetmacro{\cy}{#2}%
    \pgfmathsetmacro{\w}{#3}\pgfmathsetmacro{\h}{#4}\pgfmathsetmacro{\p}{#5}%
    \pgfmathsetmacro{\inset}{(\h-\p)/2}%
    \fill[#6] (\cx,\cy) -- (\cx+\w/2,\cy+\inset) -- (\cx+\w,\cy) -- (\cx+\w,\cy+\h)
      -- (\cx+\w/2,\cy+\h-\inset) -- (\cx,\cy+\h) -- cycle;
    \draw[black,thin] (\cx,\cy) -- (\cx+\w/2,\cy+\inset) -- (\cx+\w,\cy) -- (\cx+\w,\cy+\h)
      -- (\cx+\w/2,\cy+\h-\inset) -- (\cx,\cy+\h) -- cycle;
    \node[align=center] at (\cx+\w/2,\cy+\h/2) {#7};
  \end{scope}
}

% ---------------------------------------------------------------
% \archtokengrid{x}{y}{rows}{cols}{cell}{color} -- grid of patch tokens
% ---------------------------------------------------------------
\newcommand{\archtokengrid}[6]{%
  \begin{scope}
    \pgfmathsetmacro{\cx}{#1}\pgfmathsetmacro{\cy}{#2}\pgfmathsetmacro{\cell}{#5}%
    \pgfmathtruncatemacro{\rmax}{#3-1}%
    \pgfmathtruncatemacro{\cmax}{#4-1}%
    \foreach \r in {0,...,\rmax}{%
      \foreach \c in {0,...,\cmax}{%
        \fill[#6] ({\cx+\c*\cell},{\cy+\r*\cell}) rectangle ++(\cell*0.82,\cell*0.82);
        \draw[black,thin] ({\cx+\c*\cell},{\cy+\r*\cell}) rectangle ++(\cell*0.82,\cell*0.82);
      }%
    }%
  \end{scope}
}

% ---------------------------------------------------------------
% \archarrow{start}{end}{color} -- block arrow between stages
% ---------------------------------------------------------------
% #1 and #2 are full coordinates including parentheses, e.g. (1,2)
\newcommand{\archarrow}[3]{%
  \draw[-{Latex[length=2.4mm,width=2mm]},line width=2.2pt,#3] #1 -- #2;
}

% ---------------------------------------------------------------
% Dataflow-arrow vocabulary -- opt-in via \archarrowset{compact|detailed},
% called once near the top of a figure (after \input{arch-macros.tex}).
% \archarrow keeps its classic small size by default for any arch_*.tex
% figure that never calls \archarrowset -- this block is purely additive.
%   archflowstyle (tikz key, NOT a macro -- use bare, no backslash) --
%                                  current draw options (head + line
%                                  width), for a hand-rolled elbow \draw
%                                  that \archarrow's straight "--" can't
%                                  express (e.g. "-|" routing). Must be a
%                                  /.style key rather than a \def'd macro:
%                                  pgf's arrow-tip parser doesn't split on
%                                  a comma embedded in a plain macro's
%                                  expansion when used as a \draw option.
%   \archflowNeutral/Target/Context -- the only 3 colors a plain dataflow
%                                  arrow should use: generic same-role
%                                  passthrough; target's own real output
%                                  leaving the decoder; a real label/mask
%                                  tensor entering fusion or attention
%                                  (context's ground truth, or target's
%                                  own real prior fed back from a
%                                  previous cascade level -- same role
%                                  either way: real mask data, not a
%                                  feature)
%   \archflow{p1}{p2}{role}     -- \archarrow wrapper keyed by role name
%                                  (role in {neutral, target, context})
% ---------------------------------------------------------------
\newcommand{\archarrowset}[1]{%
  \ifthenelse{\equal{#1}{detailed}}{%
    \tikzset{archflowstyle/.style={-{Latex[length=3.6mm,width=3.2mm]},line width=3.2pt}}%
  }{%
    \tikzset{archflowstyle/.style={-{Latex[length=2.4mm,width=2mm]},line width=2.2pt}}%
  }%
  \renewcommand{\archarrow}[3]{\draw[archflowstyle,##3] ##1 -- ##2;}%
}
\newcommand{\archflowNeutral}{black!30}
\newcommand{\archflowTarget}{archTargetAccent}
\newcommand{\archflowContext}{archContextAccent}
\newcommand{\archflow}[3]{%
  \ifthenelse{\equal{#3}{target}}{\archarrow{#1}{#2}{\archflowTarget}}{%
  \ifthenelse{\equal{#3}{context}}{\archarrow{#1}{#2}{\archflowContext}}{%
  \archarrow{#1}{#2}{\archflowNeutral}}}%
}
 after \usepackage{tikz,ifthen} and
% \usetikzlibrary{arrows.meta,calc,positioning}. All macro names are
% prefixed `arch` to avoid clashing with pgf/tikz internals.

% ---- palette (approximating the PowerPoint drafts) ----
\definecolor{archTargetGreen}{HTML}{9BC7AC}
\definecolor{archContextPink}{HTML}{E8B8B0}
\definecolor{archAttnBlue}{HTML}{BBD6EC}
\definecolor{archEncGray}{HTML}{E4E4E4}
\definecolor{archArrowPale}{HTML}{CFE6D6}
\definecolor{archCtBg}{HTML}{2B2B2B}
\definecolor{archCtBlob}{HTML}{8C8C8C}
\definecolor{archCtxYellow}{HTML}{E8D24C}
% darker, more saturated siblings of archTargetGreen/archContextPink --
% pastel fills read poorly as 1pt strokes, and mask/label content needs to
% read as "which volume owns this" at a glance, so masks/labels and any
% line (arrow, brace) touching that volume's row use these instead of the
% pastel role color or the old flat archMaskGold.
\definecolor{archTargetAccent}{HTML}{2F6B4A}  % target: darker green
\definecolor{archContextAccent}{HTML}{A8453E} % context: darker red

% fake-isometric depth skew shared by every cube
\newcommand{\archdx}{0.30}
\newcommand{\archdy}{0.20}

% ---- shared text-role macros ----
% Every arch_*.tex figure is compiled standalone at the SAME base size
% (documentclass[tikz,border=4pt,12pt]{standalone}, matching the thesis
% body's 12pt), so \normalsize etc. mean the same absolute point size in
% every figure. Using these three named roles instead of ad hoc \Large /
% \bfseries\normalsize / \scriptsize per figure keeps "a box title" (etc.)
% the same size everywhere it appears, regardless of which figure it's in
% or how that figure's canvas happens to be scaled by \includegraphics.
%   \archtitle    -- figure/method/panel titles, the "xL" repeat marker
%   \archboxlabel -- sub-block header labels (Fuse, Encoder, attention, ...)
%   \archnote     -- small secondary annotations (placeholder captions)
% Plain, unlabeled node text (cube/token labels, row/lane labels) stays at
% the bare \normalsize default -- no macro needed for that role.
%
% DECLARATIONS, not argument-taking macros -- deliberately, so they can
% precede a raw `\\` inside a tikz node's text (as `{\archboxlabel line
% one\\line two}`) without hiding that `\\` one macro-expansion level deep.
% A `\\` buried inside a wrapped {\archboxlabel{...}} argument breaks
% TikZ's own align=center line-splitting (it scans for `\\` before the
% argument expands) -- use these the same way `\bfseries`/`\Large` are
% already used directly inline elsewhere in these figures.
\newcommand{\archtitle}{\Large\bfseries}
\newcommand{\archboxlabel}{\bfseries\normalsize}
\newcommand{\archnote}{\scriptsize}

% ---------------------------------------------------------------
% \archvolcube{x}{y}{w}{h}{topcolor}{content}{label}
%   content: ct | ctmask | predmask | <flat color name>
% ---------------------------------------------------------------
\newcommand{\archvolcube}[7]{%
  \begin{scope}
    \pgfmathsetmacro{\cx}{#1}\pgfmathsetmacro{\cy}{#2}%
    \pgfmathsetmacro{\w}{#3}\pgfmathsetmacro{\h}{#4}%
    \fill[#5] (\cx,\cy+\h) -- ++(\archdx,\archdy) -- ++(\w,0) -- ++(-\archdx,-\archdy) -- cycle;
    \fill[#5!60!black] (\cx+\w,\cy) -- ++(\archdx,\archdy) -- ++(0,\h) -- ++(-\archdx,-\archdy) -- cycle;
    \begin{scope}
      \clip (\cx,\cy) rectangle ++(\w,\h);
      \ifthenelse{\equal{#6}{ct}}{%
        \fill[archCtBg] (\cx,\cy) rectangle ++(\w,\h);
        \fill[archCtBlob] plot[smooth cycle] coordinates
          {(\cx+.25*\w,\cy+.55*\h) (\cx+.45*\w,\cy+.75*\h) (\cx+.7*\w,\cy+.6*\h)
           (\cx+.6*\w,\cy+.3*\h) (\cx+.35*\w,\cy+.25*\h)};
      }{\ifthenelse{\equal{#6}{ctmask}}{%
        \fill[archCtBg] (\cx,\cy) rectangle ++(\w,\h);
        \fill[archCtBlob] plot[smooth cycle] coordinates
          {(\cx+.25*\w,\cy+.55*\h) (\cx+.45*\w,\cy+.75*\h) (\cx+.7*\w,\cy+.6*\h)
           (\cx+.6*\w,\cy+.3*\h) (\cx+.35*\w,\cy+.25*\h)};
        \fill[archCtxYellow] plot[smooth cycle] coordinates
          {(\cx+.42*\w,\cy+.42*\h) (\cx+.52*\w,\cy+.5*\h) (\cx+.5*\w,\cy+.35*\h) (\cx+.44*\w,\cy+.3*\h)};
      }{\ifthenelse{\equal{#6}{predmask}}{%
        \fill[black] (\cx,\cy) rectangle ++(\w,\h);
        \fill[white] plot[smooth cycle] coordinates
          {(\cx+.4*\w,\cy+.65*\h) (\cx+.6*\w,\cy+.7*\h) (\cx+.62*\w,\cy+.45*\h)
           (\cx+.48*\w,\cy+.3*\h) (\cx+.38*\w,\cy+.4*\h)};
      }{%
        \fill[#6] (\cx,\cy) rectangle ++(\w,\h);
      }}}%
    \end{scope}
    \draw[black,thin] (\cx,\cy) rectangle ++(\w,\h);
    \draw[black,thin] (\cx,\cy+\h) -- ++(\archdx,\archdy) -- ++(\w,0) -- ++(-\archdx,-\archdy) -- cycle;
    \draw[black,thin] (\cx+\w,\cy) -- ++(\archdx,\archdy) -- ++(0,\h) -- ++(-\archdx,-\archdy) -- cycle;
    \node[anchor=north,align=center] at (\cx+\w/2,\cy-0.1) {#7};
  \end{scope}
}

% ---------------------------------------------------------------
% \archvolcubeimg{x}{y}{w}{h}{topcolor}{imgpath}{label} -- same cube as
% \archvolcube, but the front face shows a real raster image (clipped to
% exactly fill the face, aspect ratio not preserved) instead of one of
% \archvolcube's procedural ct/ctmask/predmask blobs. Needs graphicx
% loaded by the including figure. A separate macro rather than a new
% \archvolcube content mode, so existing ct/ctmask/predmask figures are
% untouched.
% ---------------------------------------------------------------
\newcommand{\archvolcubeimg}[7]{%
  \begin{scope}
    \pgfmathsetmacro{\cx}{#1}\pgfmathsetmacro{\cy}{#2}%
    \pgfmathsetmacro{\w}{#3}\pgfmathsetmacro{\h}{#4}%
    \fill[#5] (\cx,\cy+\h) -- ++(\archdx,\archdy) -- ++(\w,0) -- ++(-\archdx,-\archdy) -- cycle;
    \fill[#5!60!black] (\cx+\w,\cy) -- ++(\archdx,\archdy) -- ++(0,\h) -- ++(-\archdx,-\archdy) -- cycle;
    \begin{scope}
      \clip (\cx,\cy) rectangle ++(\w,\h);
      \node[anchor=south west, inner sep=0pt] at (\cx,\cy)
        {\includegraphics[width=\w cm,height=\h cm]{#6}};
    \end{scope}
    \draw[black,thin] (\cx,\cy) rectangle ++(\w,\h);
    \draw[black,thin] (\cx,\cy+\h) -- ++(\archdx,\archdy) -- ++(\w,0) -- ++(-\archdx,-\archdy) -- cycle;
    \draw[black,thin] (\cx+\w,\cy) -- ++(\archdx,\archdy) -- ++(0,\h) -- ++(-\archdx,-\archdy) -- cycle;
    \node[anchor=north,align=center] at (\cx+\w/2,\cy-0.1) {#7};
  \end{scope}
}

% ---------------------------------------------------------------
% \archfeatcube{x}{y}{w}{h}{color}{label} -- flat-colored feature volume
% ---------------------------------------------------------------
\newcommand{\archfeatcube}[6]{%
  \begin{scope}
    \pgfmathsetmacro{\cx}{#1}\pgfmathsetmacro{\cy}{#2}%
    \pgfmathsetmacro{\w}{#3}\pgfmathsetmacro{\h}{#4}%
    \fill[#5] (\cx,\cy) rectangle ++(\w,\h);
    \fill[#5!80!white] (\cx,\cy+\h) -- ++(\archdx,\archdy) -- ++(\w,0) -- ++(-\archdx,-\archdy) -- cycle;
    \fill[#5!60!black] (\cx+\w,\cy) -- ++(\archdx,\archdy) -- ++(0,\h) -- ++(-\archdx,-\archdy) -- cycle;
    \draw[black,thin] (\cx,\cy) rectangle ++(\w,\h);
    \draw[black,thin] (\cx,\cy+\h) -- ++(\archdx,\archdy) -- ++(\w,0) -- ++(-\archdx,-\archdy) -- cycle;
    \draw[black,thin] (\cx+\w,\cy) -- ++(\archdx,\archdy) -- ++(0,\h) -- ++(-\archdx,-\archdy) -- cycle;
    \node[align=center] at (\cx+\w/2,\cy+\h/2) {#6};
  \end{scope}
}

% ---------------------------------------------------------------
% \archfunnel{x}{y}{h}{wwide}{wnarrow}{dir}{color}{label}
%   dir=1  : wide on left, narrow on right  (encoder)
%   dir=-1 : narrow on left, wide on right  (decoder)
% ---------------------------------------------------------------
\newcommand{\archfunnel}[8]{%
  \begin{scope}
    \pgfmathsetmacro{\cx}{#1}\pgfmathsetmacro{\cy}{#2}\pgfmathsetmacro{\h}{#3}%
    \pgfmathsetmacro{\ww}{#4}\pgfmathsetmacro{\wn}{#5}%
    \pgfmathsetmacro{\inset}{(\ww-\wn)/2}%
    \ifnum#6=1
      \fill[#7] (\cx,\cy) -- (\cx+\ww,\cy+\inset) -- (\cx+\ww,\cy+\h-\inset) -- (\cx,\cy+\h) -- cycle;
      \draw[black,thin] (\cx,\cy) -- (\cx+\ww,\cy+\inset) -- (\cx+\ww,\cy+\h-\inset) -- (\cx,\cy+\h) -- cycle;
    \else
      \fill[#7] (\cx,\cy+\inset) -- (\cx+\ww,\cy) -- (\cx+\ww,\cy+\h) -- (\cx,\cy+\h-\inset) -- cycle;
      \draw[black,thin] (\cx,\cy+\inset) -- (\cx+\ww,\cy) -- (\cx+\ww,\cy+\h) -- (\cx,\cy+\h-\inset) -- cycle;
    \fi
    \node[align=center] at (\cx+\ww/2,\cy+\h/2) {#8};
  \end{scope}
}

% ---------------------------------------------------------------
% \archbowtie{x}{y}{w}{h}{pinch}{color}{label} -- Medverse-style
% single-row "dual U-Net" hourglass (funnel-in then funnel-out).
% ---------------------------------------------------------------
\newcommand{\archbowtie}[7]{%
  \begin{scope}
    \pgfmathsetmacro{\cx}{#1}\pgfmathsetmacro{\cy}{#2}%
    \pgfmathsetmacro{\w}{#3}\pgfmathsetmacro{\h}{#4}\pgfmathsetmacro{\p}{#5}%
    \pgfmathsetmacro{\inset}{(\h-\p)/2}%
    \fill[#6] (\cx,\cy) -- (\cx+\w/2,\cy+\inset) -- (\cx+\w,\cy) -- (\cx+\w,\cy+\h)
      -- (\cx+\w/2,\cy+\h-\inset) -- (\cx,\cy+\h) -- cycle;
    \draw[black,thin] (\cx,\cy) -- (\cx+\w/2,\cy+\inset) -- (\cx+\w,\cy) -- (\cx+\w,\cy+\h)
      -- (\cx+\w/2,\cy+\h-\inset) -- (\cx,\cy+\h) -- cycle;
    \node[align=center] at (\cx+\w/2,\cy+\h/2) {#7};
  \end{scope}
}

% ---------------------------------------------------------------
% \archtokengrid{x}{y}{rows}{cols}{cell}{color} -- grid of patch tokens
% ---------------------------------------------------------------
\newcommand{\archtokengrid}[6]{%
  \begin{scope}
    \pgfmathsetmacro{\cx}{#1}\pgfmathsetmacro{\cy}{#2}\pgfmathsetmacro{\cell}{#5}%
    \pgfmathtruncatemacro{\rmax}{#3-1}%
    \pgfmathtruncatemacro{\cmax}{#4-1}%
    \foreach \r in {0,...,\rmax}{%
      \foreach \c in {0,...,\cmax}{%
        \fill[#6] ({\cx+\c*\cell},{\cy+\r*\cell}) rectangle ++(\cell*0.82,\cell*0.82);
        \draw[black,thin] ({\cx+\c*\cell},{\cy+\r*\cell}) rectangle ++(\cell*0.82,\cell*0.82);
      }%
    }%
  \end{scope}
}

% ---------------------------------------------------------------
% \archarrow{start}{end}{color} -- block arrow between stages
% ---------------------------------------------------------------
% #1 and #2 are full coordinates including parentheses, e.g. (1,2)
\newcommand{\archarrow}[3]{%
  \draw[-{Latex[length=2.4mm,width=2mm]},line width=2.2pt,#3] #1 -- #2;
}

% ---------------------------------------------------------------
% Dataflow-arrow vocabulary -- opt-in via \archarrowset{compact|detailed},
% called once near the top of a figure (after % Shared TikZ macros for the architecture figures (Medverse vs. Ours).
% Load with \input{arch-macros.tex} after \usepackage{tikz,ifthen} and
% \usetikzlibrary{arrows.meta,calc,positioning}. All macro names are
% prefixed `arch` to avoid clashing with pgf/tikz internals.

% ---- palette (approximating the PowerPoint drafts) ----
\definecolor{archTargetGreen}{HTML}{9BC7AC}
\definecolor{archContextPink}{HTML}{E8B8B0}
\definecolor{archAttnBlue}{HTML}{BBD6EC}
\definecolor{archEncGray}{HTML}{E4E4E4}
\definecolor{archArrowPale}{HTML}{CFE6D6}
\definecolor{archCtBg}{HTML}{2B2B2B}
\definecolor{archCtBlob}{HTML}{8C8C8C}
\definecolor{archCtxYellow}{HTML}{E8D24C}
% darker, more saturated siblings of archTargetGreen/archContextPink --
% pastel fills read poorly as 1pt strokes, and mask/label content needs to
% read as "which volume owns this" at a glance, so masks/labels and any
% line (arrow, brace) touching that volume's row use these instead of the
% pastel role color or the old flat archMaskGold.
\definecolor{archTargetAccent}{HTML}{2F6B4A}  % target: darker green
\definecolor{archContextAccent}{HTML}{A8453E} % context: darker red

% fake-isometric depth skew shared by every cube
\newcommand{\archdx}{0.30}
\newcommand{\archdy}{0.20}

% ---- shared text-role macros ----
% Every arch_*.tex figure is compiled standalone at the SAME base size
% (documentclass[tikz,border=4pt,12pt]{standalone}, matching the thesis
% body's 12pt), so \normalsize etc. mean the same absolute point size in
% every figure. Using these three named roles instead of ad hoc \Large /
% \bfseries\normalsize / \scriptsize per figure keeps "a box title" (etc.)
% the same size everywhere it appears, regardless of which figure it's in
% or how that figure's canvas happens to be scaled by \includegraphics.
%   \archtitle    -- figure/method/panel titles, the "xL" repeat marker
%   \archboxlabel -- sub-block header labels (Fuse, Encoder, attention, ...)
%   \archnote     -- small secondary annotations (placeholder captions)
% Plain, unlabeled node text (cube/token labels, row/lane labels) stays at
% the bare \normalsize default -- no macro needed for that role.
%
% DECLARATIONS, not argument-taking macros -- deliberately, so they can
% precede a raw `\\` inside a tikz node's text (as `{\archboxlabel line
% one\\line two}`) without hiding that `\\` one macro-expansion level deep.
% A `\\` buried inside a wrapped {\archboxlabel{...}} argument breaks
% TikZ's own align=center line-splitting (it scans for `\\` before the
% argument expands) -- use these the same way `\bfseries`/`\Large` are
% already used directly inline elsewhere in these figures.
\newcommand{\archtitle}{\Large\bfseries}
\newcommand{\archboxlabel}{\bfseries\normalsize}
\newcommand{\archnote}{\scriptsize}

% ---------------------------------------------------------------
% \archvolcube{x}{y}{w}{h}{topcolor}{content}{label}
%   content: ct | ctmask | predmask | <flat color name>
% ---------------------------------------------------------------
\newcommand{\archvolcube}[7]{%
  \begin{scope}
    \pgfmathsetmacro{\cx}{#1}\pgfmathsetmacro{\cy}{#2}%
    \pgfmathsetmacro{\w}{#3}\pgfmathsetmacro{\h}{#4}%
    \fill[#5] (\cx,\cy+\h) -- ++(\archdx,\archdy) -- ++(\w,0) -- ++(-\archdx,-\archdy) -- cycle;
    \fill[#5!60!black] (\cx+\w,\cy) -- ++(\archdx,\archdy) -- ++(0,\h) -- ++(-\archdx,-\archdy) -- cycle;
    \begin{scope}
      \clip (\cx,\cy) rectangle ++(\w,\h);
      \ifthenelse{\equal{#6}{ct}}{%
        \fill[archCtBg] (\cx,\cy) rectangle ++(\w,\h);
        \fill[archCtBlob] plot[smooth cycle] coordinates
          {(\cx+.25*\w,\cy+.55*\h) (\cx+.45*\w,\cy+.75*\h) (\cx+.7*\w,\cy+.6*\h)
           (\cx+.6*\w,\cy+.3*\h) (\cx+.35*\w,\cy+.25*\h)};
      }{\ifthenelse{\equal{#6}{ctmask}}{%
        \fill[archCtBg] (\cx,\cy) rectangle ++(\w,\h);
        \fill[archCtBlob] plot[smooth cycle] coordinates
          {(\cx+.25*\w,\cy+.55*\h) (\cx+.45*\w,\cy+.75*\h) (\cx+.7*\w,\cy+.6*\h)
           (\cx+.6*\w,\cy+.3*\h) (\cx+.35*\w,\cy+.25*\h)};
        \fill[archCtxYellow] plot[smooth cycle] coordinates
          {(\cx+.42*\w,\cy+.42*\h) (\cx+.52*\w,\cy+.5*\h) (\cx+.5*\w,\cy+.35*\h) (\cx+.44*\w,\cy+.3*\h)};
      }{\ifthenelse{\equal{#6}{predmask}}{%
        \fill[black] (\cx,\cy) rectangle ++(\w,\h);
        \fill[white] plot[smooth cycle] coordinates
          {(\cx+.4*\w,\cy+.65*\h) (\cx+.6*\w,\cy+.7*\h) (\cx+.62*\w,\cy+.45*\h)
           (\cx+.48*\w,\cy+.3*\h) (\cx+.38*\w,\cy+.4*\h)};
      }{%
        \fill[#6] (\cx,\cy) rectangle ++(\w,\h);
      }}}%
    \end{scope}
    \draw[black,thin] (\cx,\cy) rectangle ++(\w,\h);
    \draw[black,thin] (\cx,\cy+\h) -- ++(\archdx,\archdy) -- ++(\w,0) -- ++(-\archdx,-\archdy) -- cycle;
    \draw[black,thin] (\cx+\w,\cy) -- ++(\archdx,\archdy) -- ++(0,\h) -- ++(-\archdx,-\archdy) -- cycle;
    \node[anchor=north,align=center] at (\cx+\w/2,\cy-0.1) {#7};
  \end{scope}
}

% ---------------------------------------------------------------
% \archvolcubeimg{x}{y}{w}{h}{topcolor}{imgpath}{label} -- same cube as
% \archvolcube, but the front face shows a real raster image (clipped to
% exactly fill the face, aspect ratio not preserved) instead of one of
% \archvolcube's procedural ct/ctmask/predmask blobs. Needs graphicx
% loaded by the including figure. A separate macro rather than a new
% \archvolcube content mode, so existing ct/ctmask/predmask figures are
% untouched.
% ---------------------------------------------------------------
\newcommand{\archvolcubeimg}[7]{%
  \begin{scope}
    \pgfmathsetmacro{\cx}{#1}\pgfmathsetmacro{\cy}{#2}%
    \pgfmathsetmacro{\w}{#3}\pgfmathsetmacro{\h}{#4}%
    \fill[#5] (\cx,\cy+\h) -- ++(\archdx,\archdy) -- ++(\w,0) -- ++(-\archdx,-\archdy) -- cycle;
    \fill[#5!60!black] (\cx+\w,\cy) -- ++(\archdx,\archdy) -- ++(0,\h) -- ++(-\archdx,-\archdy) -- cycle;
    \begin{scope}
      \clip (\cx,\cy) rectangle ++(\w,\h);
      \node[anchor=south west, inner sep=0pt] at (\cx,\cy)
        {\includegraphics[width=\w cm,height=\h cm]{#6}};
    \end{scope}
    \draw[black,thin] (\cx,\cy) rectangle ++(\w,\h);
    \draw[black,thin] (\cx,\cy+\h) -- ++(\archdx,\archdy) -- ++(\w,0) -- ++(-\archdx,-\archdy) -- cycle;
    \draw[black,thin] (\cx+\w,\cy) -- ++(\archdx,\archdy) -- ++(0,\h) -- ++(-\archdx,-\archdy) -- cycle;
    \node[anchor=north,align=center] at (\cx+\w/2,\cy-0.1) {#7};
  \end{scope}
}

% ---------------------------------------------------------------
% \archfeatcube{x}{y}{w}{h}{color}{label} -- flat-colored feature volume
% ---------------------------------------------------------------
\newcommand{\archfeatcube}[6]{%
  \begin{scope}
    \pgfmathsetmacro{\cx}{#1}\pgfmathsetmacro{\cy}{#2}%
    \pgfmathsetmacro{\w}{#3}\pgfmathsetmacro{\h}{#4}%
    \fill[#5] (\cx,\cy) rectangle ++(\w,\h);
    \fill[#5!80!white] (\cx,\cy+\h) -- ++(\archdx,\archdy) -- ++(\w,0) -- ++(-\archdx,-\archdy) -- cycle;
    \fill[#5!60!black] (\cx+\w,\cy) -- ++(\archdx,\archdy) -- ++(0,\h) -- ++(-\archdx,-\archdy) -- cycle;
    \draw[black,thin] (\cx,\cy) rectangle ++(\w,\h);
    \draw[black,thin] (\cx,\cy+\h) -- ++(\archdx,\archdy) -- ++(\w,0) -- ++(-\archdx,-\archdy) -- cycle;
    \draw[black,thin] (\cx+\w,\cy) -- ++(\archdx,\archdy) -- ++(0,\h) -- ++(-\archdx,-\archdy) -- cycle;
    \node[align=center] at (\cx+\w/2,\cy+\h/2) {#6};
  \end{scope}
}

% ---------------------------------------------------------------
% \archfunnel{x}{y}{h}{wwide}{wnarrow}{dir}{color}{label}
%   dir=1  : wide on left, narrow on right  (encoder)
%   dir=-1 : narrow on left, wide on right  (decoder)
% ---------------------------------------------------------------
\newcommand{\archfunnel}[8]{%
  \begin{scope}
    \pgfmathsetmacro{\cx}{#1}\pgfmathsetmacro{\cy}{#2}\pgfmathsetmacro{\h}{#3}%
    \pgfmathsetmacro{\ww}{#4}\pgfmathsetmacro{\wn}{#5}%
    \pgfmathsetmacro{\inset}{(\ww-\wn)/2}%
    \ifnum#6=1
      \fill[#7] (\cx,\cy) -- (\cx+\ww,\cy+\inset) -- (\cx+\ww,\cy+\h-\inset) -- (\cx,\cy+\h) -- cycle;
      \draw[black,thin] (\cx,\cy) -- (\cx+\ww,\cy+\inset) -- (\cx+\ww,\cy+\h-\inset) -- (\cx,\cy+\h) -- cycle;
    \else
      \fill[#7] (\cx,\cy+\inset) -- (\cx+\ww,\cy) -- (\cx+\ww,\cy+\h) -- (\cx,\cy+\h-\inset) -- cycle;
      \draw[black,thin] (\cx,\cy+\inset) -- (\cx+\ww,\cy) -- (\cx+\ww,\cy+\h) -- (\cx,\cy+\h-\inset) -- cycle;
    \fi
    \node[align=center] at (\cx+\ww/2,\cy+\h/2) {#8};
  \end{scope}
}

% ---------------------------------------------------------------
% \archbowtie{x}{y}{w}{h}{pinch}{color}{label} -- Medverse-style
% single-row "dual U-Net" hourglass (funnel-in then funnel-out).
% ---------------------------------------------------------------
\newcommand{\archbowtie}[7]{%
  \begin{scope}
    \pgfmathsetmacro{\cx}{#1}\pgfmathsetmacro{\cy}{#2}%
    \pgfmathsetmacro{\w}{#3}\pgfmathsetmacro{\h}{#4}\pgfmathsetmacro{\p}{#5}%
    \pgfmathsetmacro{\inset}{(\h-\p)/2}%
    \fill[#6] (\cx,\cy) -- (\cx+\w/2,\cy+\inset) -- (\cx+\w,\cy) -- (\cx+\w,\cy+\h)
      -- (\cx+\w/2,\cy+\h-\inset) -- (\cx,\cy+\h) -- cycle;
    \draw[black,thin] (\cx,\cy) -- (\cx+\w/2,\cy+\inset) -- (\cx+\w,\cy) -- (\cx+\w,\cy+\h)
      -- (\cx+\w/2,\cy+\h-\inset) -- (\cx,\cy+\h) -- cycle;
    \node[align=center] at (\cx+\w/2,\cy+\h/2) {#7};
  \end{scope}
}

% ---------------------------------------------------------------
% \archtokengrid{x}{y}{rows}{cols}{cell}{color} -- grid of patch tokens
% ---------------------------------------------------------------
\newcommand{\archtokengrid}[6]{%
  \begin{scope}
    \pgfmathsetmacro{\cx}{#1}\pgfmathsetmacro{\cy}{#2}\pgfmathsetmacro{\cell}{#5}%
    \pgfmathtruncatemacro{\rmax}{#3-1}%
    \pgfmathtruncatemacro{\cmax}{#4-1}%
    \foreach \r in {0,...,\rmax}{%
      \foreach \c in {0,...,\cmax}{%
        \fill[#6] ({\cx+\c*\cell},{\cy+\r*\cell}) rectangle ++(\cell*0.82,\cell*0.82);
        \draw[black,thin] ({\cx+\c*\cell},{\cy+\r*\cell}) rectangle ++(\cell*0.82,\cell*0.82);
      }%
    }%
  \end{scope}
}

% ---------------------------------------------------------------
% \archarrow{start}{end}{color} -- block arrow between stages
% ---------------------------------------------------------------
% #1 and #2 are full coordinates including parentheses, e.g. (1,2)
\newcommand{\archarrow}[3]{%
  \draw[-{Latex[length=2.4mm,width=2mm]},line width=2.2pt,#3] #1 -- #2;
}

% ---------------------------------------------------------------
% Dataflow-arrow vocabulary -- opt-in via \archarrowset{compact|detailed},
% called once near the top of a figure (after \input{arch-macros.tex}).
% \archarrow keeps its classic small size by default for any arch_*.tex
% figure that never calls \archarrowset -- this block is purely additive.
%   archflowstyle (tikz key, NOT a macro -- use bare, no backslash) --
%                                  current draw options (head + line
%                                  width), for a hand-rolled elbow \draw
%                                  that \archarrow's straight "--" can't
%                                  express (e.g. "-|" routing). Must be a
%                                  /.style key rather than a \def'd macro:
%                                  pgf's arrow-tip parser doesn't split on
%                                  a comma embedded in a plain macro's
%                                  expansion when used as a \draw option.
%   \archflowNeutral/Target/Context -- the only 3 colors a plain dataflow
%                                  arrow should use: generic same-role
%                                  passthrough; target's own real output
%                                  leaving the decoder; a real label/mask
%                                  tensor entering fusion or attention
%                                  (context's ground truth, or target's
%                                  own real prior fed back from a
%                                  previous cascade level -- same role
%                                  either way: real mask data, not a
%                                  feature)
%   \archflow{p1}{p2}{role}     -- \archarrow wrapper keyed by role name
%                                  (role in {neutral, target, context})
% ---------------------------------------------------------------
\newcommand{\archarrowset}[1]{%
  \ifthenelse{\equal{#1}{detailed}}{%
    \tikzset{archflowstyle/.style={-{Latex[length=3.6mm,width=3.2mm]},line width=3.2pt}}%
  }{%
    \tikzset{archflowstyle/.style={-{Latex[length=2.4mm,width=2mm]},line width=2.2pt}}%
  }%
  \renewcommand{\archarrow}[3]{\draw[archflowstyle,##3] ##1 -- ##2;}%
}
\newcommand{\archflowNeutral}{black!30}
\newcommand{\archflowTarget}{archTargetAccent}
\newcommand{\archflowContext}{archContextAccent}
\newcommand{\archflow}[3]{%
  \ifthenelse{\equal{#3}{target}}{\archarrow{#1}{#2}{\archflowTarget}}{%
  \ifthenelse{\equal{#3}{context}}{\archarrow{#1}{#2}{\archflowContext}}{%
  \archarrow{#1}{#2}{\archflowNeutral}}}%
}
).
% \archarrow keeps its classic small size by default for any arch_*.tex
% figure that never calls \archarrowset -- this block is purely additive.
%   archflowstyle (tikz key, NOT a macro -- use bare, no backslash) --
%                                  current draw options (head + line
%                                  width), for a hand-rolled elbow \draw
%                                  that \archarrow's straight "--" can't
%                                  express (e.g. "-|" routing). Must be a
%                                  /.style key rather than a \def'd macro:
%                                  pgf's arrow-tip parser doesn't split on
%                                  a comma embedded in a plain macro's
%                                  expansion when used as a \draw option.
%   \archflowNeutral/Target/Context -- the only 3 colors a plain dataflow
%                                  arrow should use: generic same-role
%                                  passthrough; target's own real output
%                                  leaving the decoder; a real label/mask
%                                  tensor entering fusion or attention
%                                  (context's ground truth, or target's
%                                  own real prior fed back from a
%                                  previous cascade level -- same role
%                                  either way: real mask data, not a
%                                  feature)
%   \archflow{p1}{p2}{role}     -- \archarrow wrapper keyed by role name
%                                  (role in {neutral, target, context})
% ---------------------------------------------------------------
\newcommand{\archarrowset}[1]{%
  \ifthenelse{\equal{#1}{detailed}}{%
    \tikzset{archflowstyle/.style={-{Latex[length=3.6mm,width=3.2mm]},line width=3.2pt}}%
  }{%
    \tikzset{archflowstyle/.style={-{Latex[length=2.4mm,width=2mm]},line width=2.2pt}}%
  }%
  \renewcommand{\archarrow}[3]{\draw[archflowstyle,##3] ##1 -- ##2;}%
}
\newcommand{\archflowNeutral}{black!30}
\newcommand{\archflowTarget}{archTargetAccent}
\newcommand{\archflowContext}{archContextAccent}
\newcommand{\archflow}[3]{%
  \ifthenelse{\equal{#3}{target}}{\archarrow{#1}{#2}{\archflowTarget}}{%
  \ifthenelse{\equal{#3}{context}}{\archarrow{#1}{#2}{\archflowContext}}{%
  \archarrow{#1}{#2}{\archflowNeutral}}}%
}
 after \usepackage{tikz,ifthen} and
% \usetikzlibrary{arrows.meta,calc,positioning}. All macro names are
% prefixed `arch` to avoid clashing with pgf/tikz internals.

% ---- palette (approximating the PowerPoint drafts) ----
\definecolor{archTargetGreen}{HTML}{9BC7AC}
\definecolor{archContextPink}{HTML}{E8B8B0}
\definecolor{archAttnBlue}{HTML}{BBD6EC}
\definecolor{archEncGray}{HTML}{E4E4E4}
\definecolor{archArrowPale}{HTML}{CFE6D6}
\definecolor{archCtBg}{HTML}{2B2B2B}
\definecolor{archCtBlob}{HTML}{8C8C8C}
\definecolor{archCtxYellow}{HTML}{E8D24C}
% darker, more saturated siblings of archTargetGreen/archContextPink --
% pastel fills read poorly as 1pt strokes, and mask/label content needs to
% read as "which volume owns this" at a glance, so masks/labels and any
% line (arrow, brace) touching that volume's row use these instead of the
% pastel role color or the old flat archMaskGold.
\definecolor{archTargetAccent}{HTML}{2F6B4A}  % target: darker green
\definecolor{archContextAccent}{HTML}{A8453E} % context: darker red

% fake-isometric depth skew shared by every cube
\newcommand{\archdx}{0.30}
\newcommand{\archdy}{0.20}

% ---- shared text-role macros ----
% Every arch_*.tex figure is compiled standalone at the SAME base size
% (documentclass[tikz,border=4pt,12pt]{standalone}, matching the thesis
% body's 12pt), so \normalsize etc. mean the same absolute point size in
% every figure. Using these three named roles instead of ad hoc \Large /
% \bfseries\normalsize / \scriptsize per figure keeps "a box title" (etc.)
% the same size everywhere it appears, regardless of which figure it's in
% or how that figure's canvas happens to be scaled by \includegraphics.
%   \archtitle    -- figure/method/panel titles, the "xL" repeat marker
%   \archboxlabel -- sub-block header labels (Fuse, Encoder, attention, ...)
%   \archnote     -- small secondary annotations (placeholder captions)
% Plain, unlabeled node text (cube/token labels, row/lane labels) stays at
% the bare \normalsize default -- no macro needed for that role.
%
% DECLARATIONS, not argument-taking macros -- deliberately, so they can
% precede a raw `\\` inside a tikz node's text (as `{\archboxlabel line
% one\\line two}`) without hiding that `\\` one macro-expansion level deep.
% A `\\` buried inside a wrapped {\archboxlabel{...}} argument breaks
% TikZ's own align=center line-splitting (it scans for `\\` before the
% argument expands) -- use these the same way `\bfseries`/`\Large` are
% already used directly inline elsewhere in these figures.
\newcommand{\archtitle}{\Large\bfseries}
\newcommand{\archboxlabel}{\bfseries\normalsize}
\newcommand{\archnote}{\scriptsize}

% ---------------------------------------------------------------
% \archvolcube{x}{y}{w}{h}{topcolor}{content}{label}
%   content: ct | ctmask | predmask | <flat color name>
% ---------------------------------------------------------------
\newcommand{\archvolcube}[7]{%
  \begin{scope}
    \pgfmathsetmacro{\cx}{#1}\pgfmathsetmacro{\cy}{#2}%
    \pgfmathsetmacro{\w}{#3}\pgfmathsetmacro{\h}{#4}%
    \fill[#5] (\cx,\cy+\h) -- ++(\archdx,\archdy) -- ++(\w,0) -- ++(-\archdx,-\archdy) -- cycle;
    \fill[#5!60!black] (\cx+\w,\cy) -- ++(\archdx,\archdy) -- ++(0,\h) -- ++(-\archdx,-\archdy) -- cycle;
    \begin{scope}
      \clip (\cx,\cy) rectangle ++(\w,\h);
      \ifthenelse{\equal{#6}{ct}}{%
        \fill[archCtBg] (\cx,\cy) rectangle ++(\w,\h);
        \fill[archCtBlob] plot[smooth cycle] coordinates
          {(\cx+.25*\w,\cy+.55*\h) (\cx+.45*\w,\cy+.75*\h) (\cx+.7*\w,\cy+.6*\h)
           (\cx+.6*\w,\cy+.3*\h) (\cx+.35*\w,\cy+.25*\h)};
      }{\ifthenelse{\equal{#6}{ctmask}}{%
        \fill[archCtBg] (\cx,\cy) rectangle ++(\w,\h);
        \fill[archCtBlob] plot[smooth cycle] coordinates
          {(\cx+.25*\w,\cy+.55*\h) (\cx+.45*\w,\cy+.75*\h) (\cx+.7*\w,\cy+.6*\h)
           (\cx+.6*\w,\cy+.3*\h) (\cx+.35*\w,\cy+.25*\h)};
        \fill[archCtxYellow] plot[smooth cycle] coordinates
          {(\cx+.42*\w,\cy+.42*\h) (\cx+.52*\w,\cy+.5*\h) (\cx+.5*\w,\cy+.35*\h) (\cx+.44*\w,\cy+.3*\h)};
      }{\ifthenelse{\equal{#6}{predmask}}{%
        \fill[black] (\cx,\cy) rectangle ++(\w,\h);
        \fill[white] plot[smooth cycle] coordinates
          {(\cx+.4*\w,\cy+.65*\h) (\cx+.6*\w,\cy+.7*\h) (\cx+.62*\w,\cy+.45*\h)
           (\cx+.48*\w,\cy+.3*\h) (\cx+.38*\w,\cy+.4*\h)};
      }{%
        \fill[#6] (\cx,\cy) rectangle ++(\w,\h);
      }}}%
    \end{scope}
    \draw[black,thin] (\cx,\cy) rectangle ++(\w,\h);
    \draw[black,thin] (\cx,\cy+\h) -- ++(\archdx,\archdy) -- ++(\w,0) -- ++(-\archdx,-\archdy) -- cycle;
    \draw[black,thin] (\cx+\w,\cy) -- ++(\archdx,\archdy) -- ++(0,\h) -- ++(-\archdx,-\archdy) -- cycle;
    \node[anchor=north,align=center] at (\cx+\w/2,\cy-0.1) {#7};
  \end{scope}
}

% ---------------------------------------------------------------
% \archvolcubeimg{x}{y}{w}{h}{topcolor}{imgpath}{label} -- same cube as
% \archvolcube, but the front face shows a real raster image (clipped to
% exactly fill the face, aspect ratio not preserved) instead of one of
% \archvolcube's procedural ct/ctmask/predmask blobs. Needs graphicx
% loaded by the including figure. A separate macro rather than a new
% \archvolcube content mode, so existing ct/ctmask/predmask figures are
% untouched.
% ---------------------------------------------------------------
\newcommand{\archvolcubeimg}[7]{%
  \begin{scope}
    \pgfmathsetmacro{\cx}{#1}\pgfmathsetmacro{\cy}{#2}%
    \pgfmathsetmacro{\w}{#3}\pgfmathsetmacro{\h}{#4}%
    \fill[#5] (\cx,\cy+\h) -- ++(\archdx,\archdy) -- ++(\w,0) -- ++(-\archdx,-\archdy) -- cycle;
    \fill[#5!60!black] (\cx+\w,\cy) -- ++(\archdx,\archdy) -- ++(0,\h) -- ++(-\archdx,-\archdy) -- cycle;
    \begin{scope}
      \clip (\cx,\cy) rectangle ++(\w,\h);
      \node[anchor=south west, inner sep=0pt] at (\cx,\cy)
        {\includegraphics[width=\w cm,height=\h cm]{#6}};
    \end{scope}
    \draw[black,thin] (\cx,\cy) rectangle ++(\w,\h);
    \draw[black,thin] (\cx,\cy+\h) -- ++(\archdx,\archdy) -- ++(\w,0) -- ++(-\archdx,-\archdy) -- cycle;
    \draw[black,thin] (\cx+\w,\cy) -- ++(\archdx,\archdy) -- ++(0,\h) -- ++(-\archdx,-\archdy) -- cycle;
    \node[anchor=north,align=center] at (\cx+\w/2,\cy-0.1) {#7};
  \end{scope}
}

% ---------------------------------------------------------------
% \archfeatcube{x}{y}{w}{h}{color}{label} -- flat-colored feature volume
% ---------------------------------------------------------------
\newcommand{\archfeatcube}[6]{%
  \begin{scope}
    \pgfmathsetmacro{\cx}{#1}\pgfmathsetmacro{\cy}{#2}%
    \pgfmathsetmacro{\w}{#3}\pgfmathsetmacro{\h}{#4}%
    \fill[#5] (\cx,\cy) rectangle ++(\w,\h);
    \fill[#5!80!white] (\cx,\cy+\h) -- ++(\archdx,\archdy) -- ++(\w,0) -- ++(-\archdx,-\archdy) -- cycle;
    \fill[#5!60!black] (\cx+\w,\cy) -- ++(\archdx,\archdy) -- ++(0,\h) -- ++(-\archdx,-\archdy) -- cycle;
    \draw[black,thin] (\cx,\cy) rectangle ++(\w,\h);
    \draw[black,thin] (\cx,\cy+\h) -- ++(\archdx,\archdy) -- ++(\w,0) -- ++(-\archdx,-\archdy) -- cycle;
    \draw[black,thin] (\cx+\w,\cy) -- ++(\archdx,\archdy) -- ++(0,\h) -- ++(-\archdx,-\archdy) -- cycle;
    \node[align=center] at (\cx+\w/2,\cy+\h/2) {#6};
  \end{scope}
}

% ---------------------------------------------------------------
% \archfunnel{x}{y}{h}{wwide}{wnarrow}{dir}{color}{label}
%   dir=1  : wide on left, narrow on right  (encoder)
%   dir=-1 : narrow on left, wide on right  (decoder)
% ---------------------------------------------------------------
\newcommand{\archfunnel}[8]{%
  \begin{scope}
    \pgfmathsetmacro{\cx}{#1}\pgfmathsetmacro{\cy}{#2}\pgfmathsetmacro{\h}{#3}%
    \pgfmathsetmacro{\ww}{#4}\pgfmathsetmacro{\wn}{#5}%
    \pgfmathsetmacro{\inset}{(\ww-\wn)/2}%
    \ifnum#6=1
      \fill[#7] (\cx,\cy) -- (\cx+\ww,\cy+\inset) -- (\cx+\ww,\cy+\h-\inset) -- (\cx,\cy+\h) -- cycle;
      \draw[black,thin] (\cx,\cy) -- (\cx+\ww,\cy+\inset) -- (\cx+\ww,\cy+\h-\inset) -- (\cx,\cy+\h) -- cycle;
    \else
      \fill[#7] (\cx,\cy+\inset) -- (\cx+\ww,\cy) -- (\cx+\ww,\cy+\h) -- (\cx,\cy+\h-\inset) -- cycle;
      \draw[black,thin] (\cx,\cy+\inset) -- (\cx+\ww,\cy) -- (\cx+\ww,\cy+\h) -- (\cx,\cy+\h-\inset) -- cycle;
    \fi
    \node[align=center] at (\cx+\ww/2,\cy+\h/2) {#8};
  \end{scope}
}

% ---------------------------------------------------------------
% \archbowtie{x}{y}{w}{h}{pinch}{color}{label} -- Medverse-style
% single-row "dual U-Net" hourglass (funnel-in then funnel-out).
% ---------------------------------------------------------------
\newcommand{\archbowtie}[7]{%
  \begin{scope}
    \pgfmathsetmacro{\cx}{#1}\pgfmathsetmacro{\cy}{#2}%
    \pgfmathsetmacro{\w}{#3}\pgfmathsetmacro{\h}{#4}\pgfmathsetmacro{\p}{#5}%
    \pgfmathsetmacro{\inset}{(\h-\p)/2}%
    \fill[#6] (\cx,\cy) -- (\cx+\w/2,\cy+\inset) -- (\cx+\w,\cy) -- (\cx+\w,\cy+\h)
      -- (\cx+\w/2,\cy+\h-\inset) -- (\cx,\cy+\h) -- cycle;
    \draw[black,thin] (\cx,\cy) -- (\cx+\w/2,\cy+\inset) -- (\cx+\w,\cy) -- (\cx+\w,\cy+\h)
      -- (\cx+\w/2,\cy+\h-\inset) -- (\cx,\cy+\h) -- cycle;
    \node[align=center] at (\cx+\w/2,\cy+\h/2) {#7};
  \end{scope}
}

% ---------------------------------------------------------------
% \archtokengrid{x}{y}{rows}{cols}{cell}{color} -- grid of patch tokens
% ---------------------------------------------------------------
\newcommand{\archtokengrid}[6]{%
  \begin{scope}
    \pgfmathsetmacro{\cx}{#1}\pgfmathsetmacro{\cy}{#2}\pgfmathsetmacro{\cell}{#5}%
    \pgfmathtruncatemacro{\rmax}{#3-1}%
    \pgfmathtruncatemacro{\cmax}{#4-1}%
    \foreach \r in {0,...,\rmax}{%
      \foreach \c in {0,...,\cmax}{%
        \fill[#6] ({\cx+\c*\cell},{\cy+\r*\cell}) rectangle ++(\cell*0.82,\cell*0.82);
        \draw[black,thin] ({\cx+\c*\cell},{\cy+\r*\cell}) rectangle ++(\cell*0.82,\cell*0.82);
      }%
    }%
  \end{scope}
}

% ---------------------------------------------------------------
% \archarrow{start}{end}{color} -- block arrow between stages
% ---------------------------------------------------------------
% #1 and #2 are full coordinates including parentheses, e.g. (1,2)
\newcommand{\archarrow}[3]{%
  \draw[-{Latex[length=2.4mm,width=2mm]},line width=2.2pt,#3] #1 -- #2;
}

% ---------------------------------------------------------------
% Dataflow-arrow vocabulary -- opt-in via \archarrowset{compact|detailed},
% called once near the top of a figure (after % Shared TikZ macros for the architecture figures (Medverse vs. Ours).
% Load with % Shared TikZ macros for the architecture figures (Medverse vs. Ours).
% Load with \input{arch-macros.tex} after \usepackage{tikz,ifthen} and
% \usetikzlibrary{arrows.meta,calc,positioning}. All macro names are
% prefixed `arch` to avoid clashing with pgf/tikz internals.

% ---- palette (approximating the PowerPoint drafts) ----
\definecolor{archTargetGreen}{HTML}{9BC7AC}
\definecolor{archContextPink}{HTML}{E8B8B0}
\definecolor{archAttnBlue}{HTML}{BBD6EC}
\definecolor{archEncGray}{HTML}{E4E4E4}
\definecolor{archArrowPale}{HTML}{CFE6D6}
\definecolor{archCtBg}{HTML}{2B2B2B}
\definecolor{archCtBlob}{HTML}{8C8C8C}
\definecolor{archCtxYellow}{HTML}{E8D24C}
% darker, more saturated siblings of archTargetGreen/archContextPink --
% pastel fills read poorly as 1pt strokes, and mask/label content needs to
% read as "which volume owns this" at a glance, so masks/labels and any
% line (arrow, brace) touching that volume's row use these instead of the
% pastel role color or the old flat archMaskGold.
\definecolor{archTargetAccent}{HTML}{2F6B4A}  % target: darker green
\definecolor{archContextAccent}{HTML}{A8453E} % context: darker red

% fake-isometric depth skew shared by every cube
\newcommand{\archdx}{0.30}
\newcommand{\archdy}{0.20}

% ---- shared text-role macros ----
% Every arch_*.tex figure is compiled standalone at the SAME base size
% (documentclass[tikz,border=4pt,12pt]{standalone}, matching the thesis
% body's 12pt), so \normalsize etc. mean the same absolute point size in
% every figure. Using these three named roles instead of ad hoc \Large /
% \bfseries\normalsize / \scriptsize per figure keeps "a box title" (etc.)
% the same size everywhere it appears, regardless of which figure it's in
% or how that figure's canvas happens to be scaled by \includegraphics.
%   \archtitle    -- figure/method/panel titles, the "xL" repeat marker
%   \archboxlabel -- sub-block header labels (Fuse, Encoder, attention, ...)
%   \archnote     -- small secondary annotations (placeholder captions)
% Plain, unlabeled node text (cube/token labels, row/lane labels) stays at
% the bare \normalsize default -- no macro needed for that role.
%
% DECLARATIONS, not argument-taking macros -- deliberately, so they can
% precede a raw `\\` inside a tikz node's text (as `{\archboxlabel line
% one\\line two}`) without hiding that `\\` one macro-expansion level deep.
% A `\\` buried inside a wrapped {\archboxlabel{...}} argument breaks
% TikZ's own align=center line-splitting (it scans for `\\` before the
% argument expands) -- use these the same way `\bfseries`/`\Large` are
% already used directly inline elsewhere in these figures.
\newcommand{\archtitle}{\Large\bfseries}
\newcommand{\archboxlabel}{\bfseries\normalsize}
\newcommand{\archnote}{\scriptsize}

% ---------------------------------------------------------------
% \archvolcube{x}{y}{w}{h}{topcolor}{content}{label}
%   content: ct | ctmask | predmask | <flat color name>
% ---------------------------------------------------------------
\newcommand{\archvolcube}[7]{%
  \begin{scope}
    \pgfmathsetmacro{\cx}{#1}\pgfmathsetmacro{\cy}{#2}%
    \pgfmathsetmacro{\w}{#3}\pgfmathsetmacro{\h}{#4}%
    \fill[#5] (\cx,\cy+\h) -- ++(\archdx,\archdy) -- ++(\w,0) -- ++(-\archdx,-\archdy) -- cycle;
    \fill[#5!60!black] (\cx+\w,\cy) -- ++(\archdx,\archdy) -- ++(0,\h) -- ++(-\archdx,-\archdy) -- cycle;
    \begin{scope}
      \clip (\cx,\cy) rectangle ++(\w,\h);
      \ifthenelse{\equal{#6}{ct}}{%
        \fill[archCtBg] (\cx,\cy) rectangle ++(\w,\h);
        \fill[archCtBlob] plot[smooth cycle] coordinates
          {(\cx+.25*\w,\cy+.55*\h) (\cx+.45*\w,\cy+.75*\h) (\cx+.7*\w,\cy+.6*\h)
           (\cx+.6*\w,\cy+.3*\h) (\cx+.35*\w,\cy+.25*\h)};
      }{\ifthenelse{\equal{#6}{ctmask}}{%
        \fill[archCtBg] (\cx,\cy) rectangle ++(\w,\h);
        \fill[archCtBlob] plot[smooth cycle] coordinates
          {(\cx+.25*\w,\cy+.55*\h) (\cx+.45*\w,\cy+.75*\h) (\cx+.7*\w,\cy+.6*\h)
           (\cx+.6*\w,\cy+.3*\h) (\cx+.35*\w,\cy+.25*\h)};
        \fill[archCtxYellow] plot[smooth cycle] coordinates
          {(\cx+.42*\w,\cy+.42*\h) (\cx+.52*\w,\cy+.5*\h) (\cx+.5*\w,\cy+.35*\h) (\cx+.44*\w,\cy+.3*\h)};
      }{\ifthenelse{\equal{#6}{predmask}}{%
        \fill[black] (\cx,\cy) rectangle ++(\w,\h);
        \fill[white] plot[smooth cycle] coordinates
          {(\cx+.4*\w,\cy+.65*\h) (\cx+.6*\w,\cy+.7*\h) (\cx+.62*\w,\cy+.45*\h)
           (\cx+.48*\w,\cy+.3*\h) (\cx+.38*\w,\cy+.4*\h)};
      }{%
        \fill[#6] (\cx,\cy) rectangle ++(\w,\h);
      }}}%
    \end{scope}
    \draw[black,thin] (\cx,\cy) rectangle ++(\w,\h);
    \draw[black,thin] (\cx,\cy+\h) -- ++(\archdx,\archdy) -- ++(\w,0) -- ++(-\archdx,-\archdy) -- cycle;
    \draw[black,thin] (\cx+\w,\cy) -- ++(\archdx,\archdy) -- ++(0,\h) -- ++(-\archdx,-\archdy) -- cycle;
    \node[anchor=north,align=center] at (\cx+\w/2,\cy-0.1) {#7};
  \end{scope}
}

% ---------------------------------------------------------------
% \archvolcubeimg{x}{y}{w}{h}{topcolor}{imgpath}{label} -- same cube as
% \archvolcube, but the front face shows a real raster image (clipped to
% exactly fill the face, aspect ratio not preserved) instead of one of
% \archvolcube's procedural ct/ctmask/predmask blobs. Needs graphicx
% loaded by the including figure. A separate macro rather than a new
% \archvolcube content mode, so existing ct/ctmask/predmask figures are
% untouched.
% ---------------------------------------------------------------
\newcommand{\archvolcubeimg}[7]{%
  \begin{scope}
    \pgfmathsetmacro{\cx}{#1}\pgfmathsetmacro{\cy}{#2}%
    \pgfmathsetmacro{\w}{#3}\pgfmathsetmacro{\h}{#4}%
    \fill[#5] (\cx,\cy+\h) -- ++(\archdx,\archdy) -- ++(\w,0) -- ++(-\archdx,-\archdy) -- cycle;
    \fill[#5!60!black] (\cx+\w,\cy) -- ++(\archdx,\archdy) -- ++(0,\h) -- ++(-\archdx,-\archdy) -- cycle;
    \begin{scope}
      \clip (\cx,\cy) rectangle ++(\w,\h);
      \node[anchor=south west, inner sep=0pt] at (\cx,\cy)
        {\includegraphics[width=\w cm,height=\h cm]{#6}};
    \end{scope}
    \draw[black,thin] (\cx,\cy) rectangle ++(\w,\h);
    \draw[black,thin] (\cx,\cy+\h) -- ++(\archdx,\archdy) -- ++(\w,0) -- ++(-\archdx,-\archdy) -- cycle;
    \draw[black,thin] (\cx+\w,\cy) -- ++(\archdx,\archdy) -- ++(0,\h) -- ++(-\archdx,-\archdy) -- cycle;
    \node[anchor=north,align=center] at (\cx+\w/2,\cy-0.1) {#7};
  \end{scope}
}

% ---------------------------------------------------------------
% \archfeatcube{x}{y}{w}{h}{color}{label} -- flat-colored feature volume
% ---------------------------------------------------------------
\newcommand{\archfeatcube}[6]{%
  \begin{scope}
    \pgfmathsetmacro{\cx}{#1}\pgfmathsetmacro{\cy}{#2}%
    \pgfmathsetmacro{\w}{#3}\pgfmathsetmacro{\h}{#4}%
    \fill[#5] (\cx,\cy) rectangle ++(\w,\h);
    \fill[#5!80!white] (\cx,\cy+\h) -- ++(\archdx,\archdy) -- ++(\w,0) -- ++(-\archdx,-\archdy) -- cycle;
    \fill[#5!60!black] (\cx+\w,\cy) -- ++(\archdx,\archdy) -- ++(0,\h) -- ++(-\archdx,-\archdy) -- cycle;
    \draw[black,thin] (\cx,\cy) rectangle ++(\w,\h);
    \draw[black,thin] (\cx,\cy+\h) -- ++(\archdx,\archdy) -- ++(\w,0) -- ++(-\archdx,-\archdy) -- cycle;
    \draw[black,thin] (\cx+\w,\cy) -- ++(\archdx,\archdy) -- ++(0,\h) -- ++(-\archdx,-\archdy) -- cycle;
    \node[align=center] at (\cx+\w/2,\cy+\h/2) {#6};
  \end{scope}
}

% ---------------------------------------------------------------
% \archfunnel{x}{y}{h}{wwide}{wnarrow}{dir}{color}{label}
%   dir=1  : wide on left, narrow on right  (encoder)
%   dir=-1 : narrow on left, wide on right  (decoder)
% ---------------------------------------------------------------
\newcommand{\archfunnel}[8]{%
  \begin{scope}
    \pgfmathsetmacro{\cx}{#1}\pgfmathsetmacro{\cy}{#2}\pgfmathsetmacro{\h}{#3}%
    \pgfmathsetmacro{\ww}{#4}\pgfmathsetmacro{\wn}{#5}%
    \pgfmathsetmacro{\inset}{(\ww-\wn)/2}%
    \ifnum#6=1
      \fill[#7] (\cx,\cy) -- (\cx+\ww,\cy+\inset) -- (\cx+\ww,\cy+\h-\inset) -- (\cx,\cy+\h) -- cycle;
      \draw[black,thin] (\cx,\cy) -- (\cx+\ww,\cy+\inset) -- (\cx+\ww,\cy+\h-\inset) -- (\cx,\cy+\h) -- cycle;
    \else
      \fill[#7] (\cx,\cy+\inset) -- (\cx+\ww,\cy) -- (\cx+\ww,\cy+\h) -- (\cx,\cy+\h-\inset) -- cycle;
      \draw[black,thin] (\cx,\cy+\inset) -- (\cx+\ww,\cy) -- (\cx+\ww,\cy+\h) -- (\cx,\cy+\h-\inset) -- cycle;
    \fi
    \node[align=center] at (\cx+\ww/2,\cy+\h/2) {#8};
  \end{scope}
}

% ---------------------------------------------------------------
% \archbowtie{x}{y}{w}{h}{pinch}{color}{label} -- Medverse-style
% single-row "dual U-Net" hourglass (funnel-in then funnel-out).
% ---------------------------------------------------------------
\newcommand{\archbowtie}[7]{%
  \begin{scope}
    \pgfmathsetmacro{\cx}{#1}\pgfmathsetmacro{\cy}{#2}%
    \pgfmathsetmacro{\w}{#3}\pgfmathsetmacro{\h}{#4}\pgfmathsetmacro{\p}{#5}%
    \pgfmathsetmacro{\inset}{(\h-\p)/2}%
    \fill[#6] (\cx,\cy) -- (\cx+\w/2,\cy+\inset) -- (\cx+\w,\cy) -- (\cx+\w,\cy+\h)
      -- (\cx+\w/2,\cy+\h-\inset) -- (\cx,\cy+\h) -- cycle;
    \draw[black,thin] (\cx,\cy) -- (\cx+\w/2,\cy+\inset) -- (\cx+\w,\cy) -- (\cx+\w,\cy+\h)
      -- (\cx+\w/2,\cy+\h-\inset) -- (\cx,\cy+\h) -- cycle;
    \node[align=center] at (\cx+\w/2,\cy+\h/2) {#7};
  \end{scope}
}

% ---------------------------------------------------------------
% \archtokengrid{x}{y}{rows}{cols}{cell}{color} -- grid of patch tokens
% ---------------------------------------------------------------
\newcommand{\archtokengrid}[6]{%
  \begin{scope}
    \pgfmathsetmacro{\cx}{#1}\pgfmathsetmacro{\cy}{#2}\pgfmathsetmacro{\cell}{#5}%
    \pgfmathtruncatemacro{\rmax}{#3-1}%
    \pgfmathtruncatemacro{\cmax}{#4-1}%
    \foreach \r in {0,...,\rmax}{%
      \foreach \c in {0,...,\cmax}{%
        \fill[#6] ({\cx+\c*\cell},{\cy+\r*\cell}) rectangle ++(\cell*0.82,\cell*0.82);
        \draw[black,thin] ({\cx+\c*\cell},{\cy+\r*\cell}) rectangle ++(\cell*0.82,\cell*0.82);
      }%
    }%
  \end{scope}
}

% ---------------------------------------------------------------
% \archarrow{start}{end}{color} -- block arrow between stages
% ---------------------------------------------------------------
% #1 and #2 are full coordinates including parentheses, e.g. (1,2)
\newcommand{\archarrow}[3]{%
  \draw[-{Latex[length=2.4mm,width=2mm]},line width=2.2pt,#3] #1 -- #2;
}

% ---------------------------------------------------------------
% Dataflow-arrow vocabulary -- opt-in via \archarrowset{compact|detailed},
% called once near the top of a figure (after \input{arch-macros.tex}).
% \archarrow keeps its classic small size by default for any arch_*.tex
% figure that never calls \archarrowset -- this block is purely additive.
%   archflowstyle (tikz key, NOT a macro -- use bare, no backslash) --
%                                  current draw options (head + line
%                                  width), for a hand-rolled elbow \draw
%                                  that \archarrow's straight "--" can't
%                                  express (e.g. "-|" routing). Must be a
%                                  /.style key rather than a \def'd macro:
%                                  pgf's arrow-tip parser doesn't split on
%                                  a comma embedded in a plain macro's
%                                  expansion when used as a \draw option.
%   \archflowNeutral/Target/Context -- the only 3 colors a plain dataflow
%                                  arrow should use: generic same-role
%                                  passthrough; target's own real output
%                                  leaving the decoder; a real label/mask
%                                  tensor entering fusion or attention
%                                  (context's ground truth, or target's
%                                  own real prior fed back from a
%                                  previous cascade level -- same role
%                                  either way: real mask data, not a
%                                  feature)
%   \archflow{p1}{p2}{role}     -- \archarrow wrapper keyed by role name
%                                  (role in {neutral, target, context})
% ---------------------------------------------------------------
\newcommand{\archarrowset}[1]{%
  \ifthenelse{\equal{#1}{detailed}}{%
    \tikzset{archflowstyle/.style={-{Latex[length=3.6mm,width=3.2mm]},line width=3.2pt}}%
  }{%
    \tikzset{archflowstyle/.style={-{Latex[length=2.4mm,width=2mm]},line width=2.2pt}}%
  }%
  \renewcommand{\archarrow}[3]{\draw[archflowstyle,##3] ##1 -- ##2;}%
}
\newcommand{\archflowNeutral}{black!30}
\newcommand{\archflowTarget}{archTargetAccent}
\newcommand{\archflowContext}{archContextAccent}
\newcommand{\archflow}[3]{%
  \ifthenelse{\equal{#3}{target}}{\archarrow{#1}{#2}{\archflowTarget}}{%
  \ifthenelse{\equal{#3}{context}}{\archarrow{#1}{#2}{\archflowContext}}{%
  \archarrow{#1}{#2}{\archflowNeutral}}}%
}
 after \usepackage{tikz,ifthen} and
% \usetikzlibrary{arrows.meta,calc,positioning}. All macro names are
% prefixed `arch` to avoid clashing with pgf/tikz internals.

% ---- palette (approximating the PowerPoint drafts) ----
\definecolor{archTargetGreen}{HTML}{9BC7AC}
\definecolor{archContextPink}{HTML}{E8B8B0}
\definecolor{archAttnBlue}{HTML}{BBD6EC}
\definecolor{archEncGray}{HTML}{E4E4E4}
\definecolor{archArrowPale}{HTML}{CFE6D6}
\definecolor{archCtBg}{HTML}{2B2B2B}
\definecolor{archCtBlob}{HTML}{8C8C8C}
\definecolor{archCtxYellow}{HTML}{E8D24C}
% darker, more saturated siblings of archTargetGreen/archContextPink --
% pastel fills read poorly as 1pt strokes, and mask/label content needs to
% read as "which volume owns this" at a glance, so masks/labels and any
% line (arrow, brace) touching that volume's row use these instead of the
% pastel role color or the old flat archMaskGold.
\definecolor{archTargetAccent}{HTML}{2F6B4A}  % target: darker green
\definecolor{archContextAccent}{HTML}{A8453E} % context: darker red

% fake-isometric depth skew shared by every cube
\newcommand{\archdx}{0.30}
\newcommand{\archdy}{0.20}

% ---- shared text-role macros ----
% Every arch_*.tex figure is compiled standalone at the SAME base size
% (documentclass[tikz,border=4pt,12pt]{standalone}, matching the thesis
% body's 12pt), so \normalsize etc. mean the same absolute point size in
% every figure. Using these three named roles instead of ad hoc \Large /
% \bfseries\normalsize / \scriptsize per figure keeps "a box title" (etc.)
% the same size everywhere it appears, regardless of which figure it's in
% or how that figure's canvas happens to be scaled by \includegraphics.
%   \archtitle    -- figure/method/panel titles, the "xL" repeat marker
%   \archboxlabel -- sub-block header labels (Fuse, Encoder, attention, ...)
%   \archnote     -- small secondary annotations (placeholder captions)
% Plain, unlabeled node text (cube/token labels, row/lane labels) stays at
% the bare \normalsize default -- no macro needed for that role.
%
% DECLARATIONS, not argument-taking macros -- deliberately, so they can
% precede a raw `\\` inside a tikz node's text (as `{\archboxlabel line
% one\\line two}`) without hiding that `\\` one macro-expansion level deep.
% A `\\` buried inside a wrapped {\archboxlabel{...}} argument breaks
% TikZ's own align=center line-splitting (it scans for `\\` before the
% argument expands) -- use these the same way `\bfseries`/`\Large` are
% already used directly inline elsewhere in these figures.
\newcommand{\archtitle}{\Large\bfseries}
\newcommand{\archboxlabel}{\bfseries\normalsize}
\newcommand{\archnote}{\scriptsize}

% ---------------------------------------------------------------
% \archvolcube{x}{y}{w}{h}{topcolor}{content}{label}
%   content: ct | ctmask | predmask | <flat color name>
% ---------------------------------------------------------------
\newcommand{\archvolcube}[7]{%
  \begin{scope}
    \pgfmathsetmacro{\cx}{#1}\pgfmathsetmacro{\cy}{#2}%
    \pgfmathsetmacro{\w}{#3}\pgfmathsetmacro{\h}{#4}%
    \fill[#5] (\cx,\cy+\h) -- ++(\archdx,\archdy) -- ++(\w,0) -- ++(-\archdx,-\archdy) -- cycle;
    \fill[#5!60!black] (\cx+\w,\cy) -- ++(\archdx,\archdy) -- ++(0,\h) -- ++(-\archdx,-\archdy) -- cycle;
    \begin{scope}
      \clip (\cx,\cy) rectangle ++(\w,\h);
      \ifthenelse{\equal{#6}{ct}}{%
        \fill[archCtBg] (\cx,\cy) rectangle ++(\w,\h);
        \fill[archCtBlob] plot[smooth cycle] coordinates
          {(\cx+.25*\w,\cy+.55*\h) (\cx+.45*\w,\cy+.75*\h) (\cx+.7*\w,\cy+.6*\h)
           (\cx+.6*\w,\cy+.3*\h) (\cx+.35*\w,\cy+.25*\h)};
      }{\ifthenelse{\equal{#6}{ctmask}}{%
        \fill[archCtBg] (\cx,\cy) rectangle ++(\w,\h);
        \fill[archCtBlob] plot[smooth cycle] coordinates
          {(\cx+.25*\w,\cy+.55*\h) (\cx+.45*\w,\cy+.75*\h) (\cx+.7*\w,\cy+.6*\h)
           (\cx+.6*\w,\cy+.3*\h) (\cx+.35*\w,\cy+.25*\h)};
        \fill[archCtxYellow] plot[smooth cycle] coordinates
          {(\cx+.42*\w,\cy+.42*\h) (\cx+.52*\w,\cy+.5*\h) (\cx+.5*\w,\cy+.35*\h) (\cx+.44*\w,\cy+.3*\h)};
      }{\ifthenelse{\equal{#6}{predmask}}{%
        \fill[black] (\cx,\cy) rectangle ++(\w,\h);
        \fill[white] plot[smooth cycle] coordinates
          {(\cx+.4*\w,\cy+.65*\h) (\cx+.6*\w,\cy+.7*\h) (\cx+.62*\w,\cy+.45*\h)
           (\cx+.48*\w,\cy+.3*\h) (\cx+.38*\w,\cy+.4*\h)};
      }{%
        \fill[#6] (\cx,\cy) rectangle ++(\w,\h);
      }}}%
    \end{scope}
    \draw[black,thin] (\cx,\cy) rectangle ++(\w,\h);
    \draw[black,thin] (\cx,\cy+\h) -- ++(\archdx,\archdy) -- ++(\w,0) -- ++(-\archdx,-\archdy) -- cycle;
    \draw[black,thin] (\cx+\w,\cy) -- ++(\archdx,\archdy) -- ++(0,\h) -- ++(-\archdx,-\archdy) -- cycle;
    \node[anchor=north,align=center] at (\cx+\w/2,\cy-0.1) {#7};
  \end{scope}
}

% ---------------------------------------------------------------
% \archvolcubeimg{x}{y}{w}{h}{topcolor}{imgpath}{label} -- same cube as
% \archvolcube, but the front face shows a real raster image (clipped to
% exactly fill the face, aspect ratio not preserved) instead of one of
% \archvolcube's procedural ct/ctmask/predmask blobs. Needs graphicx
% loaded by the including figure. A separate macro rather than a new
% \archvolcube content mode, so existing ct/ctmask/predmask figures are
% untouched.
% ---------------------------------------------------------------
\newcommand{\archvolcubeimg}[7]{%
  \begin{scope}
    \pgfmathsetmacro{\cx}{#1}\pgfmathsetmacro{\cy}{#2}%
    \pgfmathsetmacro{\w}{#3}\pgfmathsetmacro{\h}{#4}%
    \fill[#5] (\cx,\cy+\h) -- ++(\archdx,\archdy) -- ++(\w,0) -- ++(-\archdx,-\archdy) -- cycle;
    \fill[#5!60!black] (\cx+\w,\cy) -- ++(\archdx,\archdy) -- ++(0,\h) -- ++(-\archdx,-\archdy) -- cycle;
    \begin{scope}
      \clip (\cx,\cy) rectangle ++(\w,\h);
      \node[anchor=south west, inner sep=0pt] at (\cx,\cy)
        {\includegraphics[width=\w cm,height=\h cm]{#6}};
    \end{scope}
    \draw[black,thin] (\cx,\cy) rectangle ++(\w,\h);
    \draw[black,thin] (\cx,\cy+\h) -- ++(\archdx,\archdy) -- ++(\w,0) -- ++(-\archdx,-\archdy) -- cycle;
    \draw[black,thin] (\cx+\w,\cy) -- ++(\archdx,\archdy) -- ++(0,\h) -- ++(-\archdx,-\archdy) -- cycle;
    \node[anchor=north,align=center] at (\cx+\w/2,\cy-0.1) {#7};
  \end{scope}
}

% ---------------------------------------------------------------
% \archfeatcube{x}{y}{w}{h}{color}{label} -- flat-colored feature volume
% ---------------------------------------------------------------
\newcommand{\archfeatcube}[6]{%
  \begin{scope}
    \pgfmathsetmacro{\cx}{#1}\pgfmathsetmacro{\cy}{#2}%
    \pgfmathsetmacro{\w}{#3}\pgfmathsetmacro{\h}{#4}%
    \fill[#5] (\cx,\cy) rectangle ++(\w,\h);
    \fill[#5!80!white] (\cx,\cy+\h) -- ++(\archdx,\archdy) -- ++(\w,0) -- ++(-\archdx,-\archdy) -- cycle;
    \fill[#5!60!black] (\cx+\w,\cy) -- ++(\archdx,\archdy) -- ++(0,\h) -- ++(-\archdx,-\archdy) -- cycle;
    \draw[black,thin] (\cx,\cy) rectangle ++(\w,\h);
    \draw[black,thin] (\cx,\cy+\h) -- ++(\archdx,\archdy) -- ++(\w,0) -- ++(-\archdx,-\archdy) -- cycle;
    \draw[black,thin] (\cx+\w,\cy) -- ++(\archdx,\archdy) -- ++(0,\h) -- ++(-\archdx,-\archdy) -- cycle;
    \node[align=center] at (\cx+\w/2,\cy+\h/2) {#6};
  \end{scope}
}

% ---------------------------------------------------------------
% \archfunnel{x}{y}{h}{wwide}{wnarrow}{dir}{color}{label}
%   dir=1  : wide on left, narrow on right  (encoder)
%   dir=-1 : narrow on left, wide on right  (decoder)
% ---------------------------------------------------------------
\newcommand{\archfunnel}[8]{%
  \begin{scope}
    \pgfmathsetmacro{\cx}{#1}\pgfmathsetmacro{\cy}{#2}\pgfmathsetmacro{\h}{#3}%
    \pgfmathsetmacro{\ww}{#4}\pgfmathsetmacro{\wn}{#5}%
    \pgfmathsetmacro{\inset}{(\ww-\wn)/2}%
    \ifnum#6=1
      \fill[#7] (\cx,\cy) -- (\cx+\ww,\cy+\inset) -- (\cx+\ww,\cy+\h-\inset) -- (\cx,\cy+\h) -- cycle;
      \draw[black,thin] (\cx,\cy) -- (\cx+\ww,\cy+\inset) -- (\cx+\ww,\cy+\h-\inset) -- (\cx,\cy+\h) -- cycle;
    \else
      \fill[#7] (\cx,\cy+\inset) -- (\cx+\ww,\cy) -- (\cx+\ww,\cy+\h) -- (\cx,\cy+\h-\inset) -- cycle;
      \draw[black,thin] (\cx,\cy+\inset) -- (\cx+\ww,\cy) -- (\cx+\ww,\cy+\h) -- (\cx,\cy+\h-\inset) -- cycle;
    \fi
    \node[align=center] at (\cx+\ww/2,\cy+\h/2) {#8};
  \end{scope}
}

% ---------------------------------------------------------------
% \archbowtie{x}{y}{w}{h}{pinch}{color}{label} -- Medverse-style
% single-row "dual U-Net" hourglass (funnel-in then funnel-out).
% ---------------------------------------------------------------
\newcommand{\archbowtie}[7]{%
  \begin{scope}
    \pgfmathsetmacro{\cx}{#1}\pgfmathsetmacro{\cy}{#2}%
    \pgfmathsetmacro{\w}{#3}\pgfmathsetmacro{\h}{#4}\pgfmathsetmacro{\p}{#5}%
    \pgfmathsetmacro{\inset}{(\h-\p)/2}%
    \fill[#6] (\cx,\cy) -- (\cx+\w/2,\cy+\inset) -- (\cx+\w,\cy) -- (\cx+\w,\cy+\h)
      -- (\cx+\w/2,\cy+\h-\inset) -- (\cx,\cy+\h) -- cycle;
    \draw[black,thin] (\cx,\cy) -- (\cx+\w/2,\cy+\inset) -- (\cx+\w,\cy) -- (\cx+\w,\cy+\h)
      -- (\cx+\w/2,\cy+\h-\inset) -- (\cx,\cy+\h) -- cycle;
    \node[align=center] at (\cx+\w/2,\cy+\h/2) {#7};
  \end{scope}
}

% ---------------------------------------------------------------
% \archtokengrid{x}{y}{rows}{cols}{cell}{color} -- grid of patch tokens
% ---------------------------------------------------------------
\newcommand{\archtokengrid}[6]{%
  \begin{scope}
    \pgfmathsetmacro{\cx}{#1}\pgfmathsetmacro{\cy}{#2}\pgfmathsetmacro{\cell}{#5}%
    \pgfmathtruncatemacro{\rmax}{#3-1}%
    \pgfmathtruncatemacro{\cmax}{#4-1}%
    \foreach \r in {0,...,\rmax}{%
      \foreach \c in {0,...,\cmax}{%
        \fill[#6] ({\cx+\c*\cell},{\cy+\r*\cell}) rectangle ++(\cell*0.82,\cell*0.82);
        \draw[black,thin] ({\cx+\c*\cell},{\cy+\r*\cell}) rectangle ++(\cell*0.82,\cell*0.82);
      }%
    }%
  \end{scope}
}

% ---------------------------------------------------------------
% \archarrow{start}{end}{color} -- block arrow between stages
% ---------------------------------------------------------------
% #1 and #2 are full coordinates including parentheses, e.g. (1,2)
\newcommand{\archarrow}[3]{%
  \draw[-{Latex[length=2.4mm,width=2mm]},line width=2.2pt,#3] #1 -- #2;
}

% ---------------------------------------------------------------
% Dataflow-arrow vocabulary -- opt-in via \archarrowset{compact|detailed},
% called once near the top of a figure (after % Shared TikZ macros for the architecture figures (Medverse vs. Ours).
% Load with \input{arch-macros.tex} after \usepackage{tikz,ifthen} and
% \usetikzlibrary{arrows.meta,calc,positioning}. All macro names are
% prefixed `arch` to avoid clashing with pgf/tikz internals.

% ---- palette (approximating the PowerPoint drafts) ----
\definecolor{archTargetGreen}{HTML}{9BC7AC}
\definecolor{archContextPink}{HTML}{E8B8B0}
\definecolor{archAttnBlue}{HTML}{BBD6EC}
\definecolor{archEncGray}{HTML}{E4E4E4}
\definecolor{archArrowPale}{HTML}{CFE6D6}
\definecolor{archCtBg}{HTML}{2B2B2B}
\definecolor{archCtBlob}{HTML}{8C8C8C}
\definecolor{archCtxYellow}{HTML}{E8D24C}
% darker, more saturated siblings of archTargetGreen/archContextPink --
% pastel fills read poorly as 1pt strokes, and mask/label content needs to
% read as "which volume owns this" at a glance, so masks/labels and any
% line (arrow, brace) touching that volume's row use these instead of the
% pastel role color or the old flat archMaskGold.
\definecolor{archTargetAccent}{HTML}{2F6B4A}  % target: darker green
\definecolor{archContextAccent}{HTML}{A8453E} % context: darker red

% fake-isometric depth skew shared by every cube
\newcommand{\archdx}{0.30}
\newcommand{\archdy}{0.20}

% ---- shared text-role macros ----
% Every arch_*.tex figure is compiled standalone at the SAME base size
% (documentclass[tikz,border=4pt,12pt]{standalone}, matching the thesis
% body's 12pt), so \normalsize etc. mean the same absolute point size in
% every figure. Using these three named roles instead of ad hoc \Large /
% \bfseries\normalsize / \scriptsize per figure keeps "a box title" (etc.)
% the same size everywhere it appears, regardless of which figure it's in
% or how that figure's canvas happens to be scaled by \includegraphics.
%   \archtitle    -- figure/method/panel titles, the "xL" repeat marker
%   \archboxlabel -- sub-block header labels (Fuse, Encoder, attention, ...)
%   \archnote     -- small secondary annotations (placeholder captions)
% Plain, unlabeled node text (cube/token labels, row/lane labels) stays at
% the bare \normalsize default -- no macro needed for that role.
%
% DECLARATIONS, not argument-taking macros -- deliberately, so they can
% precede a raw `\\` inside a tikz node's text (as `{\archboxlabel line
% one\\line two}`) without hiding that `\\` one macro-expansion level deep.
% A `\\` buried inside a wrapped {\archboxlabel{...}} argument breaks
% TikZ's own align=center line-splitting (it scans for `\\` before the
% argument expands) -- use these the same way `\bfseries`/`\Large` are
% already used directly inline elsewhere in these figures.
\newcommand{\archtitle}{\Large\bfseries}
\newcommand{\archboxlabel}{\bfseries\normalsize}
\newcommand{\archnote}{\scriptsize}

% ---------------------------------------------------------------
% \archvolcube{x}{y}{w}{h}{topcolor}{content}{label}
%   content: ct | ctmask | predmask | <flat color name>
% ---------------------------------------------------------------
\newcommand{\archvolcube}[7]{%
  \begin{scope}
    \pgfmathsetmacro{\cx}{#1}\pgfmathsetmacro{\cy}{#2}%
    \pgfmathsetmacro{\w}{#3}\pgfmathsetmacro{\h}{#4}%
    \fill[#5] (\cx,\cy+\h) -- ++(\archdx,\archdy) -- ++(\w,0) -- ++(-\archdx,-\archdy) -- cycle;
    \fill[#5!60!black] (\cx+\w,\cy) -- ++(\archdx,\archdy) -- ++(0,\h) -- ++(-\archdx,-\archdy) -- cycle;
    \begin{scope}
      \clip (\cx,\cy) rectangle ++(\w,\h);
      \ifthenelse{\equal{#6}{ct}}{%
        \fill[archCtBg] (\cx,\cy) rectangle ++(\w,\h);
        \fill[archCtBlob] plot[smooth cycle] coordinates
          {(\cx+.25*\w,\cy+.55*\h) (\cx+.45*\w,\cy+.75*\h) (\cx+.7*\w,\cy+.6*\h)
           (\cx+.6*\w,\cy+.3*\h) (\cx+.35*\w,\cy+.25*\h)};
      }{\ifthenelse{\equal{#6}{ctmask}}{%
        \fill[archCtBg] (\cx,\cy) rectangle ++(\w,\h);
        \fill[archCtBlob] plot[smooth cycle] coordinates
          {(\cx+.25*\w,\cy+.55*\h) (\cx+.45*\w,\cy+.75*\h) (\cx+.7*\w,\cy+.6*\h)
           (\cx+.6*\w,\cy+.3*\h) (\cx+.35*\w,\cy+.25*\h)};
        \fill[archCtxYellow] plot[smooth cycle] coordinates
          {(\cx+.42*\w,\cy+.42*\h) (\cx+.52*\w,\cy+.5*\h) (\cx+.5*\w,\cy+.35*\h) (\cx+.44*\w,\cy+.3*\h)};
      }{\ifthenelse{\equal{#6}{predmask}}{%
        \fill[black] (\cx,\cy) rectangle ++(\w,\h);
        \fill[white] plot[smooth cycle] coordinates
          {(\cx+.4*\w,\cy+.65*\h) (\cx+.6*\w,\cy+.7*\h) (\cx+.62*\w,\cy+.45*\h)
           (\cx+.48*\w,\cy+.3*\h) (\cx+.38*\w,\cy+.4*\h)};
      }{%
        \fill[#6] (\cx,\cy) rectangle ++(\w,\h);
      }}}%
    \end{scope}
    \draw[black,thin] (\cx,\cy) rectangle ++(\w,\h);
    \draw[black,thin] (\cx,\cy+\h) -- ++(\archdx,\archdy) -- ++(\w,0) -- ++(-\archdx,-\archdy) -- cycle;
    \draw[black,thin] (\cx+\w,\cy) -- ++(\archdx,\archdy) -- ++(0,\h) -- ++(-\archdx,-\archdy) -- cycle;
    \node[anchor=north,align=center] at (\cx+\w/2,\cy-0.1) {#7};
  \end{scope}
}

% ---------------------------------------------------------------
% \archvolcubeimg{x}{y}{w}{h}{topcolor}{imgpath}{label} -- same cube as
% \archvolcube, but the front face shows a real raster image (clipped to
% exactly fill the face, aspect ratio not preserved) instead of one of
% \archvolcube's procedural ct/ctmask/predmask blobs. Needs graphicx
% loaded by the including figure. A separate macro rather than a new
% \archvolcube content mode, so existing ct/ctmask/predmask figures are
% untouched.
% ---------------------------------------------------------------
\newcommand{\archvolcubeimg}[7]{%
  \begin{scope}
    \pgfmathsetmacro{\cx}{#1}\pgfmathsetmacro{\cy}{#2}%
    \pgfmathsetmacro{\w}{#3}\pgfmathsetmacro{\h}{#4}%
    \fill[#5] (\cx,\cy+\h) -- ++(\archdx,\archdy) -- ++(\w,0) -- ++(-\archdx,-\archdy) -- cycle;
    \fill[#5!60!black] (\cx+\w,\cy) -- ++(\archdx,\archdy) -- ++(0,\h) -- ++(-\archdx,-\archdy) -- cycle;
    \begin{scope}
      \clip (\cx,\cy) rectangle ++(\w,\h);
      \node[anchor=south west, inner sep=0pt] at (\cx,\cy)
        {\includegraphics[width=\w cm,height=\h cm]{#6}};
    \end{scope}
    \draw[black,thin] (\cx,\cy) rectangle ++(\w,\h);
    \draw[black,thin] (\cx,\cy+\h) -- ++(\archdx,\archdy) -- ++(\w,0) -- ++(-\archdx,-\archdy) -- cycle;
    \draw[black,thin] (\cx+\w,\cy) -- ++(\archdx,\archdy) -- ++(0,\h) -- ++(-\archdx,-\archdy) -- cycle;
    \node[anchor=north,align=center] at (\cx+\w/2,\cy-0.1) {#7};
  \end{scope}
}

% ---------------------------------------------------------------
% \archfeatcube{x}{y}{w}{h}{color}{label} -- flat-colored feature volume
% ---------------------------------------------------------------
\newcommand{\archfeatcube}[6]{%
  \begin{scope}
    \pgfmathsetmacro{\cx}{#1}\pgfmathsetmacro{\cy}{#2}%
    \pgfmathsetmacro{\w}{#3}\pgfmathsetmacro{\h}{#4}%
    \fill[#5] (\cx,\cy) rectangle ++(\w,\h);
    \fill[#5!80!white] (\cx,\cy+\h) -- ++(\archdx,\archdy) -- ++(\w,0) -- ++(-\archdx,-\archdy) -- cycle;
    \fill[#5!60!black] (\cx+\w,\cy) -- ++(\archdx,\archdy) -- ++(0,\h) -- ++(-\archdx,-\archdy) -- cycle;
    \draw[black,thin] (\cx,\cy) rectangle ++(\w,\h);
    \draw[black,thin] (\cx,\cy+\h) -- ++(\archdx,\archdy) -- ++(\w,0) -- ++(-\archdx,-\archdy) -- cycle;
    \draw[black,thin] (\cx+\w,\cy) -- ++(\archdx,\archdy) -- ++(0,\h) -- ++(-\archdx,-\archdy) -- cycle;
    \node[align=center] at (\cx+\w/2,\cy+\h/2) {#6};
  \end{scope}
}

% ---------------------------------------------------------------
% \archfunnel{x}{y}{h}{wwide}{wnarrow}{dir}{color}{label}
%   dir=1  : wide on left, narrow on right  (encoder)
%   dir=-1 : narrow on left, wide on right  (decoder)
% ---------------------------------------------------------------
\newcommand{\archfunnel}[8]{%
  \begin{scope}
    \pgfmathsetmacro{\cx}{#1}\pgfmathsetmacro{\cy}{#2}\pgfmathsetmacro{\h}{#3}%
    \pgfmathsetmacro{\ww}{#4}\pgfmathsetmacro{\wn}{#5}%
    \pgfmathsetmacro{\inset}{(\ww-\wn)/2}%
    \ifnum#6=1
      \fill[#7] (\cx,\cy) -- (\cx+\ww,\cy+\inset) -- (\cx+\ww,\cy+\h-\inset) -- (\cx,\cy+\h) -- cycle;
      \draw[black,thin] (\cx,\cy) -- (\cx+\ww,\cy+\inset) -- (\cx+\ww,\cy+\h-\inset) -- (\cx,\cy+\h) -- cycle;
    \else
      \fill[#7] (\cx,\cy+\inset) -- (\cx+\ww,\cy) -- (\cx+\ww,\cy+\h) -- (\cx,\cy+\h-\inset) -- cycle;
      \draw[black,thin] (\cx,\cy+\inset) -- (\cx+\ww,\cy) -- (\cx+\ww,\cy+\h) -- (\cx,\cy+\h-\inset) -- cycle;
    \fi
    \node[align=center] at (\cx+\ww/2,\cy+\h/2) {#8};
  \end{scope}
}

% ---------------------------------------------------------------
% \archbowtie{x}{y}{w}{h}{pinch}{color}{label} -- Medverse-style
% single-row "dual U-Net" hourglass (funnel-in then funnel-out).
% ---------------------------------------------------------------
\newcommand{\archbowtie}[7]{%
  \begin{scope}
    \pgfmathsetmacro{\cx}{#1}\pgfmathsetmacro{\cy}{#2}%
    \pgfmathsetmacro{\w}{#3}\pgfmathsetmacro{\h}{#4}\pgfmathsetmacro{\p}{#5}%
    \pgfmathsetmacro{\inset}{(\h-\p)/2}%
    \fill[#6] (\cx,\cy) -- (\cx+\w/2,\cy+\inset) -- (\cx+\w,\cy) -- (\cx+\w,\cy+\h)
      -- (\cx+\w/2,\cy+\h-\inset) -- (\cx,\cy+\h) -- cycle;
    \draw[black,thin] (\cx,\cy) -- (\cx+\w/2,\cy+\inset) -- (\cx+\w,\cy) -- (\cx+\w,\cy+\h)
      -- (\cx+\w/2,\cy+\h-\inset) -- (\cx,\cy+\h) -- cycle;
    \node[align=center] at (\cx+\w/2,\cy+\h/2) {#7};
  \end{scope}
}

% ---------------------------------------------------------------
% \archtokengrid{x}{y}{rows}{cols}{cell}{color} -- grid of patch tokens
% ---------------------------------------------------------------
\newcommand{\archtokengrid}[6]{%
  \begin{scope}
    \pgfmathsetmacro{\cx}{#1}\pgfmathsetmacro{\cy}{#2}\pgfmathsetmacro{\cell}{#5}%
    \pgfmathtruncatemacro{\rmax}{#3-1}%
    \pgfmathtruncatemacro{\cmax}{#4-1}%
    \foreach \r in {0,...,\rmax}{%
      \foreach \c in {0,...,\cmax}{%
        \fill[#6] ({\cx+\c*\cell},{\cy+\r*\cell}) rectangle ++(\cell*0.82,\cell*0.82);
        \draw[black,thin] ({\cx+\c*\cell},{\cy+\r*\cell}) rectangle ++(\cell*0.82,\cell*0.82);
      }%
    }%
  \end{scope}
}

% ---------------------------------------------------------------
% \archarrow{start}{end}{color} -- block arrow between stages
% ---------------------------------------------------------------
% #1 and #2 are full coordinates including parentheses, e.g. (1,2)
\newcommand{\archarrow}[3]{%
  \draw[-{Latex[length=2.4mm,width=2mm]},line width=2.2pt,#3] #1 -- #2;
}

% ---------------------------------------------------------------
% Dataflow-arrow vocabulary -- opt-in via \archarrowset{compact|detailed},
% called once near the top of a figure (after \input{arch-macros.tex}).
% \archarrow keeps its classic small size by default for any arch_*.tex
% figure that never calls \archarrowset -- this block is purely additive.
%   archflowstyle (tikz key, NOT a macro -- use bare, no backslash) --
%                                  current draw options (head + line
%                                  width), for a hand-rolled elbow \draw
%                                  that \archarrow's straight "--" can't
%                                  express (e.g. "-|" routing). Must be a
%                                  /.style key rather than a \def'd macro:
%                                  pgf's arrow-tip parser doesn't split on
%                                  a comma embedded in a plain macro's
%                                  expansion when used as a \draw option.
%   \archflowNeutral/Target/Context -- the only 3 colors a plain dataflow
%                                  arrow should use: generic same-role
%                                  passthrough; target's own real output
%                                  leaving the decoder; a real label/mask
%                                  tensor entering fusion or attention
%                                  (context's ground truth, or target's
%                                  own real prior fed back from a
%                                  previous cascade level -- same role
%                                  either way: real mask data, not a
%                                  feature)
%   \archflow{p1}{p2}{role}     -- \archarrow wrapper keyed by role name
%                                  (role in {neutral, target, context})
% ---------------------------------------------------------------
\newcommand{\archarrowset}[1]{%
  \ifthenelse{\equal{#1}{detailed}}{%
    \tikzset{archflowstyle/.style={-{Latex[length=3.6mm,width=3.2mm]},line width=3.2pt}}%
  }{%
    \tikzset{archflowstyle/.style={-{Latex[length=2.4mm,width=2mm]},line width=2.2pt}}%
  }%
  \renewcommand{\archarrow}[3]{\draw[archflowstyle,##3] ##1 -- ##2;}%
}
\newcommand{\archflowNeutral}{black!30}
\newcommand{\archflowTarget}{archTargetAccent}
\newcommand{\archflowContext}{archContextAccent}
\newcommand{\archflow}[3]{%
  \ifthenelse{\equal{#3}{target}}{\archarrow{#1}{#2}{\archflowTarget}}{%
  \ifthenelse{\equal{#3}{context}}{\archarrow{#1}{#2}{\archflowContext}}{%
  \archarrow{#1}{#2}{\archflowNeutral}}}%
}
).
% \archarrow keeps its classic small size by default for any arch_*.tex
% figure that never calls \archarrowset -- this block is purely additive.
%   archflowstyle (tikz key, NOT a macro -- use bare, no backslash) --
%                                  current draw options (head + line
%                                  width), for a hand-rolled elbow \draw
%                                  that \archarrow's straight "--" can't
%                                  express (e.g. "-|" routing). Must be a
%                                  /.style key rather than a \def'd macro:
%                                  pgf's arrow-tip parser doesn't split on
%                                  a comma embedded in a plain macro's
%                                  expansion when used as a \draw option.
%   \archflowNeutral/Target/Context -- the only 3 colors a plain dataflow
%                                  arrow should use: generic same-role
%                                  passthrough; target's own real output
%                                  leaving the decoder; a real label/mask
%                                  tensor entering fusion or attention
%                                  (context's ground truth, or target's
%                                  own real prior fed back from a
%                                  previous cascade level -- same role
%                                  either way: real mask data, not a
%                                  feature)
%   \archflow{p1}{p2}{role}     -- \archarrow wrapper keyed by role name
%                                  (role in {neutral, target, context})
% ---------------------------------------------------------------
\newcommand{\archarrowset}[1]{%
  \ifthenelse{\equal{#1}{detailed}}{%
    \tikzset{archflowstyle/.style={-{Latex[length=3.6mm,width=3.2mm]},line width=3.2pt}}%
  }{%
    \tikzset{archflowstyle/.style={-{Latex[length=2.4mm,width=2mm]},line width=2.2pt}}%
  }%
  \renewcommand{\archarrow}[3]{\draw[archflowstyle,##3] ##1 -- ##2;}%
}
\newcommand{\archflowNeutral}{black!30}
\newcommand{\archflowTarget}{archTargetAccent}
\newcommand{\archflowContext}{archContextAccent}
\newcommand{\archflow}[3]{%
  \ifthenelse{\equal{#3}{target}}{\archarrow{#1}{#2}{\archflowTarget}}{%
  \ifthenelse{\equal{#3}{context}}{\archarrow{#1}{#2}{\archflowContext}}{%
  \archarrow{#1}{#2}{\archflowNeutral}}}%
}
).
% \archarrow keeps its classic small size by default for any arch_*.tex
% figure that never calls \archarrowset -- this block is purely additive.
%   archflowstyle (tikz key, NOT a macro -- use bare, no backslash) --
%                                  current draw options (head + line
%                                  width), for a hand-rolled elbow \draw
%                                  that \archarrow's straight "--" can't
%                                  express (e.g. "-|" routing). Must be a
%                                  /.style key rather than a \def'd macro:
%                                  pgf's arrow-tip parser doesn't split on
%                                  a comma embedded in a plain macro's
%                                  expansion when used as a \draw option.
%   \archflowNeutral/Target/Context -- the only 3 colors a plain dataflow
%                                  arrow should use: generic same-role
%                                  passthrough; target's own real output
%                                  leaving the decoder; a real label/mask
%                                  tensor entering fusion or attention
%                                  (context's ground truth, or target's
%                                  own real prior fed back from a
%                                  previous cascade level -- same role
%                                  either way: real mask data, not a
%                                  feature)
%   \archflow{p1}{p2}{role}     -- \archarrow wrapper keyed by role name
%                                  (role in {neutral, target, context})
% ---------------------------------------------------------------
\newcommand{\archarrowset}[1]{%
  \ifthenelse{\equal{#1}{detailed}}{%
    \tikzset{archflowstyle/.style={-{Latex[length=3.6mm,width=3.2mm]},line width=3.2pt}}%
  }{%
    \tikzset{archflowstyle/.style={-{Latex[length=2.4mm,width=2mm]},line width=2.2pt}}%
  }%
  \renewcommand{\archarrow}[3]{\draw[archflowstyle,##3] ##1 -- ##2;}%
}
\newcommand{\archflowNeutral}{black!30}
\newcommand{\archflowTarget}{archTargetAccent}
\newcommand{\archflowContext}{archContextAccent}
\newcommand{\archflow}[3]{%
  \ifthenelse{\equal{#3}{target}}{\archarrow{#1}{#2}{\archflowTarget}}{%
  \ifthenelse{\equal{#3}{context}}{\archarrow{#1}{#2}{\archflowContext}}{%
  \archarrow{#1}{#2}{\archflowNeutral}}}%
}
 after \usepackage{tikz,ifthen} and
% \usetikzlibrary{arrows.meta,calc,positioning}. All macro names are
% prefixed `arch` to avoid clashing with pgf/tikz internals.

% ---- palette (approximating the PowerPoint drafts) ----
\definecolor{archTargetGreen}{HTML}{9BC7AC}
\definecolor{archContextPink}{HTML}{E8B8B0}
\definecolor{archAttnBlue}{HTML}{BBD6EC}
\definecolor{archEncGray}{HTML}{E4E4E4}
\definecolor{archArrowPale}{HTML}{CFE6D6}
\definecolor{archCtBg}{HTML}{2B2B2B}
\definecolor{archCtBlob}{HTML}{8C8C8C}
\definecolor{archCtxYellow}{HTML}{E8D24C}
% darker, more saturated siblings of archTargetGreen/archContextPink --
% pastel fills read poorly as 1pt strokes, and mask/label content needs to
% read as "which volume owns this" at a glance, so masks/labels and any
% line (arrow, brace) touching that volume's row use these instead of the
% pastel role color or the old flat archMaskGold.
\definecolor{archTargetAccent}{HTML}{2F6B4A}  % target: darker green
\definecolor{archContextAccent}{HTML}{A8453E} % context: darker red

% fake-isometric depth skew shared by every cube
\newcommand{\archdx}{0.30}
\newcommand{\archdy}{0.20}

% ---- shared text-role macros ----
% Every arch_*.tex figure is compiled standalone at the SAME base size
% (documentclass[tikz,border=4pt,12pt]{standalone}, matching the thesis
% body's 12pt), so \normalsize etc. mean the same absolute point size in
% every figure. Using these three named roles instead of ad hoc \Large /
% \bfseries\normalsize / \scriptsize per figure keeps "a box title" (etc.)
% the same size everywhere it appears, regardless of which figure it's in
% or how that figure's canvas happens to be scaled by \includegraphics.
%   \archtitle    -- figure/method/panel titles, the "xL" repeat marker
%   \archboxlabel -- sub-block header labels (Fuse, Encoder, attention, ...)
%   \archnote     -- small secondary annotations (placeholder captions)
% Plain, unlabeled node text (cube/token labels, row/lane labels) stays at
% the bare \normalsize default -- no macro needed for that role.
%
% DECLARATIONS, not argument-taking macros -- deliberately, so they can
% precede a raw `\\` inside a tikz node's text (as `{\archboxlabel line
% one\\line two}`) without hiding that `\\` one macro-expansion level deep.
% A `\\` buried inside a wrapped {\archboxlabel{...}} argument breaks
% TikZ's own align=center line-splitting (it scans for `\\` before the
% argument expands) -- use these the same way `\bfseries`/`\Large` are
% already used directly inline elsewhere in these figures.
\newcommand{\archtitle}{\Large\bfseries}
\newcommand{\archboxlabel}{\bfseries\normalsize}
\newcommand{\archnote}{\scriptsize}

% ---------------------------------------------------------------
% \archvolcube{x}{y}{w}{h}{topcolor}{content}{label}
%   content: ct | ctmask | predmask | <flat color name>
% ---------------------------------------------------------------
\newcommand{\archvolcube}[7]{%
  \begin{scope}
    \pgfmathsetmacro{\cx}{#1}\pgfmathsetmacro{\cy}{#2}%
    \pgfmathsetmacro{\w}{#3}\pgfmathsetmacro{\h}{#4}%
    \fill[#5] (\cx,\cy+\h) -- ++(\archdx,\archdy) -- ++(\w,0) -- ++(-\archdx,-\archdy) -- cycle;
    \fill[#5!60!black] (\cx+\w,\cy) -- ++(\archdx,\archdy) -- ++(0,\h) -- ++(-\archdx,-\archdy) -- cycle;
    \begin{scope}
      \clip (\cx,\cy) rectangle ++(\w,\h);
      \ifthenelse{\equal{#6}{ct}}{%
        \fill[archCtBg] (\cx,\cy) rectangle ++(\w,\h);
        \fill[archCtBlob] plot[smooth cycle] coordinates
          {(\cx+.25*\w,\cy+.55*\h) (\cx+.45*\w,\cy+.75*\h) (\cx+.7*\w,\cy+.6*\h)
           (\cx+.6*\w,\cy+.3*\h) (\cx+.35*\w,\cy+.25*\h)};
      }{\ifthenelse{\equal{#6}{ctmask}}{%
        \fill[archCtBg] (\cx,\cy) rectangle ++(\w,\h);
        \fill[archCtBlob] plot[smooth cycle] coordinates
          {(\cx+.25*\w,\cy+.55*\h) (\cx+.45*\w,\cy+.75*\h) (\cx+.7*\w,\cy+.6*\h)
           (\cx+.6*\w,\cy+.3*\h) (\cx+.35*\w,\cy+.25*\h)};
        \fill[archCtxYellow] plot[smooth cycle] coordinates
          {(\cx+.42*\w,\cy+.42*\h) (\cx+.52*\w,\cy+.5*\h) (\cx+.5*\w,\cy+.35*\h) (\cx+.44*\w,\cy+.3*\h)};
      }{\ifthenelse{\equal{#6}{predmask}}{%
        \fill[black] (\cx,\cy) rectangle ++(\w,\h);
        \fill[white] plot[smooth cycle] coordinates
          {(\cx+.4*\w,\cy+.65*\h) (\cx+.6*\w,\cy+.7*\h) (\cx+.62*\w,\cy+.45*\h)
           (\cx+.48*\w,\cy+.3*\h) (\cx+.38*\w,\cy+.4*\h)};
      }{%
        \fill[#6] (\cx,\cy) rectangle ++(\w,\h);
      }}}%
    \end{scope}
    \draw[black,thin] (\cx,\cy) rectangle ++(\w,\h);
    \draw[black,thin] (\cx,\cy+\h) -- ++(\archdx,\archdy) -- ++(\w,0) -- ++(-\archdx,-\archdy) -- cycle;
    \draw[black,thin] (\cx+\w,\cy) -- ++(\archdx,\archdy) -- ++(0,\h) -- ++(-\archdx,-\archdy) -- cycle;
    \node[anchor=north,align=center] at (\cx+\w/2,\cy-0.1) {#7};
  \end{scope}
}

% ---------------------------------------------------------------
% \archvolcubeimg{x}{y}{w}{h}{topcolor}{imgpath}{label} -- same cube as
% \archvolcube, but the front face shows a real raster image (clipped to
% exactly fill the face, aspect ratio not preserved) instead of one of
% \archvolcube's procedural ct/ctmask/predmask blobs. Needs graphicx
% loaded by the including figure. A separate macro rather than a new
% \archvolcube content mode, so existing ct/ctmask/predmask figures are
% untouched.
% ---------------------------------------------------------------
\newcommand{\archvolcubeimg}[7]{%
  \begin{scope}
    \pgfmathsetmacro{\cx}{#1}\pgfmathsetmacro{\cy}{#2}%
    \pgfmathsetmacro{\w}{#3}\pgfmathsetmacro{\h}{#4}%
    \fill[#5] (\cx,\cy+\h) -- ++(\archdx,\archdy) -- ++(\w,0) -- ++(-\archdx,-\archdy) -- cycle;
    \fill[#5!60!black] (\cx+\w,\cy) -- ++(\archdx,\archdy) -- ++(0,\h) -- ++(-\archdx,-\archdy) -- cycle;
    \begin{scope}
      \clip (\cx,\cy) rectangle ++(\w,\h);
      \node[anchor=south west, inner sep=0pt] at (\cx,\cy)
        {\includegraphics[width=\w cm,height=\h cm]{#6}};
    \end{scope}
    \draw[black,thin] (\cx,\cy) rectangle ++(\w,\h);
    \draw[black,thin] (\cx,\cy+\h) -- ++(\archdx,\archdy) -- ++(\w,0) -- ++(-\archdx,-\archdy) -- cycle;
    \draw[black,thin] (\cx+\w,\cy) -- ++(\archdx,\archdy) -- ++(0,\h) -- ++(-\archdx,-\archdy) -- cycle;
    \node[anchor=north,align=center] at (\cx+\w/2,\cy-0.1) {#7};
  \end{scope}
}

% ---------------------------------------------------------------
% \archfeatcube{x}{y}{w}{h}{color}{label} -- flat-colored feature volume
% ---------------------------------------------------------------
\newcommand{\archfeatcube}[6]{%
  \begin{scope}
    \pgfmathsetmacro{\cx}{#1}\pgfmathsetmacro{\cy}{#2}%
    \pgfmathsetmacro{\w}{#3}\pgfmathsetmacro{\h}{#4}%
    \fill[#5] (\cx,\cy) rectangle ++(\w,\h);
    \fill[#5!80!white] (\cx,\cy+\h) -- ++(\archdx,\archdy) -- ++(\w,0) -- ++(-\archdx,-\archdy) -- cycle;
    \fill[#5!60!black] (\cx+\w,\cy) -- ++(\archdx,\archdy) -- ++(0,\h) -- ++(-\archdx,-\archdy) -- cycle;
    \draw[black,thin] (\cx,\cy) rectangle ++(\w,\h);
    \draw[black,thin] (\cx,\cy+\h) -- ++(\archdx,\archdy) -- ++(\w,0) -- ++(-\archdx,-\archdy) -- cycle;
    \draw[black,thin] (\cx+\w,\cy) -- ++(\archdx,\archdy) -- ++(0,\h) -- ++(-\archdx,-\archdy) -- cycle;
    \node[align=center] at (\cx+\w/2,\cy+\h/2) {#6};
  \end{scope}
}

% ---------------------------------------------------------------
% \archfunnel{x}{y}{h}{wwide}{wnarrow}{dir}{color}{label}
%   dir=1  : wide on left, narrow on right  (encoder)
%   dir=-1 : narrow on left, wide on right  (decoder)
% ---------------------------------------------------------------
\newcommand{\archfunnel}[8]{%
  \begin{scope}
    \pgfmathsetmacro{\cx}{#1}\pgfmathsetmacro{\cy}{#2}\pgfmathsetmacro{\h}{#3}%
    \pgfmathsetmacro{\ww}{#4}\pgfmathsetmacro{\wn}{#5}%
    \pgfmathsetmacro{\inset}{(\ww-\wn)/2}%
    \ifnum#6=1
      \fill[#7] (\cx,\cy) -- (\cx+\ww,\cy+\inset) -- (\cx+\ww,\cy+\h-\inset) -- (\cx,\cy+\h) -- cycle;
      \draw[black,thin] (\cx,\cy) -- (\cx+\ww,\cy+\inset) -- (\cx+\ww,\cy+\h-\inset) -- (\cx,\cy+\h) -- cycle;
    \else
      \fill[#7] (\cx,\cy+\inset) -- (\cx+\ww,\cy) -- (\cx+\ww,\cy+\h) -- (\cx,\cy+\h-\inset) -- cycle;
      \draw[black,thin] (\cx,\cy+\inset) -- (\cx+\ww,\cy) -- (\cx+\ww,\cy+\h) -- (\cx,\cy+\h-\inset) -- cycle;
    \fi
    \node[align=center] at (\cx+\ww/2,\cy+\h/2) {#8};
  \end{scope}
}

% ---------------------------------------------------------------
% \archbowtie{x}{y}{w}{h}{pinch}{color}{label} -- Medverse-style
% single-row "dual U-Net" hourglass (funnel-in then funnel-out).
% ---------------------------------------------------------------
\newcommand{\archbowtie}[7]{%
  \begin{scope}
    \pgfmathsetmacro{\cx}{#1}\pgfmathsetmacro{\cy}{#2}%
    \pgfmathsetmacro{\w}{#3}\pgfmathsetmacro{\h}{#4}\pgfmathsetmacro{\p}{#5}%
    \pgfmathsetmacro{\inset}{(\h-\p)/2}%
    \fill[#6] (\cx,\cy) -- (\cx+\w/2,\cy+\inset) -- (\cx+\w,\cy) -- (\cx+\w,\cy+\h)
      -- (\cx+\w/2,\cy+\h-\inset) -- (\cx,\cy+\h) -- cycle;
    \draw[black,thin] (\cx,\cy) -- (\cx+\w/2,\cy+\inset) -- (\cx+\w,\cy) -- (\cx+\w,\cy+\h)
      -- (\cx+\w/2,\cy+\h-\inset) -- (\cx,\cy+\h) -- cycle;
    \node[align=center] at (\cx+\w/2,\cy+\h/2) {#7};
  \end{scope}
}

% ---------------------------------------------------------------
% \archtokengrid{x}{y}{rows}{cols}{cell}{color} -- grid of patch tokens
% ---------------------------------------------------------------
\newcommand{\archtokengrid}[6]{%
  \begin{scope}
    \pgfmathsetmacro{\cx}{#1}\pgfmathsetmacro{\cy}{#2}\pgfmathsetmacro{\cell}{#5}%
    \pgfmathtruncatemacro{\rmax}{#3-1}%
    \pgfmathtruncatemacro{\cmax}{#4-1}%
    \foreach \r in {0,...,\rmax}{%
      \foreach \c in {0,...,\cmax}{%
        \fill[#6] ({\cx+\c*\cell},{\cy+\r*\cell}) rectangle ++(\cell*0.82,\cell*0.82);
        \draw[black,thin] ({\cx+\c*\cell},{\cy+\r*\cell}) rectangle ++(\cell*0.82,\cell*0.82);
      }%
    }%
  \end{scope}
}

% ---------------------------------------------------------------
% \archarrow{start}{end}{color} -- block arrow between stages
% ---------------------------------------------------------------
% #1 and #2 are full coordinates including parentheses, e.g. (1,2)
\newcommand{\archarrow}[3]{%
  \draw[-{Latex[length=2.4mm,width=2mm]},line width=2.2pt,#3] #1 -- #2;
}

% ---------------------------------------------------------------
% Dataflow-arrow vocabulary -- opt-in via \archarrowset{compact|detailed},
% called once near the top of a figure (after % Shared TikZ macros for the architecture figures (Medverse vs. Ours).
% Load with % Shared TikZ macros for the architecture figures (Medverse vs. Ours).
% Load with % Shared TikZ macros for the architecture figures (Medverse vs. Ours).
% Load with \input{arch-macros.tex} after \usepackage{tikz,ifthen} and
% \usetikzlibrary{arrows.meta,calc,positioning}. All macro names are
% prefixed `arch` to avoid clashing with pgf/tikz internals.

% ---- palette (approximating the PowerPoint drafts) ----
\definecolor{archTargetGreen}{HTML}{9BC7AC}
\definecolor{archContextPink}{HTML}{E8B8B0}
\definecolor{archAttnBlue}{HTML}{BBD6EC}
\definecolor{archEncGray}{HTML}{E4E4E4}
\definecolor{archArrowPale}{HTML}{CFE6D6}
\definecolor{archCtBg}{HTML}{2B2B2B}
\definecolor{archCtBlob}{HTML}{8C8C8C}
\definecolor{archCtxYellow}{HTML}{E8D24C}
% darker, more saturated siblings of archTargetGreen/archContextPink --
% pastel fills read poorly as 1pt strokes, and mask/label content needs to
% read as "which volume owns this" at a glance, so masks/labels and any
% line (arrow, brace) touching that volume's row use these instead of the
% pastel role color or the old flat archMaskGold.
\definecolor{archTargetAccent}{HTML}{2F6B4A}  % target: darker green
\definecolor{archContextAccent}{HTML}{A8453E} % context: darker red

% fake-isometric depth skew shared by every cube
\newcommand{\archdx}{0.30}
\newcommand{\archdy}{0.20}

% ---- shared text-role macros ----
% Every arch_*.tex figure is compiled standalone at the SAME base size
% (documentclass[tikz,border=4pt,12pt]{standalone}, matching the thesis
% body's 12pt), so \normalsize etc. mean the same absolute point size in
% every figure. Using these three named roles instead of ad hoc \Large /
% \bfseries\normalsize / \scriptsize per figure keeps "a box title" (etc.)
% the same size everywhere it appears, regardless of which figure it's in
% or how that figure's canvas happens to be scaled by \includegraphics.
%   \archtitle    -- figure/method/panel titles, the "xL" repeat marker
%   \archboxlabel -- sub-block header labels (Fuse, Encoder, attention, ...)
%   \archnote     -- small secondary annotations (placeholder captions)
% Plain, unlabeled node text (cube/token labels, row/lane labels) stays at
% the bare \normalsize default -- no macro needed for that role.
%
% DECLARATIONS, not argument-taking macros -- deliberately, so they can
% precede a raw `\\` inside a tikz node's text (as `{\archboxlabel line
% one\\line two}`) without hiding that `\\` one macro-expansion level deep.
% A `\\` buried inside a wrapped {\archboxlabel{...}} argument breaks
% TikZ's own align=center line-splitting (it scans for `\\` before the
% argument expands) -- use these the same way `\bfseries`/`\Large` are
% already used directly inline elsewhere in these figures.
\newcommand{\archtitle}{\Large\bfseries}
\newcommand{\archboxlabel}{\bfseries\normalsize}
\newcommand{\archnote}{\scriptsize}

% ---------------------------------------------------------------
% \archvolcube{x}{y}{w}{h}{topcolor}{content}{label}
%   content: ct | ctmask | predmask | <flat color name>
% ---------------------------------------------------------------
\newcommand{\archvolcube}[7]{%
  \begin{scope}
    \pgfmathsetmacro{\cx}{#1}\pgfmathsetmacro{\cy}{#2}%
    \pgfmathsetmacro{\w}{#3}\pgfmathsetmacro{\h}{#4}%
    \fill[#5] (\cx,\cy+\h) -- ++(\archdx,\archdy) -- ++(\w,0) -- ++(-\archdx,-\archdy) -- cycle;
    \fill[#5!60!black] (\cx+\w,\cy) -- ++(\archdx,\archdy) -- ++(0,\h) -- ++(-\archdx,-\archdy) -- cycle;
    \begin{scope}
      \clip (\cx,\cy) rectangle ++(\w,\h);
      \ifthenelse{\equal{#6}{ct}}{%
        \fill[archCtBg] (\cx,\cy) rectangle ++(\w,\h);
        \fill[archCtBlob] plot[smooth cycle] coordinates
          {(\cx+.25*\w,\cy+.55*\h) (\cx+.45*\w,\cy+.75*\h) (\cx+.7*\w,\cy+.6*\h)
           (\cx+.6*\w,\cy+.3*\h) (\cx+.35*\w,\cy+.25*\h)};
      }{\ifthenelse{\equal{#6}{ctmask}}{%
        \fill[archCtBg] (\cx,\cy) rectangle ++(\w,\h);
        \fill[archCtBlob] plot[smooth cycle] coordinates
          {(\cx+.25*\w,\cy+.55*\h) (\cx+.45*\w,\cy+.75*\h) (\cx+.7*\w,\cy+.6*\h)
           (\cx+.6*\w,\cy+.3*\h) (\cx+.35*\w,\cy+.25*\h)};
        \fill[archCtxYellow] plot[smooth cycle] coordinates
          {(\cx+.42*\w,\cy+.42*\h) (\cx+.52*\w,\cy+.5*\h) (\cx+.5*\w,\cy+.35*\h) (\cx+.44*\w,\cy+.3*\h)};
      }{\ifthenelse{\equal{#6}{predmask}}{%
        \fill[black] (\cx,\cy) rectangle ++(\w,\h);
        \fill[white] plot[smooth cycle] coordinates
          {(\cx+.4*\w,\cy+.65*\h) (\cx+.6*\w,\cy+.7*\h) (\cx+.62*\w,\cy+.45*\h)
           (\cx+.48*\w,\cy+.3*\h) (\cx+.38*\w,\cy+.4*\h)};
      }{%
        \fill[#6] (\cx,\cy) rectangle ++(\w,\h);
      }}}%
    \end{scope}
    \draw[black,thin] (\cx,\cy) rectangle ++(\w,\h);
    \draw[black,thin] (\cx,\cy+\h) -- ++(\archdx,\archdy) -- ++(\w,0) -- ++(-\archdx,-\archdy) -- cycle;
    \draw[black,thin] (\cx+\w,\cy) -- ++(\archdx,\archdy) -- ++(0,\h) -- ++(-\archdx,-\archdy) -- cycle;
    \node[anchor=north,align=center] at (\cx+\w/2,\cy-0.1) {#7};
  \end{scope}
}

% ---------------------------------------------------------------
% \archvolcubeimg{x}{y}{w}{h}{topcolor}{imgpath}{label} -- same cube as
% \archvolcube, but the front face shows a real raster image (clipped to
% exactly fill the face, aspect ratio not preserved) instead of one of
% \archvolcube's procedural ct/ctmask/predmask blobs. Needs graphicx
% loaded by the including figure. A separate macro rather than a new
% \archvolcube content mode, so existing ct/ctmask/predmask figures are
% untouched.
% ---------------------------------------------------------------
\newcommand{\archvolcubeimg}[7]{%
  \begin{scope}
    \pgfmathsetmacro{\cx}{#1}\pgfmathsetmacro{\cy}{#2}%
    \pgfmathsetmacro{\w}{#3}\pgfmathsetmacro{\h}{#4}%
    \fill[#5] (\cx,\cy+\h) -- ++(\archdx,\archdy) -- ++(\w,0) -- ++(-\archdx,-\archdy) -- cycle;
    \fill[#5!60!black] (\cx+\w,\cy) -- ++(\archdx,\archdy) -- ++(0,\h) -- ++(-\archdx,-\archdy) -- cycle;
    \begin{scope}
      \clip (\cx,\cy) rectangle ++(\w,\h);
      \node[anchor=south west, inner sep=0pt] at (\cx,\cy)
        {\includegraphics[width=\w cm,height=\h cm]{#6}};
    \end{scope}
    \draw[black,thin] (\cx,\cy) rectangle ++(\w,\h);
    \draw[black,thin] (\cx,\cy+\h) -- ++(\archdx,\archdy) -- ++(\w,0) -- ++(-\archdx,-\archdy) -- cycle;
    \draw[black,thin] (\cx+\w,\cy) -- ++(\archdx,\archdy) -- ++(0,\h) -- ++(-\archdx,-\archdy) -- cycle;
    \node[anchor=north,align=center] at (\cx+\w/2,\cy-0.1) {#7};
  \end{scope}
}

% ---------------------------------------------------------------
% \archfeatcube{x}{y}{w}{h}{color}{label} -- flat-colored feature volume
% ---------------------------------------------------------------
\newcommand{\archfeatcube}[6]{%
  \begin{scope}
    \pgfmathsetmacro{\cx}{#1}\pgfmathsetmacro{\cy}{#2}%
    \pgfmathsetmacro{\w}{#3}\pgfmathsetmacro{\h}{#4}%
    \fill[#5] (\cx,\cy) rectangle ++(\w,\h);
    \fill[#5!80!white] (\cx,\cy+\h) -- ++(\archdx,\archdy) -- ++(\w,0) -- ++(-\archdx,-\archdy) -- cycle;
    \fill[#5!60!black] (\cx+\w,\cy) -- ++(\archdx,\archdy) -- ++(0,\h) -- ++(-\archdx,-\archdy) -- cycle;
    \draw[black,thin] (\cx,\cy) rectangle ++(\w,\h);
    \draw[black,thin] (\cx,\cy+\h) -- ++(\archdx,\archdy) -- ++(\w,0) -- ++(-\archdx,-\archdy) -- cycle;
    \draw[black,thin] (\cx+\w,\cy) -- ++(\archdx,\archdy) -- ++(0,\h) -- ++(-\archdx,-\archdy) -- cycle;
    \node[align=center] at (\cx+\w/2,\cy+\h/2) {#6};
  \end{scope}
}

% ---------------------------------------------------------------
% \archfunnel{x}{y}{h}{wwide}{wnarrow}{dir}{color}{label}
%   dir=1  : wide on left, narrow on right  (encoder)
%   dir=-1 : narrow on left, wide on right  (decoder)
% ---------------------------------------------------------------
\newcommand{\archfunnel}[8]{%
  \begin{scope}
    \pgfmathsetmacro{\cx}{#1}\pgfmathsetmacro{\cy}{#2}\pgfmathsetmacro{\h}{#3}%
    \pgfmathsetmacro{\ww}{#4}\pgfmathsetmacro{\wn}{#5}%
    \pgfmathsetmacro{\inset}{(\ww-\wn)/2}%
    \ifnum#6=1
      \fill[#7] (\cx,\cy) -- (\cx+\ww,\cy+\inset) -- (\cx+\ww,\cy+\h-\inset) -- (\cx,\cy+\h) -- cycle;
      \draw[black,thin] (\cx,\cy) -- (\cx+\ww,\cy+\inset) -- (\cx+\ww,\cy+\h-\inset) -- (\cx,\cy+\h) -- cycle;
    \else
      \fill[#7] (\cx,\cy+\inset) -- (\cx+\ww,\cy) -- (\cx+\ww,\cy+\h) -- (\cx,\cy+\h-\inset) -- cycle;
      \draw[black,thin] (\cx,\cy+\inset) -- (\cx+\ww,\cy) -- (\cx+\ww,\cy+\h) -- (\cx,\cy+\h-\inset) -- cycle;
    \fi
    \node[align=center] at (\cx+\ww/2,\cy+\h/2) {#8};
  \end{scope}
}

% ---------------------------------------------------------------
% \archbowtie{x}{y}{w}{h}{pinch}{color}{label} -- Medverse-style
% single-row "dual U-Net" hourglass (funnel-in then funnel-out).
% ---------------------------------------------------------------
\newcommand{\archbowtie}[7]{%
  \begin{scope}
    \pgfmathsetmacro{\cx}{#1}\pgfmathsetmacro{\cy}{#2}%
    \pgfmathsetmacro{\w}{#3}\pgfmathsetmacro{\h}{#4}\pgfmathsetmacro{\p}{#5}%
    \pgfmathsetmacro{\inset}{(\h-\p)/2}%
    \fill[#6] (\cx,\cy) -- (\cx+\w/2,\cy+\inset) -- (\cx+\w,\cy) -- (\cx+\w,\cy+\h)
      -- (\cx+\w/2,\cy+\h-\inset) -- (\cx,\cy+\h) -- cycle;
    \draw[black,thin] (\cx,\cy) -- (\cx+\w/2,\cy+\inset) -- (\cx+\w,\cy) -- (\cx+\w,\cy+\h)
      -- (\cx+\w/2,\cy+\h-\inset) -- (\cx,\cy+\h) -- cycle;
    \node[align=center] at (\cx+\w/2,\cy+\h/2) {#7};
  \end{scope}
}

% ---------------------------------------------------------------
% \archtokengrid{x}{y}{rows}{cols}{cell}{color} -- grid of patch tokens
% ---------------------------------------------------------------
\newcommand{\archtokengrid}[6]{%
  \begin{scope}
    \pgfmathsetmacro{\cx}{#1}\pgfmathsetmacro{\cy}{#2}\pgfmathsetmacro{\cell}{#5}%
    \pgfmathtruncatemacro{\rmax}{#3-1}%
    \pgfmathtruncatemacro{\cmax}{#4-1}%
    \foreach \r in {0,...,\rmax}{%
      \foreach \c in {0,...,\cmax}{%
        \fill[#6] ({\cx+\c*\cell},{\cy+\r*\cell}) rectangle ++(\cell*0.82,\cell*0.82);
        \draw[black,thin] ({\cx+\c*\cell},{\cy+\r*\cell}) rectangle ++(\cell*0.82,\cell*0.82);
      }%
    }%
  \end{scope}
}

% ---------------------------------------------------------------
% \archarrow{start}{end}{color} -- block arrow between stages
% ---------------------------------------------------------------
% #1 and #2 are full coordinates including parentheses, e.g. (1,2)
\newcommand{\archarrow}[3]{%
  \draw[-{Latex[length=2.4mm,width=2mm]},line width=2.2pt,#3] #1 -- #2;
}

% ---------------------------------------------------------------
% Dataflow-arrow vocabulary -- opt-in via \archarrowset{compact|detailed},
% called once near the top of a figure (after \input{arch-macros.tex}).
% \archarrow keeps its classic small size by default for any arch_*.tex
% figure that never calls \archarrowset -- this block is purely additive.
%   archflowstyle (tikz key, NOT a macro -- use bare, no backslash) --
%                                  current draw options (head + line
%                                  width), for a hand-rolled elbow \draw
%                                  that \archarrow's straight "--" can't
%                                  express (e.g. "-|" routing). Must be a
%                                  /.style key rather than a \def'd macro:
%                                  pgf's arrow-tip parser doesn't split on
%                                  a comma embedded in a plain macro's
%                                  expansion when used as a \draw option.
%   \archflowNeutral/Target/Context -- the only 3 colors a plain dataflow
%                                  arrow should use: generic same-role
%                                  passthrough; target's own real output
%                                  leaving the decoder; a real label/mask
%                                  tensor entering fusion or attention
%                                  (context's ground truth, or target's
%                                  own real prior fed back from a
%                                  previous cascade level -- same role
%                                  either way: real mask data, not a
%                                  feature)
%   \archflow{p1}{p2}{role}     -- \archarrow wrapper keyed by role name
%                                  (role in {neutral, target, context})
% ---------------------------------------------------------------
\newcommand{\archarrowset}[1]{%
  \ifthenelse{\equal{#1}{detailed}}{%
    \tikzset{archflowstyle/.style={-{Latex[length=3.6mm,width=3.2mm]},line width=3.2pt}}%
  }{%
    \tikzset{archflowstyle/.style={-{Latex[length=2.4mm,width=2mm]},line width=2.2pt}}%
  }%
  \renewcommand{\archarrow}[3]{\draw[archflowstyle,##3] ##1 -- ##2;}%
}
\newcommand{\archflowNeutral}{black!30}
\newcommand{\archflowTarget}{archTargetAccent}
\newcommand{\archflowContext}{archContextAccent}
\newcommand{\archflow}[3]{%
  \ifthenelse{\equal{#3}{target}}{\archarrow{#1}{#2}{\archflowTarget}}{%
  \ifthenelse{\equal{#3}{context}}{\archarrow{#1}{#2}{\archflowContext}}{%
  \archarrow{#1}{#2}{\archflowNeutral}}}%
}
 after \usepackage{tikz,ifthen} and
% \usetikzlibrary{arrows.meta,calc,positioning}. All macro names are
% prefixed `arch` to avoid clashing with pgf/tikz internals.

% ---- palette (approximating the PowerPoint drafts) ----
\definecolor{archTargetGreen}{HTML}{9BC7AC}
\definecolor{archContextPink}{HTML}{E8B8B0}
\definecolor{archAttnBlue}{HTML}{BBD6EC}
\definecolor{archEncGray}{HTML}{E4E4E4}
\definecolor{archArrowPale}{HTML}{CFE6D6}
\definecolor{archCtBg}{HTML}{2B2B2B}
\definecolor{archCtBlob}{HTML}{8C8C8C}
\definecolor{archCtxYellow}{HTML}{E8D24C}
% darker, more saturated siblings of archTargetGreen/archContextPink --
% pastel fills read poorly as 1pt strokes, and mask/label content needs to
% read as "which volume owns this" at a glance, so masks/labels and any
% line (arrow, brace) touching that volume's row use these instead of the
% pastel role color or the old flat archMaskGold.
\definecolor{archTargetAccent}{HTML}{2F6B4A}  % target: darker green
\definecolor{archContextAccent}{HTML}{A8453E} % context: darker red

% fake-isometric depth skew shared by every cube
\newcommand{\archdx}{0.30}
\newcommand{\archdy}{0.20}

% ---- shared text-role macros ----
% Every arch_*.tex figure is compiled standalone at the SAME base size
% (documentclass[tikz,border=4pt,12pt]{standalone}, matching the thesis
% body's 12pt), so \normalsize etc. mean the same absolute point size in
% every figure. Using these three named roles instead of ad hoc \Large /
% \bfseries\normalsize / \scriptsize per figure keeps "a box title" (etc.)
% the same size everywhere it appears, regardless of which figure it's in
% or how that figure's canvas happens to be scaled by \includegraphics.
%   \archtitle    -- figure/method/panel titles, the "xL" repeat marker
%   \archboxlabel -- sub-block header labels (Fuse, Encoder, attention, ...)
%   \archnote     -- small secondary annotations (placeholder captions)
% Plain, unlabeled node text (cube/token labels, row/lane labels) stays at
% the bare \normalsize default -- no macro needed for that role.
%
% DECLARATIONS, not argument-taking macros -- deliberately, so they can
% precede a raw `\\` inside a tikz node's text (as `{\archboxlabel line
% one\\line two}`) without hiding that `\\` one macro-expansion level deep.
% A `\\` buried inside a wrapped {\archboxlabel{...}} argument breaks
% TikZ's own align=center line-splitting (it scans for `\\` before the
% argument expands) -- use these the same way `\bfseries`/`\Large` are
% already used directly inline elsewhere in these figures.
\newcommand{\archtitle}{\Large\bfseries}
\newcommand{\archboxlabel}{\bfseries\normalsize}
\newcommand{\archnote}{\scriptsize}

% ---------------------------------------------------------------
% \archvolcube{x}{y}{w}{h}{topcolor}{content}{label}
%   content: ct | ctmask | predmask | <flat color name>
% ---------------------------------------------------------------
\newcommand{\archvolcube}[7]{%
  \begin{scope}
    \pgfmathsetmacro{\cx}{#1}\pgfmathsetmacro{\cy}{#2}%
    \pgfmathsetmacro{\w}{#3}\pgfmathsetmacro{\h}{#4}%
    \fill[#5] (\cx,\cy+\h) -- ++(\archdx,\archdy) -- ++(\w,0) -- ++(-\archdx,-\archdy) -- cycle;
    \fill[#5!60!black] (\cx+\w,\cy) -- ++(\archdx,\archdy) -- ++(0,\h) -- ++(-\archdx,-\archdy) -- cycle;
    \begin{scope}
      \clip (\cx,\cy) rectangle ++(\w,\h);
      \ifthenelse{\equal{#6}{ct}}{%
        \fill[archCtBg] (\cx,\cy) rectangle ++(\w,\h);
        \fill[archCtBlob] plot[smooth cycle] coordinates
          {(\cx+.25*\w,\cy+.55*\h) (\cx+.45*\w,\cy+.75*\h) (\cx+.7*\w,\cy+.6*\h)
           (\cx+.6*\w,\cy+.3*\h) (\cx+.35*\w,\cy+.25*\h)};
      }{\ifthenelse{\equal{#6}{ctmask}}{%
        \fill[archCtBg] (\cx,\cy) rectangle ++(\w,\h);
        \fill[archCtBlob] plot[smooth cycle] coordinates
          {(\cx+.25*\w,\cy+.55*\h) (\cx+.45*\w,\cy+.75*\h) (\cx+.7*\w,\cy+.6*\h)
           (\cx+.6*\w,\cy+.3*\h) (\cx+.35*\w,\cy+.25*\h)};
        \fill[archCtxYellow] plot[smooth cycle] coordinates
          {(\cx+.42*\w,\cy+.42*\h) (\cx+.52*\w,\cy+.5*\h) (\cx+.5*\w,\cy+.35*\h) (\cx+.44*\w,\cy+.3*\h)};
      }{\ifthenelse{\equal{#6}{predmask}}{%
        \fill[black] (\cx,\cy) rectangle ++(\w,\h);
        \fill[white] plot[smooth cycle] coordinates
          {(\cx+.4*\w,\cy+.65*\h) (\cx+.6*\w,\cy+.7*\h) (\cx+.62*\w,\cy+.45*\h)
           (\cx+.48*\w,\cy+.3*\h) (\cx+.38*\w,\cy+.4*\h)};
      }{%
        \fill[#6] (\cx,\cy) rectangle ++(\w,\h);
      }}}%
    \end{scope}
    \draw[black,thin] (\cx,\cy) rectangle ++(\w,\h);
    \draw[black,thin] (\cx,\cy+\h) -- ++(\archdx,\archdy) -- ++(\w,0) -- ++(-\archdx,-\archdy) -- cycle;
    \draw[black,thin] (\cx+\w,\cy) -- ++(\archdx,\archdy) -- ++(0,\h) -- ++(-\archdx,-\archdy) -- cycle;
    \node[anchor=north,align=center] at (\cx+\w/2,\cy-0.1) {#7};
  \end{scope}
}

% ---------------------------------------------------------------
% \archvolcubeimg{x}{y}{w}{h}{topcolor}{imgpath}{label} -- same cube as
% \archvolcube, but the front face shows a real raster image (clipped to
% exactly fill the face, aspect ratio not preserved) instead of one of
% \archvolcube's procedural ct/ctmask/predmask blobs. Needs graphicx
% loaded by the including figure. A separate macro rather than a new
% \archvolcube content mode, so existing ct/ctmask/predmask figures are
% untouched.
% ---------------------------------------------------------------
\newcommand{\archvolcubeimg}[7]{%
  \begin{scope}
    \pgfmathsetmacro{\cx}{#1}\pgfmathsetmacro{\cy}{#2}%
    \pgfmathsetmacro{\w}{#3}\pgfmathsetmacro{\h}{#4}%
    \fill[#5] (\cx,\cy+\h) -- ++(\archdx,\archdy) -- ++(\w,0) -- ++(-\archdx,-\archdy) -- cycle;
    \fill[#5!60!black] (\cx+\w,\cy) -- ++(\archdx,\archdy) -- ++(0,\h) -- ++(-\archdx,-\archdy) -- cycle;
    \begin{scope}
      \clip (\cx,\cy) rectangle ++(\w,\h);
      \node[anchor=south west, inner sep=0pt] at (\cx,\cy)
        {\includegraphics[width=\w cm,height=\h cm]{#6}};
    \end{scope}
    \draw[black,thin] (\cx,\cy) rectangle ++(\w,\h);
    \draw[black,thin] (\cx,\cy+\h) -- ++(\archdx,\archdy) -- ++(\w,0) -- ++(-\archdx,-\archdy) -- cycle;
    \draw[black,thin] (\cx+\w,\cy) -- ++(\archdx,\archdy) -- ++(0,\h) -- ++(-\archdx,-\archdy) -- cycle;
    \node[anchor=north,align=center] at (\cx+\w/2,\cy-0.1) {#7};
  \end{scope}
}

% ---------------------------------------------------------------
% \archfeatcube{x}{y}{w}{h}{color}{label} -- flat-colored feature volume
% ---------------------------------------------------------------
\newcommand{\archfeatcube}[6]{%
  \begin{scope}
    \pgfmathsetmacro{\cx}{#1}\pgfmathsetmacro{\cy}{#2}%
    \pgfmathsetmacro{\w}{#3}\pgfmathsetmacro{\h}{#4}%
    \fill[#5] (\cx,\cy) rectangle ++(\w,\h);
    \fill[#5!80!white] (\cx,\cy+\h) -- ++(\archdx,\archdy) -- ++(\w,0) -- ++(-\archdx,-\archdy) -- cycle;
    \fill[#5!60!black] (\cx+\w,\cy) -- ++(\archdx,\archdy) -- ++(0,\h) -- ++(-\archdx,-\archdy) -- cycle;
    \draw[black,thin] (\cx,\cy) rectangle ++(\w,\h);
    \draw[black,thin] (\cx,\cy+\h) -- ++(\archdx,\archdy) -- ++(\w,0) -- ++(-\archdx,-\archdy) -- cycle;
    \draw[black,thin] (\cx+\w,\cy) -- ++(\archdx,\archdy) -- ++(0,\h) -- ++(-\archdx,-\archdy) -- cycle;
    \node[align=center] at (\cx+\w/2,\cy+\h/2) {#6};
  \end{scope}
}

% ---------------------------------------------------------------
% \archfunnel{x}{y}{h}{wwide}{wnarrow}{dir}{color}{label}
%   dir=1  : wide on left, narrow on right  (encoder)
%   dir=-1 : narrow on left, wide on right  (decoder)
% ---------------------------------------------------------------
\newcommand{\archfunnel}[8]{%
  \begin{scope}
    \pgfmathsetmacro{\cx}{#1}\pgfmathsetmacro{\cy}{#2}\pgfmathsetmacro{\h}{#3}%
    \pgfmathsetmacro{\ww}{#4}\pgfmathsetmacro{\wn}{#5}%
    \pgfmathsetmacro{\inset}{(\ww-\wn)/2}%
    \ifnum#6=1
      \fill[#7] (\cx,\cy) -- (\cx+\ww,\cy+\inset) -- (\cx+\ww,\cy+\h-\inset) -- (\cx,\cy+\h) -- cycle;
      \draw[black,thin] (\cx,\cy) -- (\cx+\ww,\cy+\inset) -- (\cx+\ww,\cy+\h-\inset) -- (\cx,\cy+\h) -- cycle;
    \else
      \fill[#7] (\cx,\cy+\inset) -- (\cx+\ww,\cy) -- (\cx+\ww,\cy+\h) -- (\cx,\cy+\h-\inset) -- cycle;
      \draw[black,thin] (\cx,\cy+\inset) -- (\cx+\ww,\cy) -- (\cx+\ww,\cy+\h) -- (\cx,\cy+\h-\inset) -- cycle;
    \fi
    \node[align=center] at (\cx+\ww/2,\cy+\h/2) {#8};
  \end{scope}
}

% ---------------------------------------------------------------
% \archbowtie{x}{y}{w}{h}{pinch}{color}{label} -- Medverse-style
% single-row "dual U-Net" hourglass (funnel-in then funnel-out).
% ---------------------------------------------------------------
\newcommand{\archbowtie}[7]{%
  \begin{scope}
    \pgfmathsetmacro{\cx}{#1}\pgfmathsetmacro{\cy}{#2}%
    \pgfmathsetmacro{\w}{#3}\pgfmathsetmacro{\h}{#4}\pgfmathsetmacro{\p}{#5}%
    \pgfmathsetmacro{\inset}{(\h-\p)/2}%
    \fill[#6] (\cx,\cy) -- (\cx+\w/2,\cy+\inset) -- (\cx+\w,\cy) -- (\cx+\w,\cy+\h)
      -- (\cx+\w/2,\cy+\h-\inset) -- (\cx,\cy+\h) -- cycle;
    \draw[black,thin] (\cx,\cy) -- (\cx+\w/2,\cy+\inset) -- (\cx+\w,\cy) -- (\cx+\w,\cy+\h)
      -- (\cx+\w/2,\cy+\h-\inset) -- (\cx,\cy+\h) -- cycle;
    \node[align=center] at (\cx+\w/2,\cy+\h/2) {#7};
  \end{scope}
}

% ---------------------------------------------------------------
% \archtokengrid{x}{y}{rows}{cols}{cell}{color} -- grid of patch tokens
% ---------------------------------------------------------------
\newcommand{\archtokengrid}[6]{%
  \begin{scope}
    \pgfmathsetmacro{\cx}{#1}\pgfmathsetmacro{\cy}{#2}\pgfmathsetmacro{\cell}{#5}%
    \pgfmathtruncatemacro{\rmax}{#3-1}%
    \pgfmathtruncatemacro{\cmax}{#4-1}%
    \foreach \r in {0,...,\rmax}{%
      \foreach \c in {0,...,\cmax}{%
        \fill[#6] ({\cx+\c*\cell},{\cy+\r*\cell}) rectangle ++(\cell*0.82,\cell*0.82);
        \draw[black,thin] ({\cx+\c*\cell},{\cy+\r*\cell}) rectangle ++(\cell*0.82,\cell*0.82);
      }%
    }%
  \end{scope}
}

% ---------------------------------------------------------------
% \archarrow{start}{end}{color} -- block arrow between stages
% ---------------------------------------------------------------
% #1 and #2 are full coordinates including parentheses, e.g. (1,2)
\newcommand{\archarrow}[3]{%
  \draw[-{Latex[length=2.4mm,width=2mm]},line width=2.2pt,#3] #1 -- #2;
}

% ---------------------------------------------------------------
% Dataflow-arrow vocabulary -- opt-in via \archarrowset{compact|detailed},
% called once near the top of a figure (after % Shared TikZ macros for the architecture figures (Medverse vs. Ours).
% Load with \input{arch-macros.tex} after \usepackage{tikz,ifthen} and
% \usetikzlibrary{arrows.meta,calc,positioning}. All macro names are
% prefixed `arch` to avoid clashing with pgf/tikz internals.

% ---- palette (approximating the PowerPoint drafts) ----
\definecolor{archTargetGreen}{HTML}{9BC7AC}
\definecolor{archContextPink}{HTML}{E8B8B0}
\definecolor{archAttnBlue}{HTML}{BBD6EC}
\definecolor{archEncGray}{HTML}{E4E4E4}
\definecolor{archArrowPale}{HTML}{CFE6D6}
\definecolor{archCtBg}{HTML}{2B2B2B}
\definecolor{archCtBlob}{HTML}{8C8C8C}
\definecolor{archCtxYellow}{HTML}{E8D24C}
% darker, more saturated siblings of archTargetGreen/archContextPink --
% pastel fills read poorly as 1pt strokes, and mask/label content needs to
% read as "which volume owns this" at a glance, so masks/labels and any
% line (arrow, brace) touching that volume's row use these instead of the
% pastel role color or the old flat archMaskGold.
\definecolor{archTargetAccent}{HTML}{2F6B4A}  % target: darker green
\definecolor{archContextAccent}{HTML}{A8453E} % context: darker red

% fake-isometric depth skew shared by every cube
\newcommand{\archdx}{0.30}
\newcommand{\archdy}{0.20}

% ---- shared text-role macros ----
% Every arch_*.tex figure is compiled standalone at the SAME base size
% (documentclass[tikz,border=4pt,12pt]{standalone}, matching the thesis
% body's 12pt), so \normalsize etc. mean the same absolute point size in
% every figure. Using these three named roles instead of ad hoc \Large /
% \bfseries\normalsize / \scriptsize per figure keeps "a box title" (etc.)
% the same size everywhere it appears, regardless of which figure it's in
% or how that figure's canvas happens to be scaled by \includegraphics.
%   \archtitle    -- figure/method/panel titles, the "xL" repeat marker
%   \archboxlabel -- sub-block header labels (Fuse, Encoder, attention, ...)
%   \archnote     -- small secondary annotations (placeholder captions)
% Plain, unlabeled node text (cube/token labels, row/lane labels) stays at
% the bare \normalsize default -- no macro needed for that role.
%
% DECLARATIONS, not argument-taking macros -- deliberately, so they can
% precede a raw `\\` inside a tikz node's text (as `{\archboxlabel line
% one\\line two}`) without hiding that `\\` one macro-expansion level deep.
% A `\\` buried inside a wrapped {\archboxlabel{...}} argument breaks
% TikZ's own align=center line-splitting (it scans for `\\` before the
% argument expands) -- use these the same way `\bfseries`/`\Large` are
% already used directly inline elsewhere in these figures.
\newcommand{\archtitle}{\Large\bfseries}
\newcommand{\archboxlabel}{\bfseries\normalsize}
\newcommand{\archnote}{\scriptsize}

% ---------------------------------------------------------------
% \archvolcube{x}{y}{w}{h}{topcolor}{content}{label}
%   content: ct | ctmask | predmask | <flat color name>
% ---------------------------------------------------------------
\newcommand{\archvolcube}[7]{%
  \begin{scope}
    \pgfmathsetmacro{\cx}{#1}\pgfmathsetmacro{\cy}{#2}%
    \pgfmathsetmacro{\w}{#3}\pgfmathsetmacro{\h}{#4}%
    \fill[#5] (\cx,\cy+\h) -- ++(\archdx,\archdy) -- ++(\w,0) -- ++(-\archdx,-\archdy) -- cycle;
    \fill[#5!60!black] (\cx+\w,\cy) -- ++(\archdx,\archdy) -- ++(0,\h) -- ++(-\archdx,-\archdy) -- cycle;
    \begin{scope}
      \clip (\cx,\cy) rectangle ++(\w,\h);
      \ifthenelse{\equal{#6}{ct}}{%
        \fill[archCtBg] (\cx,\cy) rectangle ++(\w,\h);
        \fill[archCtBlob] plot[smooth cycle] coordinates
          {(\cx+.25*\w,\cy+.55*\h) (\cx+.45*\w,\cy+.75*\h) (\cx+.7*\w,\cy+.6*\h)
           (\cx+.6*\w,\cy+.3*\h) (\cx+.35*\w,\cy+.25*\h)};
      }{\ifthenelse{\equal{#6}{ctmask}}{%
        \fill[archCtBg] (\cx,\cy) rectangle ++(\w,\h);
        \fill[archCtBlob] plot[smooth cycle] coordinates
          {(\cx+.25*\w,\cy+.55*\h) (\cx+.45*\w,\cy+.75*\h) (\cx+.7*\w,\cy+.6*\h)
           (\cx+.6*\w,\cy+.3*\h) (\cx+.35*\w,\cy+.25*\h)};
        \fill[archCtxYellow] plot[smooth cycle] coordinates
          {(\cx+.42*\w,\cy+.42*\h) (\cx+.52*\w,\cy+.5*\h) (\cx+.5*\w,\cy+.35*\h) (\cx+.44*\w,\cy+.3*\h)};
      }{\ifthenelse{\equal{#6}{predmask}}{%
        \fill[black] (\cx,\cy) rectangle ++(\w,\h);
        \fill[white] plot[smooth cycle] coordinates
          {(\cx+.4*\w,\cy+.65*\h) (\cx+.6*\w,\cy+.7*\h) (\cx+.62*\w,\cy+.45*\h)
           (\cx+.48*\w,\cy+.3*\h) (\cx+.38*\w,\cy+.4*\h)};
      }{%
        \fill[#6] (\cx,\cy) rectangle ++(\w,\h);
      }}}%
    \end{scope}
    \draw[black,thin] (\cx,\cy) rectangle ++(\w,\h);
    \draw[black,thin] (\cx,\cy+\h) -- ++(\archdx,\archdy) -- ++(\w,0) -- ++(-\archdx,-\archdy) -- cycle;
    \draw[black,thin] (\cx+\w,\cy) -- ++(\archdx,\archdy) -- ++(0,\h) -- ++(-\archdx,-\archdy) -- cycle;
    \node[anchor=north,align=center] at (\cx+\w/2,\cy-0.1) {#7};
  \end{scope}
}

% ---------------------------------------------------------------
% \archvolcubeimg{x}{y}{w}{h}{topcolor}{imgpath}{label} -- same cube as
% \archvolcube, but the front face shows a real raster image (clipped to
% exactly fill the face, aspect ratio not preserved) instead of one of
% \archvolcube's procedural ct/ctmask/predmask blobs. Needs graphicx
% loaded by the including figure. A separate macro rather than a new
% \archvolcube content mode, so existing ct/ctmask/predmask figures are
% untouched.
% ---------------------------------------------------------------
\newcommand{\archvolcubeimg}[7]{%
  \begin{scope}
    \pgfmathsetmacro{\cx}{#1}\pgfmathsetmacro{\cy}{#2}%
    \pgfmathsetmacro{\w}{#3}\pgfmathsetmacro{\h}{#4}%
    \fill[#5] (\cx,\cy+\h) -- ++(\archdx,\archdy) -- ++(\w,0) -- ++(-\archdx,-\archdy) -- cycle;
    \fill[#5!60!black] (\cx+\w,\cy) -- ++(\archdx,\archdy) -- ++(0,\h) -- ++(-\archdx,-\archdy) -- cycle;
    \begin{scope}
      \clip (\cx,\cy) rectangle ++(\w,\h);
      \node[anchor=south west, inner sep=0pt] at (\cx,\cy)
        {\includegraphics[width=\w cm,height=\h cm]{#6}};
    \end{scope}
    \draw[black,thin] (\cx,\cy) rectangle ++(\w,\h);
    \draw[black,thin] (\cx,\cy+\h) -- ++(\archdx,\archdy) -- ++(\w,0) -- ++(-\archdx,-\archdy) -- cycle;
    \draw[black,thin] (\cx+\w,\cy) -- ++(\archdx,\archdy) -- ++(0,\h) -- ++(-\archdx,-\archdy) -- cycle;
    \node[anchor=north,align=center] at (\cx+\w/2,\cy-0.1) {#7};
  \end{scope}
}

% ---------------------------------------------------------------
% \archfeatcube{x}{y}{w}{h}{color}{label} -- flat-colored feature volume
% ---------------------------------------------------------------
\newcommand{\archfeatcube}[6]{%
  \begin{scope}
    \pgfmathsetmacro{\cx}{#1}\pgfmathsetmacro{\cy}{#2}%
    \pgfmathsetmacro{\w}{#3}\pgfmathsetmacro{\h}{#4}%
    \fill[#5] (\cx,\cy) rectangle ++(\w,\h);
    \fill[#5!80!white] (\cx,\cy+\h) -- ++(\archdx,\archdy) -- ++(\w,0) -- ++(-\archdx,-\archdy) -- cycle;
    \fill[#5!60!black] (\cx+\w,\cy) -- ++(\archdx,\archdy) -- ++(0,\h) -- ++(-\archdx,-\archdy) -- cycle;
    \draw[black,thin] (\cx,\cy) rectangle ++(\w,\h);
    \draw[black,thin] (\cx,\cy+\h) -- ++(\archdx,\archdy) -- ++(\w,0) -- ++(-\archdx,-\archdy) -- cycle;
    \draw[black,thin] (\cx+\w,\cy) -- ++(\archdx,\archdy) -- ++(0,\h) -- ++(-\archdx,-\archdy) -- cycle;
    \node[align=center] at (\cx+\w/2,\cy+\h/2) {#6};
  \end{scope}
}

% ---------------------------------------------------------------
% \archfunnel{x}{y}{h}{wwide}{wnarrow}{dir}{color}{label}
%   dir=1  : wide on left, narrow on right  (encoder)
%   dir=-1 : narrow on left, wide on right  (decoder)
% ---------------------------------------------------------------
\newcommand{\archfunnel}[8]{%
  \begin{scope}
    \pgfmathsetmacro{\cx}{#1}\pgfmathsetmacro{\cy}{#2}\pgfmathsetmacro{\h}{#3}%
    \pgfmathsetmacro{\ww}{#4}\pgfmathsetmacro{\wn}{#5}%
    \pgfmathsetmacro{\inset}{(\ww-\wn)/2}%
    \ifnum#6=1
      \fill[#7] (\cx,\cy) -- (\cx+\ww,\cy+\inset) -- (\cx+\ww,\cy+\h-\inset) -- (\cx,\cy+\h) -- cycle;
      \draw[black,thin] (\cx,\cy) -- (\cx+\ww,\cy+\inset) -- (\cx+\ww,\cy+\h-\inset) -- (\cx,\cy+\h) -- cycle;
    \else
      \fill[#7] (\cx,\cy+\inset) -- (\cx+\ww,\cy) -- (\cx+\ww,\cy+\h) -- (\cx,\cy+\h-\inset) -- cycle;
      \draw[black,thin] (\cx,\cy+\inset) -- (\cx+\ww,\cy) -- (\cx+\ww,\cy+\h) -- (\cx,\cy+\h-\inset) -- cycle;
    \fi
    \node[align=center] at (\cx+\ww/2,\cy+\h/2) {#8};
  \end{scope}
}

% ---------------------------------------------------------------
% \archbowtie{x}{y}{w}{h}{pinch}{color}{label} -- Medverse-style
% single-row "dual U-Net" hourglass (funnel-in then funnel-out).
% ---------------------------------------------------------------
\newcommand{\archbowtie}[7]{%
  \begin{scope}
    \pgfmathsetmacro{\cx}{#1}\pgfmathsetmacro{\cy}{#2}%
    \pgfmathsetmacro{\w}{#3}\pgfmathsetmacro{\h}{#4}\pgfmathsetmacro{\p}{#5}%
    \pgfmathsetmacro{\inset}{(\h-\p)/2}%
    \fill[#6] (\cx,\cy) -- (\cx+\w/2,\cy+\inset) -- (\cx+\w,\cy) -- (\cx+\w,\cy+\h)
      -- (\cx+\w/2,\cy+\h-\inset) -- (\cx,\cy+\h) -- cycle;
    \draw[black,thin] (\cx,\cy) -- (\cx+\w/2,\cy+\inset) -- (\cx+\w,\cy) -- (\cx+\w,\cy+\h)
      -- (\cx+\w/2,\cy+\h-\inset) -- (\cx,\cy+\h) -- cycle;
    \node[align=center] at (\cx+\w/2,\cy+\h/2) {#7};
  \end{scope}
}

% ---------------------------------------------------------------
% \archtokengrid{x}{y}{rows}{cols}{cell}{color} -- grid of patch tokens
% ---------------------------------------------------------------
\newcommand{\archtokengrid}[6]{%
  \begin{scope}
    \pgfmathsetmacro{\cx}{#1}\pgfmathsetmacro{\cy}{#2}\pgfmathsetmacro{\cell}{#5}%
    \pgfmathtruncatemacro{\rmax}{#3-1}%
    \pgfmathtruncatemacro{\cmax}{#4-1}%
    \foreach \r in {0,...,\rmax}{%
      \foreach \c in {0,...,\cmax}{%
        \fill[#6] ({\cx+\c*\cell},{\cy+\r*\cell}) rectangle ++(\cell*0.82,\cell*0.82);
        \draw[black,thin] ({\cx+\c*\cell},{\cy+\r*\cell}) rectangle ++(\cell*0.82,\cell*0.82);
      }%
    }%
  \end{scope}
}

% ---------------------------------------------------------------
% \archarrow{start}{end}{color} -- block arrow between stages
% ---------------------------------------------------------------
% #1 and #2 are full coordinates including parentheses, e.g. (1,2)
\newcommand{\archarrow}[3]{%
  \draw[-{Latex[length=2.4mm,width=2mm]},line width=2.2pt,#3] #1 -- #2;
}

% ---------------------------------------------------------------
% Dataflow-arrow vocabulary -- opt-in via \archarrowset{compact|detailed},
% called once near the top of a figure (after \input{arch-macros.tex}).
% \archarrow keeps its classic small size by default for any arch_*.tex
% figure that never calls \archarrowset -- this block is purely additive.
%   archflowstyle (tikz key, NOT a macro -- use bare, no backslash) --
%                                  current draw options (head + line
%                                  width), for a hand-rolled elbow \draw
%                                  that \archarrow's straight "--" can't
%                                  express (e.g. "-|" routing). Must be a
%                                  /.style key rather than a \def'd macro:
%                                  pgf's arrow-tip parser doesn't split on
%                                  a comma embedded in a plain macro's
%                                  expansion when used as a \draw option.
%   \archflowNeutral/Target/Context -- the only 3 colors a plain dataflow
%                                  arrow should use: generic same-role
%                                  passthrough; target's own real output
%                                  leaving the decoder; a real label/mask
%                                  tensor entering fusion or attention
%                                  (context's ground truth, or target's
%                                  own real prior fed back from a
%                                  previous cascade level -- same role
%                                  either way: real mask data, not a
%                                  feature)
%   \archflow{p1}{p2}{role}     -- \archarrow wrapper keyed by role name
%                                  (role in {neutral, target, context})
% ---------------------------------------------------------------
\newcommand{\archarrowset}[1]{%
  \ifthenelse{\equal{#1}{detailed}}{%
    \tikzset{archflowstyle/.style={-{Latex[length=3.6mm,width=3.2mm]},line width=3.2pt}}%
  }{%
    \tikzset{archflowstyle/.style={-{Latex[length=2.4mm,width=2mm]},line width=2.2pt}}%
  }%
  \renewcommand{\archarrow}[3]{\draw[archflowstyle,##3] ##1 -- ##2;}%
}
\newcommand{\archflowNeutral}{black!30}
\newcommand{\archflowTarget}{archTargetAccent}
\newcommand{\archflowContext}{archContextAccent}
\newcommand{\archflow}[3]{%
  \ifthenelse{\equal{#3}{target}}{\archarrow{#1}{#2}{\archflowTarget}}{%
  \ifthenelse{\equal{#3}{context}}{\archarrow{#1}{#2}{\archflowContext}}{%
  \archarrow{#1}{#2}{\archflowNeutral}}}%
}
).
% \archarrow keeps its classic small size by default for any arch_*.tex
% figure that never calls \archarrowset -- this block is purely additive.
%   archflowstyle (tikz key, NOT a macro -- use bare, no backslash) --
%                                  current draw options (head + line
%                                  width), for a hand-rolled elbow \draw
%                                  that \archarrow's straight "--" can't
%                                  express (e.g. "-|" routing). Must be a
%                                  /.style key rather than a \def'd macro:
%                                  pgf's arrow-tip parser doesn't split on
%                                  a comma embedded in a plain macro's
%                                  expansion when used as a \draw option.
%   \archflowNeutral/Target/Context -- the only 3 colors a plain dataflow
%                                  arrow should use: generic same-role
%                                  passthrough; target's own real output
%                                  leaving the decoder; a real label/mask
%                                  tensor entering fusion or attention
%                                  (context's ground truth, or target's
%                                  own real prior fed back from a
%                                  previous cascade level -- same role
%                                  either way: real mask data, not a
%                                  feature)
%   \archflow{p1}{p2}{role}     -- \archarrow wrapper keyed by role name
%                                  (role in {neutral, target, context})
% ---------------------------------------------------------------
\newcommand{\archarrowset}[1]{%
  \ifthenelse{\equal{#1}{detailed}}{%
    \tikzset{archflowstyle/.style={-{Latex[length=3.6mm,width=3.2mm]},line width=3.2pt}}%
  }{%
    \tikzset{archflowstyle/.style={-{Latex[length=2.4mm,width=2mm]},line width=2.2pt}}%
  }%
  \renewcommand{\archarrow}[3]{\draw[archflowstyle,##3] ##1 -- ##2;}%
}
\newcommand{\archflowNeutral}{black!30}
\newcommand{\archflowTarget}{archTargetAccent}
\newcommand{\archflowContext}{archContextAccent}
\newcommand{\archflow}[3]{%
  \ifthenelse{\equal{#3}{target}}{\archarrow{#1}{#2}{\archflowTarget}}{%
  \ifthenelse{\equal{#3}{context}}{\archarrow{#1}{#2}{\archflowContext}}{%
  \archarrow{#1}{#2}{\archflowNeutral}}}%
}
 after \usepackage{tikz,ifthen} and
% \usetikzlibrary{arrows.meta,calc,positioning}. All macro names are
% prefixed `arch` to avoid clashing with pgf/tikz internals.

% ---- palette (approximating the PowerPoint drafts) ----
\definecolor{archTargetGreen}{HTML}{9BC7AC}
\definecolor{archContextPink}{HTML}{E8B8B0}
\definecolor{archAttnBlue}{HTML}{BBD6EC}
\definecolor{archEncGray}{HTML}{E4E4E4}
\definecolor{archArrowPale}{HTML}{CFE6D6}
\definecolor{archCtBg}{HTML}{2B2B2B}
\definecolor{archCtBlob}{HTML}{8C8C8C}
\definecolor{archCtxYellow}{HTML}{E8D24C}
% darker, more saturated siblings of archTargetGreen/archContextPink --
% pastel fills read poorly as 1pt strokes, and mask/label content needs to
% read as "which volume owns this" at a glance, so masks/labels and any
% line (arrow, brace) touching that volume's row use these instead of the
% pastel role color or the old flat archMaskGold.
\definecolor{archTargetAccent}{HTML}{2F6B4A}  % target: darker green
\definecolor{archContextAccent}{HTML}{A8453E} % context: darker red

% fake-isometric depth skew shared by every cube
\newcommand{\archdx}{0.30}
\newcommand{\archdy}{0.20}

% ---- shared text-role macros ----
% Every arch_*.tex figure is compiled standalone at the SAME base size
% (documentclass[tikz,border=4pt,12pt]{standalone}, matching the thesis
% body's 12pt), so \normalsize etc. mean the same absolute point size in
% every figure. Using these three named roles instead of ad hoc \Large /
% \bfseries\normalsize / \scriptsize per figure keeps "a box title" (etc.)
% the same size everywhere it appears, regardless of which figure it's in
% or how that figure's canvas happens to be scaled by \includegraphics.
%   \archtitle    -- figure/method/panel titles, the "xL" repeat marker
%   \archboxlabel -- sub-block header labels (Fuse, Encoder, attention, ...)
%   \archnote     -- small secondary annotations (placeholder captions)
% Plain, unlabeled node text (cube/token labels, row/lane labels) stays at
% the bare \normalsize default -- no macro needed for that role.
%
% DECLARATIONS, not argument-taking macros -- deliberately, so they can
% precede a raw `\\` inside a tikz node's text (as `{\archboxlabel line
% one\\line two}`) without hiding that `\\` one macro-expansion level deep.
% A `\\` buried inside a wrapped {\archboxlabel{...}} argument breaks
% TikZ's own align=center line-splitting (it scans for `\\` before the
% argument expands) -- use these the same way `\bfseries`/`\Large` are
% already used directly inline elsewhere in these figures.
\newcommand{\archtitle}{\Large\bfseries}
\newcommand{\archboxlabel}{\bfseries\normalsize}
\newcommand{\archnote}{\scriptsize}

% ---------------------------------------------------------------
% \archvolcube{x}{y}{w}{h}{topcolor}{content}{label}
%   content: ct | ctmask | predmask | <flat color name>
% ---------------------------------------------------------------
\newcommand{\archvolcube}[7]{%
  \begin{scope}
    \pgfmathsetmacro{\cx}{#1}\pgfmathsetmacro{\cy}{#2}%
    \pgfmathsetmacro{\w}{#3}\pgfmathsetmacro{\h}{#4}%
    \fill[#5] (\cx,\cy+\h) -- ++(\archdx,\archdy) -- ++(\w,0) -- ++(-\archdx,-\archdy) -- cycle;
    \fill[#5!60!black] (\cx+\w,\cy) -- ++(\archdx,\archdy) -- ++(0,\h) -- ++(-\archdx,-\archdy) -- cycle;
    \begin{scope}
      \clip (\cx,\cy) rectangle ++(\w,\h);
      \ifthenelse{\equal{#6}{ct}}{%
        \fill[archCtBg] (\cx,\cy) rectangle ++(\w,\h);
        \fill[archCtBlob] plot[smooth cycle] coordinates
          {(\cx+.25*\w,\cy+.55*\h) (\cx+.45*\w,\cy+.75*\h) (\cx+.7*\w,\cy+.6*\h)
           (\cx+.6*\w,\cy+.3*\h) (\cx+.35*\w,\cy+.25*\h)};
      }{\ifthenelse{\equal{#6}{ctmask}}{%
        \fill[archCtBg] (\cx,\cy) rectangle ++(\w,\h);
        \fill[archCtBlob] plot[smooth cycle] coordinates
          {(\cx+.25*\w,\cy+.55*\h) (\cx+.45*\w,\cy+.75*\h) (\cx+.7*\w,\cy+.6*\h)
           (\cx+.6*\w,\cy+.3*\h) (\cx+.35*\w,\cy+.25*\h)};
        \fill[archCtxYellow] plot[smooth cycle] coordinates
          {(\cx+.42*\w,\cy+.42*\h) (\cx+.52*\w,\cy+.5*\h) (\cx+.5*\w,\cy+.35*\h) (\cx+.44*\w,\cy+.3*\h)};
      }{\ifthenelse{\equal{#6}{predmask}}{%
        \fill[black] (\cx,\cy) rectangle ++(\w,\h);
        \fill[white] plot[smooth cycle] coordinates
          {(\cx+.4*\w,\cy+.65*\h) (\cx+.6*\w,\cy+.7*\h) (\cx+.62*\w,\cy+.45*\h)
           (\cx+.48*\w,\cy+.3*\h) (\cx+.38*\w,\cy+.4*\h)};
      }{%
        \fill[#6] (\cx,\cy) rectangle ++(\w,\h);
      }}}%
    \end{scope}
    \draw[black,thin] (\cx,\cy) rectangle ++(\w,\h);
    \draw[black,thin] (\cx,\cy+\h) -- ++(\archdx,\archdy) -- ++(\w,0) -- ++(-\archdx,-\archdy) -- cycle;
    \draw[black,thin] (\cx+\w,\cy) -- ++(\archdx,\archdy) -- ++(0,\h) -- ++(-\archdx,-\archdy) -- cycle;
    \node[anchor=north,align=center] at (\cx+\w/2,\cy-0.1) {#7};
  \end{scope}
}

% ---------------------------------------------------------------
% \archvolcubeimg{x}{y}{w}{h}{topcolor}{imgpath}{label} -- same cube as
% \archvolcube, but the front face shows a real raster image (clipped to
% exactly fill the face, aspect ratio not preserved) instead of one of
% \archvolcube's procedural ct/ctmask/predmask blobs. Needs graphicx
% loaded by the including figure. A separate macro rather than a new
% \archvolcube content mode, so existing ct/ctmask/predmask figures are
% untouched.
% ---------------------------------------------------------------
\newcommand{\archvolcubeimg}[7]{%
  \begin{scope}
    \pgfmathsetmacro{\cx}{#1}\pgfmathsetmacro{\cy}{#2}%
    \pgfmathsetmacro{\w}{#3}\pgfmathsetmacro{\h}{#4}%
    \fill[#5] (\cx,\cy+\h) -- ++(\archdx,\archdy) -- ++(\w,0) -- ++(-\archdx,-\archdy) -- cycle;
    \fill[#5!60!black] (\cx+\w,\cy) -- ++(\archdx,\archdy) -- ++(0,\h) -- ++(-\archdx,-\archdy) -- cycle;
    \begin{scope}
      \clip (\cx,\cy) rectangle ++(\w,\h);
      \node[anchor=south west, inner sep=0pt] at (\cx,\cy)
        {\includegraphics[width=\w cm,height=\h cm]{#6}};
    \end{scope}
    \draw[black,thin] (\cx,\cy) rectangle ++(\w,\h);
    \draw[black,thin] (\cx,\cy+\h) -- ++(\archdx,\archdy) -- ++(\w,0) -- ++(-\archdx,-\archdy) -- cycle;
    \draw[black,thin] (\cx+\w,\cy) -- ++(\archdx,\archdy) -- ++(0,\h) -- ++(-\archdx,-\archdy) -- cycle;
    \node[anchor=north,align=center] at (\cx+\w/2,\cy-0.1) {#7};
  \end{scope}
}

% ---------------------------------------------------------------
% \archfeatcube{x}{y}{w}{h}{color}{label} -- flat-colored feature volume
% ---------------------------------------------------------------
\newcommand{\archfeatcube}[6]{%
  \begin{scope}
    \pgfmathsetmacro{\cx}{#1}\pgfmathsetmacro{\cy}{#2}%
    \pgfmathsetmacro{\w}{#3}\pgfmathsetmacro{\h}{#4}%
    \fill[#5] (\cx,\cy) rectangle ++(\w,\h);
    \fill[#5!80!white] (\cx,\cy+\h) -- ++(\archdx,\archdy) -- ++(\w,0) -- ++(-\archdx,-\archdy) -- cycle;
    \fill[#5!60!black] (\cx+\w,\cy) -- ++(\archdx,\archdy) -- ++(0,\h) -- ++(-\archdx,-\archdy) -- cycle;
    \draw[black,thin] (\cx,\cy) rectangle ++(\w,\h);
    \draw[black,thin] (\cx,\cy+\h) -- ++(\archdx,\archdy) -- ++(\w,0) -- ++(-\archdx,-\archdy) -- cycle;
    \draw[black,thin] (\cx+\w,\cy) -- ++(\archdx,\archdy) -- ++(0,\h) -- ++(-\archdx,-\archdy) -- cycle;
    \node[align=center] at (\cx+\w/2,\cy+\h/2) {#6};
  \end{scope}
}

% ---------------------------------------------------------------
% \archfunnel{x}{y}{h}{wwide}{wnarrow}{dir}{color}{label}
%   dir=1  : wide on left, narrow on right  (encoder)
%   dir=-1 : narrow on left, wide on right  (decoder)
% ---------------------------------------------------------------
\newcommand{\archfunnel}[8]{%
  \begin{scope}
    \pgfmathsetmacro{\cx}{#1}\pgfmathsetmacro{\cy}{#2}\pgfmathsetmacro{\h}{#3}%
    \pgfmathsetmacro{\ww}{#4}\pgfmathsetmacro{\wn}{#5}%
    \pgfmathsetmacro{\inset}{(\ww-\wn)/2}%
    \ifnum#6=1
      \fill[#7] (\cx,\cy) -- (\cx+\ww,\cy+\inset) -- (\cx+\ww,\cy+\h-\inset) -- (\cx,\cy+\h) -- cycle;
      \draw[black,thin] (\cx,\cy) -- (\cx+\ww,\cy+\inset) -- (\cx+\ww,\cy+\h-\inset) -- (\cx,\cy+\h) -- cycle;
    \else
      \fill[#7] (\cx,\cy+\inset) -- (\cx+\ww,\cy) -- (\cx+\ww,\cy+\h) -- (\cx,\cy+\h-\inset) -- cycle;
      \draw[black,thin] (\cx,\cy+\inset) -- (\cx+\ww,\cy) -- (\cx+\ww,\cy+\h) -- (\cx,\cy+\h-\inset) -- cycle;
    \fi
    \node[align=center] at (\cx+\ww/2,\cy+\h/2) {#8};
  \end{scope}
}

% ---------------------------------------------------------------
% \archbowtie{x}{y}{w}{h}{pinch}{color}{label} -- Medverse-style
% single-row "dual U-Net" hourglass (funnel-in then funnel-out).
% ---------------------------------------------------------------
\newcommand{\archbowtie}[7]{%
  \begin{scope}
    \pgfmathsetmacro{\cx}{#1}\pgfmathsetmacro{\cy}{#2}%
    \pgfmathsetmacro{\w}{#3}\pgfmathsetmacro{\h}{#4}\pgfmathsetmacro{\p}{#5}%
    \pgfmathsetmacro{\inset}{(\h-\p)/2}%
    \fill[#6] (\cx,\cy) -- (\cx+\w/2,\cy+\inset) -- (\cx+\w,\cy) -- (\cx+\w,\cy+\h)
      -- (\cx+\w/2,\cy+\h-\inset) -- (\cx,\cy+\h) -- cycle;
    \draw[black,thin] (\cx,\cy) -- (\cx+\w/2,\cy+\inset) -- (\cx+\w,\cy) -- (\cx+\w,\cy+\h)
      -- (\cx+\w/2,\cy+\h-\inset) -- (\cx,\cy+\h) -- cycle;
    \node[align=center] at (\cx+\w/2,\cy+\h/2) {#7};
  \end{scope}
}

% ---------------------------------------------------------------
% \archtokengrid{x}{y}{rows}{cols}{cell}{color} -- grid of patch tokens
% ---------------------------------------------------------------
\newcommand{\archtokengrid}[6]{%
  \begin{scope}
    \pgfmathsetmacro{\cx}{#1}\pgfmathsetmacro{\cy}{#2}\pgfmathsetmacro{\cell}{#5}%
    \pgfmathtruncatemacro{\rmax}{#3-1}%
    \pgfmathtruncatemacro{\cmax}{#4-1}%
    \foreach \r in {0,...,\rmax}{%
      \foreach \c in {0,...,\cmax}{%
        \fill[#6] ({\cx+\c*\cell},{\cy+\r*\cell}) rectangle ++(\cell*0.82,\cell*0.82);
        \draw[black,thin] ({\cx+\c*\cell},{\cy+\r*\cell}) rectangle ++(\cell*0.82,\cell*0.82);
      }%
    }%
  \end{scope}
}

% ---------------------------------------------------------------
% \archarrow{start}{end}{color} -- block arrow between stages
% ---------------------------------------------------------------
% #1 and #2 are full coordinates including parentheses, e.g. (1,2)
\newcommand{\archarrow}[3]{%
  \draw[-{Latex[length=2.4mm,width=2mm]},line width=2.2pt,#3] #1 -- #2;
}

% ---------------------------------------------------------------
% Dataflow-arrow vocabulary -- opt-in via \archarrowset{compact|detailed},
% called once near the top of a figure (after % Shared TikZ macros for the architecture figures (Medverse vs. Ours).
% Load with % Shared TikZ macros for the architecture figures (Medverse vs. Ours).
% Load with \input{arch-macros.tex} after \usepackage{tikz,ifthen} and
% \usetikzlibrary{arrows.meta,calc,positioning}. All macro names are
% prefixed `arch` to avoid clashing with pgf/tikz internals.

% ---- palette (approximating the PowerPoint drafts) ----
\definecolor{archTargetGreen}{HTML}{9BC7AC}
\definecolor{archContextPink}{HTML}{E8B8B0}
\definecolor{archAttnBlue}{HTML}{BBD6EC}
\definecolor{archEncGray}{HTML}{E4E4E4}
\definecolor{archArrowPale}{HTML}{CFE6D6}
\definecolor{archCtBg}{HTML}{2B2B2B}
\definecolor{archCtBlob}{HTML}{8C8C8C}
\definecolor{archCtxYellow}{HTML}{E8D24C}
% darker, more saturated siblings of archTargetGreen/archContextPink --
% pastel fills read poorly as 1pt strokes, and mask/label content needs to
% read as "which volume owns this" at a glance, so masks/labels and any
% line (arrow, brace) touching that volume's row use these instead of the
% pastel role color or the old flat archMaskGold.
\definecolor{archTargetAccent}{HTML}{2F6B4A}  % target: darker green
\definecolor{archContextAccent}{HTML}{A8453E} % context: darker red

% fake-isometric depth skew shared by every cube
\newcommand{\archdx}{0.30}
\newcommand{\archdy}{0.20}

% ---- shared text-role macros ----
% Every arch_*.tex figure is compiled standalone at the SAME base size
% (documentclass[tikz,border=4pt,12pt]{standalone}, matching the thesis
% body's 12pt), so \normalsize etc. mean the same absolute point size in
% every figure. Using these three named roles instead of ad hoc \Large /
% \bfseries\normalsize / \scriptsize per figure keeps "a box title" (etc.)
% the same size everywhere it appears, regardless of which figure it's in
% or how that figure's canvas happens to be scaled by \includegraphics.
%   \archtitle    -- figure/method/panel titles, the "xL" repeat marker
%   \archboxlabel -- sub-block header labels (Fuse, Encoder, attention, ...)
%   \archnote     -- small secondary annotations (placeholder captions)
% Plain, unlabeled node text (cube/token labels, row/lane labels) stays at
% the bare \normalsize default -- no macro needed for that role.
%
% DECLARATIONS, not argument-taking macros -- deliberately, so they can
% precede a raw `\\` inside a tikz node's text (as `{\archboxlabel line
% one\\line two}`) without hiding that `\\` one macro-expansion level deep.
% A `\\` buried inside a wrapped {\archboxlabel{...}} argument breaks
% TikZ's own align=center line-splitting (it scans for `\\` before the
% argument expands) -- use these the same way `\bfseries`/`\Large` are
% already used directly inline elsewhere in these figures.
\newcommand{\archtitle}{\Large\bfseries}
\newcommand{\archboxlabel}{\bfseries\normalsize}
\newcommand{\archnote}{\scriptsize}

% ---------------------------------------------------------------
% \archvolcube{x}{y}{w}{h}{topcolor}{content}{label}
%   content: ct | ctmask | predmask | <flat color name>
% ---------------------------------------------------------------
\newcommand{\archvolcube}[7]{%
  \begin{scope}
    \pgfmathsetmacro{\cx}{#1}\pgfmathsetmacro{\cy}{#2}%
    \pgfmathsetmacro{\w}{#3}\pgfmathsetmacro{\h}{#4}%
    \fill[#5] (\cx,\cy+\h) -- ++(\archdx,\archdy) -- ++(\w,0) -- ++(-\archdx,-\archdy) -- cycle;
    \fill[#5!60!black] (\cx+\w,\cy) -- ++(\archdx,\archdy) -- ++(0,\h) -- ++(-\archdx,-\archdy) -- cycle;
    \begin{scope}
      \clip (\cx,\cy) rectangle ++(\w,\h);
      \ifthenelse{\equal{#6}{ct}}{%
        \fill[archCtBg] (\cx,\cy) rectangle ++(\w,\h);
        \fill[archCtBlob] plot[smooth cycle] coordinates
          {(\cx+.25*\w,\cy+.55*\h) (\cx+.45*\w,\cy+.75*\h) (\cx+.7*\w,\cy+.6*\h)
           (\cx+.6*\w,\cy+.3*\h) (\cx+.35*\w,\cy+.25*\h)};
      }{\ifthenelse{\equal{#6}{ctmask}}{%
        \fill[archCtBg] (\cx,\cy) rectangle ++(\w,\h);
        \fill[archCtBlob] plot[smooth cycle] coordinates
          {(\cx+.25*\w,\cy+.55*\h) (\cx+.45*\w,\cy+.75*\h) (\cx+.7*\w,\cy+.6*\h)
           (\cx+.6*\w,\cy+.3*\h) (\cx+.35*\w,\cy+.25*\h)};
        \fill[archCtxYellow] plot[smooth cycle] coordinates
          {(\cx+.42*\w,\cy+.42*\h) (\cx+.52*\w,\cy+.5*\h) (\cx+.5*\w,\cy+.35*\h) (\cx+.44*\w,\cy+.3*\h)};
      }{\ifthenelse{\equal{#6}{predmask}}{%
        \fill[black] (\cx,\cy) rectangle ++(\w,\h);
        \fill[white] plot[smooth cycle] coordinates
          {(\cx+.4*\w,\cy+.65*\h) (\cx+.6*\w,\cy+.7*\h) (\cx+.62*\w,\cy+.45*\h)
           (\cx+.48*\w,\cy+.3*\h) (\cx+.38*\w,\cy+.4*\h)};
      }{%
        \fill[#6] (\cx,\cy) rectangle ++(\w,\h);
      }}}%
    \end{scope}
    \draw[black,thin] (\cx,\cy) rectangle ++(\w,\h);
    \draw[black,thin] (\cx,\cy+\h) -- ++(\archdx,\archdy) -- ++(\w,0) -- ++(-\archdx,-\archdy) -- cycle;
    \draw[black,thin] (\cx+\w,\cy) -- ++(\archdx,\archdy) -- ++(0,\h) -- ++(-\archdx,-\archdy) -- cycle;
    \node[anchor=north,align=center] at (\cx+\w/2,\cy-0.1) {#7};
  \end{scope}
}

% ---------------------------------------------------------------
% \archvolcubeimg{x}{y}{w}{h}{topcolor}{imgpath}{label} -- same cube as
% \archvolcube, but the front face shows a real raster image (clipped to
% exactly fill the face, aspect ratio not preserved) instead of one of
% \archvolcube's procedural ct/ctmask/predmask blobs. Needs graphicx
% loaded by the including figure. A separate macro rather than a new
% \archvolcube content mode, so existing ct/ctmask/predmask figures are
% untouched.
% ---------------------------------------------------------------
\newcommand{\archvolcubeimg}[7]{%
  \begin{scope}
    \pgfmathsetmacro{\cx}{#1}\pgfmathsetmacro{\cy}{#2}%
    \pgfmathsetmacro{\w}{#3}\pgfmathsetmacro{\h}{#4}%
    \fill[#5] (\cx,\cy+\h) -- ++(\archdx,\archdy) -- ++(\w,0) -- ++(-\archdx,-\archdy) -- cycle;
    \fill[#5!60!black] (\cx+\w,\cy) -- ++(\archdx,\archdy) -- ++(0,\h) -- ++(-\archdx,-\archdy) -- cycle;
    \begin{scope}
      \clip (\cx,\cy) rectangle ++(\w,\h);
      \node[anchor=south west, inner sep=0pt] at (\cx,\cy)
        {\includegraphics[width=\w cm,height=\h cm]{#6}};
    \end{scope}
    \draw[black,thin] (\cx,\cy) rectangle ++(\w,\h);
    \draw[black,thin] (\cx,\cy+\h) -- ++(\archdx,\archdy) -- ++(\w,0) -- ++(-\archdx,-\archdy) -- cycle;
    \draw[black,thin] (\cx+\w,\cy) -- ++(\archdx,\archdy) -- ++(0,\h) -- ++(-\archdx,-\archdy) -- cycle;
    \node[anchor=north,align=center] at (\cx+\w/2,\cy-0.1) {#7};
  \end{scope}
}

% ---------------------------------------------------------------
% \archfeatcube{x}{y}{w}{h}{color}{label} -- flat-colored feature volume
% ---------------------------------------------------------------
\newcommand{\archfeatcube}[6]{%
  \begin{scope}
    \pgfmathsetmacro{\cx}{#1}\pgfmathsetmacro{\cy}{#2}%
    \pgfmathsetmacro{\w}{#3}\pgfmathsetmacro{\h}{#4}%
    \fill[#5] (\cx,\cy) rectangle ++(\w,\h);
    \fill[#5!80!white] (\cx,\cy+\h) -- ++(\archdx,\archdy) -- ++(\w,0) -- ++(-\archdx,-\archdy) -- cycle;
    \fill[#5!60!black] (\cx+\w,\cy) -- ++(\archdx,\archdy) -- ++(0,\h) -- ++(-\archdx,-\archdy) -- cycle;
    \draw[black,thin] (\cx,\cy) rectangle ++(\w,\h);
    \draw[black,thin] (\cx,\cy+\h) -- ++(\archdx,\archdy) -- ++(\w,0) -- ++(-\archdx,-\archdy) -- cycle;
    \draw[black,thin] (\cx+\w,\cy) -- ++(\archdx,\archdy) -- ++(0,\h) -- ++(-\archdx,-\archdy) -- cycle;
    \node[align=center] at (\cx+\w/2,\cy+\h/2) {#6};
  \end{scope}
}

% ---------------------------------------------------------------
% \archfunnel{x}{y}{h}{wwide}{wnarrow}{dir}{color}{label}
%   dir=1  : wide on left, narrow on right  (encoder)
%   dir=-1 : narrow on left, wide on right  (decoder)
% ---------------------------------------------------------------
\newcommand{\archfunnel}[8]{%
  \begin{scope}
    \pgfmathsetmacro{\cx}{#1}\pgfmathsetmacro{\cy}{#2}\pgfmathsetmacro{\h}{#3}%
    \pgfmathsetmacro{\ww}{#4}\pgfmathsetmacro{\wn}{#5}%
    \pgfmathsetmacro{\inset}{(\ww-\wn)/2}%
    \ifnum#6=1
      \fill[#7] (\cx,\cy) -- (\cx+\ww,\cy+\inset) -- (\cx+\ww,\cy+\h-\inset) -- (\cx,\cy+\h) -- cycle;
      \draw[black,thin] (\cx,\cy) -- (\cx+\ww,\cy+\inset) -- (\cx+\ww,\cy+\h-\inset) -- (\cx,\cy+\h) -- cycle;
    \else
      \fill[#7] (\cx,\cy+\inset) -- (\cx+\ww,\cy) -- (\cx+\ww,\cy+\h) -- (\cx,\cy+\h-\inset) -- cycle;
      \draw[black,thin] (\cx,\cy+\inset) -- (\cx+\ww,\cy) -- (\cx+\ww,\cy+\h) -- (\cx,\cy+\h-\inset) -- cycle;
    \fi
    \node[align=center] at (\cx+\ww/2,\cy+\h/2) {#8};
  \end{scope}
}

% ---------------------------------------------------------------
% \archbowtie{x}{y}{w}{h}{pinch}{color}{label} -- Medverse-style
% single-row "dual U-Net" hourglass (funnel-in then funnel-out).
% ---------------------------------------------------------------
\newcommand{\archbowtie}[7]{%
  \begin{scope}
    \pgfmathsetmacro{\cx}{#1}\pgfmathsetmacro{\cy}{#2}%
    \pgfmathsetmacro{\w}{#3}\pgfmathsetmacro{\h}{#4}\pgfmathsetmacro{\p}{#5}%
    \pgfmathsetmacro{\inset}{(\h-\p)/2}%
    \fill[#6] (\cx,\cy) -- (\cx+\w/2,\cy+\inset) -- (\cx+\w,\cy) -- (\cx+\w,\cy+\h)
      -- (\cx+\w/2,\cy+\h-\inset) -- (\cx,\cy+\h) -- cycle;
    \draw[black,thin] (\cx,\cy) -- (\cx+\w/2,\cy+\inset) -- (\cx+\w,\cy) -- (\cx+\w,\cy+\h)
      -- (\cx+\w/2,\cy+\h-\inset) -- (\cx,\cy+\h) -- cycle;
    \node[align=center] at (\cx+\w/2,\cy+\h/2) {#7};
  \end{scope}
}

% ---------------------------------------------------------------
% \archtokengrid{x}{y}{rows}{cols}{cell}{color} -- grid of patch tokens
% ---------------------------------------------------------------
\newcommand{\archtokengrid}[6]{%
  \begin{scope}
    \pgfmathsetmacro{\cx}{#1}\pgfmathsetmacro{\cy}{#2}\pgfmathsetmacro{\cell}{#5}%
    \pgfmathtruncatemacro{\rmax}{#3-1}%
    \pgfmathtruncatemacro{\cmax}{#4-1}%
    \foreach \r in {0,...,\rmax}{%
      \foreach \c in {0,...,\cmax}{%
        \fill[#6] ({\cx+\c*\cell},{\cy+\r*\cell}) rectangle ++(\cell*0.82,\cell*0.82);
        \draw[black,thin] ({\cx+\c*\cell},{\cy+\r*\cell}) rectangle ++(\cell*0.82,\cell*0.82);
      }%
    }%
  \end{scope}
}

% ---------------------------------------------------------------
% \archarrow{start}{end}{color} -- block arrow between stages
% ---------------------------------------------------------------
% #1 and #2 are full coordinates including parentheses, e.g. (1,2)
\newcommand{\archarrow}[3]{%
  \draw[-{Latex[length=2.4mm,width=2mm]},line width=2.2pt,#3] #1 -- #2;
}

% ---------------------------------------------------------------
% Dataflow-arrow vocabulary -- opt-in via \archarrowset{compact|detailed},
% called once near the top of a figure (after \input{arch-macros.tex}).
% \archarrow keeps its classic small size by default for any arch_*.tex
% figure that never calls \archarrowset -- this block is purely additive.
%   archflowstyle (tikz key, NOT a macro -- use bare, no backslash) --
%                                  current draw options (head + line
%                                  width), for a hand-rolled elbow \draw
%                                  that \archarrow's straight "--" can't
%                                  express (e.g. "-|" routing). Must be a
%                                  /.style key rather than a \def'd macro:
%                                  pgf's arrow-tip parser doesn't split on
%                                  a comma embedded in a plain macro's
%                                  expansion when used as a \draw option.
%   \archflowNeutral/Target/Context -- the only 3 colors a plain dataflow
%                                  arrow should use: generic same-role
%                                  passthrough; target's own real output
%                                  leaving the decoder; a real label/mask
%                                  tensor entering fusion or attention
%                                  (context's ground truth, or target's
%                                  own real prior fed back from a
%                                  previous cascade level -- same role
%                                  either way: real mask data, not a
%                                  feature)
%   \archflow{p1}{p2}{role}     -- \archarrow wrapper keyed by role name
%                                  (role in {neutral, target, context})
% ---------------------------------------------------------------
\newcommand{\archarrowset}[1]{%
  \ifthenelse{\equal{#1}{detailed}}{%
    \tikzset{archflowstyle/.style={-{Latex[length=3.6mm,width=3.2mm]},line width=3.2pt}}%
  }{%
    \tikzset{archflowstyle/.style={-{Latex[length=2.4mm,width=2mm]},line width=2.2pt}}%
  }%
  \renewcommand{\archarrow}[3]{\draw[archflowstyle,##3] ##1 -- ##2;}%
}
\newcommand{\archflowNeutral}{black!30}
\newcommand{\archflowTarget}{archTargetAccent}
\newcommand{\archflowContext}{archContextAccent}
\newcommand{\archflow}[3]{%
  \ifthenelse{\equal{#3}{target}}{\archarrow{#1}{#2}{\archflowTarget}}{%
  \ifthenelse{\equal{#3}{context}}{\archarrow{#1}{#2}{\archflowContext}}{%
  \archarrow{#1}{#2}{\archflowNeutral}}}%
}
 after \usepackage{tikz,ifthen} and
% \usetikzlibrary{arrows.meta,calc,positioning}. All macro names are
% prefixed `arch` to avoid clashing with pgf/tikz internals.

% ---- palette (approximating the PowerPoint drafts) ----
\definecolor{archTargetGreen}{HTML}{9BC7AC}
\definecolor{archContextPink}{HTML}{E8B8B0}
\definecolor{archAttnBlue}{HTML}{BBD6EC}
\definecolor{archEncGray}{HTML}{E4E4E4}
\definecolor{archArrowPale}{HTML}{CFE6D6}
\definecolor{archCtBg}{HTML}{2B2B2B}
\definecolor{archCtBlob}{HTML}{8C8C8C}
\definecolor{archCtxYellow}{HTML}{E8D24C}
% darker, more saturated siblings of archTargetGreen/archContextPink --
% pastel fills read poorly as 1pt strokes, and mask/label content needs to
% read as "which volume owns this" at a glance, so masks/labels and any
% line (arrow, brace) touching that volume's row use these instead of the
% pastel role color or the old flat archMaskGold.
\definecolor{archTargetAccent}{HTML}{2F6B4A}  % target: darker green
\definecolor{archContextAccent}{HTML}{A8453E} % context: darker red

% fake-isometric depth skew shared by every cube
\newcommand{\archdx}{0.30}
\newcommand{\archdy}{0.20}

% ---- shared text-role macros ----
% Every arch_*.tex figure is compiled standalone at the SAME base size
% (documentclass[tikz,border=4pt,12pt]{standalone}, matching the thesis
% body's 12pt), so \normalsize etc. mean the same absolute point size in
% every figure. Using these three named roles instead of ad hoc \Large /
% \bfseries\normalsize / \scriptsize per figure keeps "a box title" (etc.)
% the same size everywhere it appears, regardless of which figure it's in
% or how that figure's canvas happens to be scaled by \includegraphics.
%   \archtitle    -- figure/method/panel titles, the "xL" repeat marker
%   \archboxlabel -- sub-block header labels (Fuse, Encoder, attention, ...)
%   \archnote     -- small secondary annotations (placeholder captions)
% Plain, unlabeled node text (cube/token labels, row/lane labels) stays at
% the bare \normalsize default -- no macro needed for that role.
%
% DECLARATIONS, not argument-taking macros -- deliberately, so they can
% precede a raw `\\` inside a tikz node's text (as `{\archboxlabel line
% one\\line two}`) without hiding that `\\` one macro-expansion level deep.
% A `\\` buried inside a wrapped {\archboxlabel{...}} argument breaks
% TikZ's own align=center line-splitting (it scans for `\\` before the
% argument expands) -- use these the same way `\bfseries`/`\Large` are
% already used directly inline elsewhere in these figures.
\newcommand{\archtitle}{\Large\bfseries}
\newcommand{\archboxlabel}{\bfseries\normalsize}
\newcommand{\archnote}{\scriptsize}

% ---------------------------------------------------------------
% \archvolcube{x}{y}{w}{h}{topcolor}{content}{label}
%   content: ct | ctmask | predmask | <flat color name>
% ---------------------------------------------------------------
\newcommand{\archvolcube}[7]{%
  \begin{scope}
    \pgfmathsetmacro{\cx}{#1}\pgfmathsetmacro{\cy}{#2}%
    \pgfmathsetmacro{\w}{#3}\pgfmathsetmacro{\h}{#4}%
    \fill[#5] (\cx,\cy+\h) -- ++(\archdx,\archdy) -- ++(\w,0) -- ++(-\archdx,-\archdy) -- cycle;
    \fill[#5!60!black] (\cx+\w,\cy) -- ++(\archdx,\archdy) -- ++(0,\h) -- ++(-\archdx,-\archdy) -- cycle;
    \begin{scope}
      \clip (\cx,\cy) rectangle ++(\w,\h);
      \ifthenelse{\equal{#6}{ct}}{%
        \fill[archCtBg] (\cx,\cy) rectangle ++(\w,\h);
        \fill[archCtBlob] plot[smooth cycle] coordinates
          {(\cx+.25*\w,\cy+.55*\h) (\cx+.45*\w,\cy+.75*\h) (\cx+.7*\w,\cy+.6*\h)
           (\cx+.6*\w,\cy+.3*\h) (\cx+.35*\w,\cy+.25*\h)};
      }{\ifthenelse{\equal{#6}{ctmask}}{%
        \fill[archCtBg] (\cx,\cy) rectangle ++(\w,\h);
        \fill[archCtBlob] plot[smooth cycle] coordinates
          {(\cx+.25*\w,\cy+.55*\h) (\cx+.45*\w,\cy+.75*\h) (\cx+.7*\w,\cy+.6*\h)
           (\cx+.6*\w,\cy+.3*\h) (\cx+.35*\w,\cy+.25*\h)};
        \fill[archCtxYellow] plot[smooth cycle] coordinates
          {(\cx+.42*\w,\cy+.42*\h) (\cx+.52*\w,\cy+.5*\h) (\cx+.5*\w,\cy+.35*\h) (\cx+.44*\w,\cy+.3*\h)};
      }{\ifthenelse{\equal{#6}{predmask}}{%
        \fill[black] (\cx,\cy) rectangle ++(\w,\h);
        \fill[white] plot[smooth cycle] coordinates
          {(\cx+.4*\w,\cy+.65*\h) (\cx+.6*\w,\cy+.7*\h) (\cx+.62*\w,\cy+.45*\h)
           (\cx+.48*\w,\cy+.3*\h) (\cx+.38*\w,\cy+.4*\h)};
      }{%
        \fill[#6] (\cx,\cy) rectangle ++(\w,\h);
      }}}%
    \end{scope}
    \draw[black,thin] (\cx,\cy) rectangle ++(\w,\h);
    \draw[black,thin] (\cx,\cy+\h) -- ++(\archdx,\archdy) -- ++(\w,0) -- ++(-\archdx,-\archdy) -- cycle;
    \draw[black,thin] (\cx+\w,\cy) -- ++(\archdx,\archdy) -- ++(0,\h) -- ++(-\archdx,-\archdy) -- cycle;
    \node[anchor=north,align=center] at (\cx+\w/2,\cy-0.1) {#7};
  \end{scope}
}

% ---------------------------------------------------------------
% \archvolcubeimg{x}{y}{w}{h}{topcolor}{imgpath}{label} -- same cube as
% \archvolcube, but the front face shows a real raster image (clipped to
% exactly fill the face, aspect ratio not preserved) instead of one of
% \archvolcube's procedural ct/ctmask/predmask blobs. Needs graphicx
% loaded by the including figure. A separate macro rather than a new
% \archvolcube content mode, so existing ct/ctmask/predmask figures are
% untouched.
% ---------------------------------------------------------------
\newcommand{\archvolcubeimg}[7]{%
  \begin{scope}
    \pgfmathsetmacro{\cx}{#1}\pgfmathsetmacro{\cy}{#2}%
    \pgfmathsetmacro{\w}{#3}\pgfmathsetmacro{\h}{#4}%
    \fill[#5] (\cx,\cy+\h) -- ++(\archdx,\archdy) -- ++(\w,0) -- ++(-\archdx,-\archdy) -- cycle;
    \fill[#5!60!black] (\cx+\w,\cy) -- ++(\archdx,\archdy) -- ++(0,\h) -- ++(-\archdx,-\archdy) -- cycle;
    \begin{scope}
      \clip (\cx,\cy) rectangle ++(\w,\h);
      \node[anchor=south west, inner sep=0pt] at (\cx,\cy)
        {\includegraphics[width=\w cm,height=\h cm]{#6}};
    \end{scope}
    \draw[black,thin] (\cx,\cy) rectangle ++(\w,\h);
    \draw[black,thin] (\cx,\cy+\h) -- ++(\archdx,\archdy) -- ++(\w,0) -- ++(-\archdx,-\archdy) -- cycle;
    \draw[black,thin] (\cx+\w,\cy) -- ++(\archdx,\archdy) -- ++(0,\h) -- ++(-\archdx,-\archdy) -- cycle;
    \node[anchor=north,align=center] at (\cx+\w/2,\cy-0.1) {#7};
  \end{scope}
}

% ---------------------------------------------------------------
% \archfeatcube{x}{y}{w}{h}{color}{label} -- flat-colored feature volume
% ---------------------------------------------------------------
\newcommand{\archfeatcube}[6]{%
  \begin{scope}
    \pgfmathsetmacro{\cx}{#1}\pgfmathsetmacro{\cy}{#2}%
    \pgfmathsetmacro{\w}{#3}\pgfmathsetmacro{\h}{#4}%
    \fill[#5] (\cx,\cy) rectangle ++(\w,\h);
    \fill[#5!80!white] (\cx,\cy+\h) -- ++(\archdx,\archdy) -- ++(\w,0) -- ++(-\archdx,-\archdy) -- cycle;
    \fill[#5!60!black] (\cx+\w,\cy) -- ++(\archdx,\archdy) -- ++(0,\h) -- ++(-\archdx,-\archdy) -- cycle;
    \draw[black,thin] (\cx,\cy) rectangle ++(\w,\h);
    \draw[black,thin] (\cx,\cy+\h) -- ++(\archdx,\archdy) -- ++(\w,0) -- ++(-\archdx,-\archdy) -- cycle;
    \draw[black,thin] (\cx+\w,\cy) -- ++(\archdx,\archdy) -- ++(0,\h) -- ++(-\archdx,-\archdy) -- cycle;
    \node[align=center] at (\cx+\w/2,\cy+\h/2) {#6};
  \end{scope}
}

% ---------------------------------------------------------------
% \archfunnel{x}{y}{h}{wwide}{wnarrow}{dir}{color}{label}
%   dir=1  : wide on left, narrow on right  (encoder)
%   dir=-1 : narrow on left, wide on right  (decoder)
% ---------------------------------------------------------------
\newcommand{\archfunnel}[8]{%
  \begin{scope}
    \pgfmathsetmacro{\cx}{#1}\pgfmathsetmacro{\cy}{#2}\pgfmathsetmacro{\h}{#3}%
    \pgfmathsetmacro{\ww}{#4}\pgfmathsetmacro{\wn}{#5}%
    \pgfmathsetmacro{\inset}{(\ww-\wn)/2}%
    \ifnum#6=1
      \fill[#7] (\cx,\cy) -- (\cx+\ww,\cy+\inset) -- (\cx+\ww,\cy+\h-\inset) -- (\cx,\cy+\h) -- cycle;
      \draw[black,thin] (\cx,\cy) -- (\cx+\ww,\cy+\inset) -- (\cx+\ww,\cy+\h-\inset) -- (\cx,\cy+\h) -- cycle;
    \else
      \fill[#7] (\cx,\cy+\inset) -- (\cx+\ww,\cy) -- (\cx+\ww,\cy+\h) -- (\cx,\cy+\h-\inset) -- cycle;
      \draw[black,thin] (\cx,\cy+\inset) -- (\cx+\ww,\cy) -- (\cx+\ww,\cy+\h) -- (\cx,\cy+\h-\inset) -- cycle;
    \fi
    \node[align=center] at (\cx+\ww/2,\cy+\h/2) {#8};
  \end{scope}
}

% ---------------------------------------------------------------
% \archbowtie{x}{y}{w}{h}{pinch}{color}{label} -- Medverse-style
% single-row "dual U-Net" hourglass (funnel-in then funnel-out).
% ---------------------------------------------------------------
\newcommand{\archbowtie}[7]{%
  \begin{scope}
    \pgfmathsetmacro{\cx}{#1}\pgfmathsetmacro{\cy}{#2}%
    \pgfmathsetmacro{\w}{#3}\pgfmathsetmacro{\h}{#4}\pgfmathsetmacro{\p}{#5}%
    \pgfmathsetmacro{\inset}{(\h-\p)/2}%
    \fill[#6] (\cx,\cy) -- (\cx+\w/2,\cy+\inset) -- (\cx+\w,\cy) -- (\cx+\w,\cy+\h)
      -- (\cx+\w/2,\cy+\h-\inset) -- (\cx,\cy+\h) -- cycle;
    \draw[black,thin] (\cx,\cy) -- (\cx+\w/2,\cy+\inset) -- (\cx+\w,\cy) -- (\cx+\w,\cy+\h)
      -- (\cx+\w/2,\cy+\h-\inset) -- (\cx,\cy+\h) -- cycle;
    \node[align=center] at (\cx+\w/2,\cy+\h/2) {#7};
  \end{scope}
}

% ---------------------------------------------------------------
% \archtokengrid{x}{y}{rows}{cols}{cell}{color} -- grid of patch tokens
% ---------------------------------------------------------------
\newcommand{\archtokengrid}[6]{%
  \begin{scope}
    \pgfmathsetmacro{\cx}{#1}\pgfmathsetmacro{\cy}{#2}\pgfmathsetmacro{\cell}{#5}%
    \pgfmathtruncatemacro{\rmax}{#3-1}%
    \pgfmathtruncatemacro{\cmax}{#4-1}%
    \foreach \r in {0,...,\rmax}{%
      \foreach \c in {0,...,\cmax}{%
        \fill[#6] ({\cx+\c*\cell},{\cy+\r*\cell}) rectangle ++(\cell*0.82,\cell*0.82);
        \draw[black,thin] ({\cx+\c*\cell},{\cy+\r*\cell}) rectangle ++(\cell*0.82,\cell*0.82);
      }%
    }%
  \end{scope}
}

% ---------------------------------------------------------------
% \archarrow{start}{end}{color} -- block arrow between stages
% ---------------------------------------------------------------
% #1 and #2 are full coordinates including parentheses, e.g. (1,2)
\newcommand{\archarrow}[3]{%
  \draw[-{Latex[length=2.4mm,width=2mm]},line width=2.2pt,#3] #1 -- #2;
}

% ---------------------------------------------------------------
% Dataflow-arrow vocabulary -- opt-in via \archarrowset{compact|detailed},
% called once near the top of a figure (after % Shared TikZ macros for the architecture figures (Medverse vs. Ours).
% Load with \input{arch-macros.tex} after \usepackage{tikz,ifthen} and
% \usetikzlibrary{arrows.meta,calc,positioning}. All macro names are
% prefixed `arch` to avoid clashing with pgf/tikz internals.

% ---- palette (approximating the PowerPoint drafts) ----
\definecolor{archTargetGreen}{HTML}{9BC7AC}
\definecolor{archContextPink}{HTML}{E8B8B0}
\definecolor{archAttnBlue}{HTML}{BBD6EC}
\definecolor{archEncGray}{HTML}{E4E4E4}
\definecolor{archArrowPale}{HTML}{CFE6D6}
\definecolor{archCtBg}{HTML}{2B2B2B}
\definecolor{archCtBlob}{HTML}{8C8C8C}
\definecolor{archCtxYellow}{HTML}{E8D24C}
% darker, more saturated siblings of archTargetGreen/archContextPink --
% pastel fills read poorly as 1pt strokes, and mask/label content needs to
% read as "which volume owns this" at a glance, so masks/labels and any
% line (arrow, brace) touching that volume's row use these instead of the
% pastel role color or the old flat archMaskGold.
\definecolor{archTargetAccent}{HTML}{2F6B4A}  % target: darker green
\definecolor{archContextAccent}{HTML}{A8453E} % context: darker red

% fake-isometric depth skew shared by every cube
\newcommand{\archdx}{0.30}
\newcommand{\archdy}{0.20}

% ---- shared text-role macros ----
% Every arch_*.tex figure is compiled standalone at the SAME base size
% (documentclass[tikz,border=4pt,12pt]{standalone}, matching the thesis
% body's 12pt), so \normalsize etc. mean the same absolute point size in
% every figure. Using these three named roles instead of ad hoc \Large /
% \bfseries\normalsize / \scriptsize per figure keeps "a box title" (etc.)
% the same size everywhere it appears, regardless of which figure it's in
% or how that figure's canvas happens to be scaled by \includegraphics.
%   \archtitle    -- figure/method/panel titles, the "xL" repeat marker
%   \archboxlabel -- sub-block header labels (Fuse, Encoder, attention, ...)
%   \archnote     -- small secondary annotations (placeholder captions)
% Plain, unlabeled node text (cube/token labels, row/lane labels) stays at
% the bare \normalsize default -- no macro needed for that role.
%
% DECLARATIONS, not argument-taking macros -- deliberately, so they can
% precede a raw `\\` inside a tikz node's text (as `{\archboxlabel line
% one\\line two}`) without hiding that `\\` one macro-expansion level deep.
% A `\\` buried inside a wrapped {\archboxlabel{...}} argument breaks
% TikZ's own align=center line-splitting (it scans for `\\` before the
% argument expands) -- use these the same way `\bfseries`/`\Large` are
% already used directly inline elsewhere in these figures.
\newcommand{\archtitle}{\Large\bfseries}
\newcommand{\archboxlabel}{\bfseries\normalsize}
\newcommand{\archnote}{\scriptsize}

% ---------------------------------------------------------------
% \archvolcube{x}{y}{w}{h}{topcolor}{content}{label}
%   content: ct | ctmask | predmask | <flat color name>
% ---------------------------------------------------------------
\newcommand{\archvolcube}[7]{%
  \begin{scope}
    \pgfmathsetmacro{\cx}{#1}\pgfmathsetmacro{\cy}{#2}%
    \pgfmathsetmacro{\w}{#3}\pgfmathsetmacro{\h}{#4}%
    \fill[#5] (\cx,\cy+\h) -- ++(\archdx,\archdy) -- ++(\w,0) -- ++(-\archdx,-\archdy) -- cycle;
    \fill[#5!60!black] (\cx+\w,\cy) -- ++(\archdx,\archdy) -- ++(0,\h) -- ++(-\archdx,-\archdy) -- cycle;
    \begin{scope}
      \clip (\cx,\cy) rectangle ++(\w,\h);
      \ifthenelse{\equal{#6}{ct}}{%
        \fill[archCtBg] (\cx,\cy) rectangle ++(\w,\h);
        \fill[archCtBlob] plot[smooth cycle] coordinates
          {(\cx+.25*\w,\cy+.55*\h) (\cx+.45*\w,\cy+.75*\h) (\cx+.7*\w,\cy+.6*\h)
           (\cx+.6*\w,\cy+.3*\h) (\cx+.35*\w,\cy+.25*\h)};
      }{\ifthenelse{\equal{#6}{ctmask}}{%
        \fill[archCtBg] (\cx,\cy) rectangle ++(\w,\h);
        \fill[archCtBlob] plot[smooth cycle] coordinates
          {(\cx+.25*\w,\cy+.55*\h) (\cx+.45*\w,\cy+.75*\h) (\cx+.7*\w,\cy+.6*\h)
           (\cx+.6*\w,\cy+.3*\h) (\cx+.35*\w,\cy+.25*\h)};
        \fill[archCtxYellow] plot[smooth cycle] coordinates
          {(\cx+.42*\w,\cy+.42*\h) (\cx+.52*\w,\cy+.5*\h) (\cx+.5*\w,\cy+.35*\h) (\cx+.44*\w,\cy+.3*\h)};
      }{\ifthenelse{\equal{#6}{predmask}}{%
        \fill[black] (\cx,\cy) rectangle ++(\w,\h);
        \fill[white] plot[smooth cycle] coordinates
          {(\cx+.4*\w,\cy+.65*\h) (\cx+.6*\w,\cy+.7*\h) (\cx+.62*\w,\cy+.45*\h)
           (\cx+.48*\w,\cy+.3*\h) (\cx+.38*\w,\cy+.4*\h)};
      }{%
        \fill[#6] (\cx,\cy) rectangle ++(\w,\h);
      }}}%
    \end{scope}
    \draw[black,thin] (\cx,\cy) rectangle ++(\w,\h);
    \draw[black,thin] (\cx,\cy+\h) -- ++(\archdx,\archdy) -- ++(\w,0) -- ++(-\archdx,-\archdy) -- cycle;
    \draw[black,thin] (\cx+\w,\cy) -- ++(\archdx,\archdy) -- ++(0,\h) -- ++(-\archdx,-\archdy) -- cycle;
    \node[anchor=north,align=center] at (\cx+\w/2,\cy-0.1) {#7};
  \end{scope}
}

% ---------------------------------------------------------------
% \archvolcubeimg{x}{y}{w}{h}{topcolor}{imgpath}{label} -- same cube as
% \archvolcube, but the front face shows a real raster image (clipped to
% exactly fill the face, aspect ratio not preserved) instead of one of
% \archvolcube's procedural ct/ctmask/predmask blobs. Needs graphicx
% loaded by the including figure. A separate macro rather than a new
% \archvolcube content mode, so existing ct/ctmask/predmask figures are
% untouched.
% ---------------------------------------------------------------
\newcommand{\archvolcubeimg}[7]{%
  \begin{scope}
    \pgfmathsetmacro{\cx}{#1}\pgfmathsetmacro{\cy}{#2}%
    \pgfmathsetmacro{\w}{#3}\pgfmathsetmacro{\h}{#4}%
    \fill[#5] (\cx,\cy+\h) -- ++(\archdx,\archdy) -- ++(\w,0) -- ++(-\archdx,-\archdy) -- cycle;
    \fill[#5!60!black] (\cx+\w,\cy) -- ++(\archdx,\archdy) -- ++(0,\h) -- ++(-\archdx,-\archdy) -- cycle;
    \begin{scope}
      \clip (\cx,\cy) rectangle ++(\w,\h);
      \node[anchor=south west, inner sep=0pt] at (\cx,\cy)
        {\includegraphics[width=\w cm,height=\h cm]{#6}};
    \end{scope}
    \draw[black,thin] (\cx,\cy) rectangle ++(\w,\h);
    \draw[black,thin] (\cx,\cy+\h) -- ++(\archdx,\archdy) -- ++(\w,0) -- ++(-\archdx,-\archdy) -- cycle;
    \draw[black,thin] (\cx+\w,\cy) -- ++(\archdx,\archdy) -- ++(0,\h) -- ++(-\archdx,-\archdy) -- cycle;
    \node[anchor=north,align=center] at (\cx+\w/2,\cy-0.1) {#7};
  \end{scope}
}

% ---------------------------------------------------------------
% \archfeatcube{x}{y}{w}{h}{color}{label} -- flat-colored feature volume
% ---------------------------------------------------------------
\newcommand{\archfeatcube}[6]{%
  \begin{scope}
    \pgfmathsetmacro{\cx}{#1}\pgfmathsetmacro{\cy}{#2}%
    \pgfmathsetmacro{\w}{#3}\pgfmathsetmacro{\h}{#4}%
    \fill[#5] (\cx,\cy) rectangle ++(\w,\h);
    \fill[#5!80!white] (\cx,\cy+\h) -- ++(\archdx,\archdy) -- ++(\w,0) -- ++(-\archdx,-\archdy) -- cycle;
    \fill[#5!60!black] (\cx+\w,\cy) -- ++(\archdx,\archdy) -- ++(0,\h) -- ++(-\archdx,-\archdy) -- cycle;
    \draw[black,thin] (\cx,\cy) rectangle ++(\w,\h);
    \draw[black,thin] (\cx,\cy+\h) -- ++(\archdx,\archdy) -- ++(\w,0) -- ++(-\archdx,-\archdy) -- cycle;
    \draw[black,thin] (\cx+\w,\cy) -- ++(\archdx,\archdy) -- ++(0,\h) -- ++(-\archdx,-\archdy) -- cycle;
    \node[align=center] at (\cx+\w/2,\cy+\h/2) {#6};
  \end{scope}
}

% ---------------------------------------------------------------
% \archfunnel{x}{y}{h}{wwide}{wnarrow}{dir}{color}{label}
%   dir=1  : wide on left, narrow on right  (encoder)
%   dir=-1 : narrow on left, wide on right  (decoder)
% ---------------------------------------------------------------
\newcommand{\archfunnel}[8]{%
  \begin{scope}
    \pgfmathsetmacro{\cx}{#1}\pgfmathsetmacro{\cy}{#2}\pgfmathsetmacro{\h}{#3}%
    \pgfmathsetmacro{\ww}{#4}\pgfmathsetmacro{\wn}{#5}%
    \pgfmathsetmacro{\inset}{(\ww-\wn)/2}%
    \ifnum#6=1
      \fill[#7] (\cx,\cy) -- (\cx+\ww,\cy+\inset) -- (\cx+\ww,\cy+\h-\inset) -- (\cx,\cy+\h) -- cycle;
      \draw[black,thin] (\cx,\cy) -- (\cx+\ww,\cy+\inset) -- (\cx+\ww,\cy+\h-\inset) -- (\cx,\cy+\h) -- cycle;
    \else
      \fill[#7] (\cx,\cy+\inset) -- (\cx+\ww,\cy) -- (\cx+\ww,\cy+\h) -- (\cx,\cy+\h-\inset) -- cycle;
      \draw[black,thin] (\cx,\cy+\inset) -- (\cx+\ww,\cy) -- (\cx+\ww,\cy+\h) -- (\cx,\cy+\h-\inset) -- cycle;
    \fi
    \node[align=center] at (\cx+\ww/2,\cy+\h/2) {#8};
  \end{scope}
}

% ---------------------------------------------------------------
% \archbowtie{x}{y}{w}{h}{pinch}{color}{label} -- Medverse-style
% single-row "dual U-Net" hourglass (funnel-in then funnel-out).
% ---------------------------------------------------------------
\newcommand{\archbowtie}[7]{%
  \begin{scope}
    \pgfmathsetmacro{\cx}{#1}\pgfmathsetmacro{\cy}{#2}%
    \pgfmathsetmacro{\w}{#3}\pgfmathsetmacro{\h}{#4}\pgfmathsetmacro{\p}{#5}%
    \pgfmathsetmacro{\inset}{(\h-\p)/2}%
    \fill[#6] (\cx,\cy) -- (\cx+\w/2,\cy+\inset) -- (\cx+\w,\cy) -- (\cx+\w,\cy+\h)
      -- (\cx+\w/2,\cy+\h-\inset) -- (\cx,\cy+\h) -- cycle;
    \draw[black,thin] (\cx,\cy) -- (\cx+\w/2,\cy+\inset) -- (\cx+\w,\cy) -- (\cx+\w,\cy+\h)
      -- (\cx+\w/2,\cy+\h-\inset) -- (\cx,\cy+\h) -- cycle;
    \node[align=center] at (\cx+\w/2,\cy+\h/2) {#7};
  \end{scope}
}

% ---------------------------------------------------------------
% \archtokengrid{x}{y}{rows}{cols}{cell}{color} -- grid of patch tokens
% ---------------------------------------------------------------
\newcommand{\archtokengrid}[6]{%
  \begin{scope}
    \pgfmathsetmacro{\cx}{#1}\pgfmathsetmacro{\cy}{#2}\pgfmathsetmacro{\cell}{#5}%
    \pgfmathtruncatemacro{\rmax}{#3-1}%
    \pgfmathtruncatemacro{\cmax}{#4-1}%
    \foreach \r in {0,...,\rmax}{%
      \foreach \c in {0,...,\cmax}{%
        \fill[#6] ({\cx+\c*\cell},{\cy+\r*\cell}) rectangle ++(\cell*0.82,\cell*0.82);
        \draw[black,thin] ({\cx+\c*\cell},{\cy+\r*\cell}) rectangle ++(\cell*0.82,\cell*0.82);
      }%
    }%
  \end{scope}
}

% ---------------------------------------------------------------
% \archarrow{start}{end}{color} -- block arrow between stages
% ---------------------------------------------------------------
% #1 and #2 are full coordinates including parentheses, e.g. (1,2)
\newcommand{\archarrow}[3]{%
  \draw[-{Latex[length=2.4mm,width=2mm]},line width=2.2pt,#3] #1 -- #2;
}

% ---------------------------------------------------------------
% Dataflow-arrow vocabulary -- opt-in via \archarrowset{compact|detailed},
% called once near the top of a figure (after \input{arch-macros.tex}).
% \archarrow keeps its classic small size by default for any arch_*.tex
% figure that never calls \archarrowset -- this block is purely additive.
%   archflowstyle (tikz key, NOT a macro -- use bare, no backslash) --
%                                  current draw options (head + line
%                                  width), for a hand-rolled elbow \draw
%                                  that \archarrow's straight "--" can't
%                                  express (e.g. "-|" routing). Must be a
%                                  /.style key rather than a \def'd macro:
%                                  pgf's arrow-tip parser doesn't split on
%                                  a comma embedded in a plain macro's
%                                  expansion when used as a \draw option.
%   \archflowNeutral/Target/Context -- the only 3 colors a plain dataflow
%                                  arrow should use: generic same-role
%                                  passthrough; target's own real output
%                                  leaving the decoder; a real label/mask
%                                  tensor entering fusion or attention
%                                  (context's ground truth, or target's
%                                  own real prior fed back from a
%                                  previous cascade level -- same role
%                                  either way: real mask data, not a
%                                  feature)
%   \archflow{p1}{p2}{role}     -- \archarrow wrapper keyed by role name
%                                  (role in {neutral, target, context})
% ---------------------------------------------------------------
\newcommand{\archarrowset}[1]{%
  \ifthenelse{\equal{#1}{detailed}}{%
    \tikzset{archflowstyle/.style={-{Latex[length=3.6mm,width=3.2mm]},line width=3.2pt}}%
  }{%
    \tikzset{archflowstyle/.style={-{Latex[length=2.4mm,width=2mm]},line width=2.2pt}}%
  }%
  \renewcommand{\archarrow}[3]{\draw[archflowstyle,##3] ##1 -- ##2;}%
}
\newcommand{\archflowNeutral}{black!30}
\newcommand{\archflowTarget}{archTargetAccent}
\newcommand{\archflowContext}{archContextAccent}
\newcommand{\archflow}[3]{%
  \ifthenelse{\equal{#3}{target}}{\archarrow{#1}{#2}{\archflowTarget}}{%
  \ifthenelse{\equal{#3}{context}}{\archarrow{#1}{#2}{\archflowContext}}{%
  \archarrow{#1}{#2}{\archflowNeutral}}}%
}
).
% \archarrow keeps its classic small size by default for any arch_*.tex
% figure that never calls \archarrowset -- this block is purely additive.
%   archflowstyle (tikz key, NOT a macro -- use bare, no backslash) --
%                                  current draw options (head + line
%                                  width), for a hand-rolled elbow \draw
%                                  that \archarrow's straight "--" can't
%                                  express (e.g. "-|" routing). Must be a
%                                  /.style key rather than a \def'd macro:
%                                  pgf's arrow-tip parser doesn't split on
%                                  a comma embedded in a plain macro's
%                                  expansion when used as a \draw option.
%   \archflowNeutral/Target/Context -- the only 3 colors a plain dataflow
%                                  arrow should use: generic same-role
%                                  passthrough; target's own real output
%                                  leaving the decoder; a real label/mask
%                                  tensor entering fusion or attention
%                                  (context's ground truth, or target's
%                                  own real prior fed back from a
%                                  previous cascade level -- same role
%                                  either way: real mask data, not a
%                                  feature)
%   \archflow{p1}{p2}{role}     -- \archarrow wrapper keyed by role name
%                                  (role in {neutral, target, context})
% ---------------------------------------------------------------
\newcommand{\archarrowset}[1]{%
  \ifthenelse{\equal{#1}{detailed}}{%
    \tikzset{archflowstyle/.style={-{Latex[length=3.6mm,width=3.2mm]},line width=3.2pt}}%
  }{%
    \tikzset{archflowstyle/.style={-{Latex[length=2.4mm,width=2mm]},line width=2.2pt}}%
  }%
  \renewcommand{\archarrow}[3]{\draw[archflowstyle,##3] ##1 -- ##2;}%
}
\newcommand{\archflowNeutral}{black!30}
\newcommand{\archflowTarget}{archTargetAccent}
\newcommand{\archflowContext}{archContextAccent}
\newcommand{\archflow}[3]{%
  \ifthenelse{\equal{#3}{target}}{\archarrow{#1}{#2}{\archflowTarget}}{%
  \ifthenelse{\equal{#3}{context}}{\archarrow{#1}{#2}{\archflowContext}}{%
  \archarrow{#1}{#2}{\archflowNeutral}}}%
}
).
% \archarrow keeps its classic small size by default for any arch_*.tex
% figure that never calls \archarrowset -- this block is purely additive.
%   archflowstyle (tikz key, NOT a macro -- use bare, no backslash) --
%                                  current draw options (head + line
%                                  width), for a hand-rolled elbow \draw
%                                  that \archarrow's straight "--" can't
%                                  express (e.g. "-|" routing). Must be a
%                                  /.style key rather than a \def'd macro:
%                                  pgf's arrow-tip parser doesn't split on
%                                  a comma embedded in a plain macro's
%                                  expansion when used as a \draw option.
%   \archflowNeutral/Target/Context -- the only 3 colors a plain dataflow
%                                  arrow should use: generic same-role
%                                  passthrough; target's own real output
%                                  leaving the decoder; a real label/mask
%                                  tensor entering fusion or attention
%                                  (context's ground truth, or target's
%                                  own real prior fed back from a
%                                  previous cascade level -- same role
%                                  either way: real mask data, not a
%                                  feature)
%   \archflow{p1}{p2}{role}     -- \archarrow wrapper keyed by role name
%                                  (role in {neutral, target, context})
% ---------------------------------------------------------------
\newcommand{\archarrowset}[1]{%
  \ifthenelse{\equal{#1}{detailed}}{%
    \tikzset{archflowstyle/.style={-{Latex[length=3.6mm,width=3.2mm]},line width=3.2pt}}%
  }{%
    \tikzset{archflowstyle/.style={-{Latex[length=2.4mm,width=2mm]},line width=2.2pt}}%
  }%
  \renewcommand{\archarrow}[3]{\draw[archflowstyle,##3] ##1 -- ##2;}%
}
\newcommand{\archflowNeutral}{black!30}
\newcommand{\archflowTarget}{archTargetAccent}
\newcommand{\archflowContext}{archContextAccent}
\newcommand{\archflow}[3]{%
  \ifthenelse{\equal{#3}{target}}{\archarrow{#1}{#2}{\archflowTarget}}{%
  \ifthenelse{\equal{#3}{context}}{\archarrow{#1}{#2}{\archflowContext}}{%
  \archarrow{#1}{#2}{\archflowNeutral}}}%
}
).
% \archarrow keeps its classic small size by default for any arch_*.tex
% figure that never calls \archarrowset -- this block is purely additive.
%   archflowstyle (tikz key, NOT a macro -- use bare, no backslash) --
%                                  current draw options (head + line
%                                  width), for a hand-rolled elbow \draw
%                                  that \archarrow's straight "--" can't
%                                  express (e.g. "-|" routing). Must be a
%                                  /.style key rather than a \def'd macro:
%                                  pgf's arrow-tip parser doesn't split on
%                                  a comma embedded in a plain macro's
%                                  expansion when used as a \draw option.
%   \archflowNeutral/Target/Context -- the only 3 colors a plain dataflow
%                                  arrow should use: generic same-role
%                                  passthrough; target's own real output
%                                  leaving the decoder; a real label/mask
%                                  tensor entering fusion or attention
%                                  (context's ground truth, or target's
%                                  own real prior fed back from a
%                                  previous cascade level -- same role
%                                  either way: real mask data, not a
%                                  feature)
%   \archflow{p1}{p2}{role}     -- \archarrow wrapper keyed by role name
%                                  (role in {neutral, target, context})
% ---------------------------------------------------------------
\newcommand{\archarrowset}[1]{%
  \ifthenelse{\equal{#1}{detailed}}{%
    \tikzset{archflowstyle/.style={-{Latex[length=3.6mm,width=3.2mm]},line width=3.2pt}}%
  }{%
    \tikzset{archflowstyle/.style={-{Latex[length=2.4mm,width=2mm]},line width=2.2pt}}%
  }%
  \renewcommand{\archarrow}[3]{\draw[archflowstyle,##3] ##1 -- ##2;}%
}
\newcommand{\archflowNeutral}{black!30}
\newcommand{\archflowTarget}{archTargetAccent}
\newcommand{\archflowContext}{archContextAccent}
\newcommand{\archflow}[3]{%
  \ifthenelse{\equal{#3}{target}}{\archarrow{#1}{#2}{\archflowTarget}}{%
  \ifthenelse{\equal{#3}{context}}{\archarrow{#1}{#2}{\archflowContext}}{%
  \archarrow{#1}{#2}{\archflowNeutral}}}%
}
